\documentclass[11pt]{article}
\usepackage[T1]{fontenc}
\usepackage{lmodern}
\usepackage[margin=1in]{geometry}
\usepackage{amsmath,amssymb,amsthm,mathtools,mathrsfs,booktabs,longtable,array,microtype,xcolor}
\usepackage[colorlinks=true,linkcolor=blue!45!black,citecolor=blue!45!black,urlcolor=blue!45!black]{hyperref}
\hypersetup{pdftitle={Euler--Gompertz and Bessel Values: Joint Independence, Integral Intersections, and Critical Limits},pdfauthor={Keio}}
\newtheorem{theorem}{Theorem}[section]
\newtheorem{proposition}[theorem]{Proposition}
\newtheorem{corollary}[theorem]{Corollary}
\newtheorem{lemma}[theorem]{Lemma}
\theoremstyle{definition}
\newtheorem{sourceinput}[theorem]{Input}
\newtheorem{definition}[theorem]{Definition}
\theoremstyle{remark}
\newtheorem{remark}[theorem]{Remark}
\newcommand{\Qa}{\overline{\mathbb Q}}
\newcommand{\Qbar}{\overline{\mathbb Q}}
\newcommand{\E}{\mathbf E}
\newcommand{\Ebar}{\overline{\mathbf E}}
\newcommand{\C}{\mathbb C}
\newcommand{\trdeg}{\operatorname{trdeg}}
\newcommand{\Ein}{\operatorname{Ein}}
\newcommand{\Kexp}{K_{\mathrm{exp}}}
\newcommand{\Si}{\operatorname{Si}}
\newcommand{\Ci}{\operatorname{Ci}}
\newcommand{\Cov}{\operatorname{Cov}}
\newcommand{\Ex}{\mathbb E}
\setlength{\emergencystretch}{3em}
\setcounter{tocdepth}{1}
\title{Euler--Gompertz and Bessel Values:\\Joint Independence, Integral Intersections,\\and Critical Limits}
\author{Keio}
\date{22 October 2026\\\small Revised integrated research manuscript}
\begin{document}
\maketitle
\begin{abstract}
We develop joint independence and exact integral-intersection theorems
for Euler--Gompertz, exponential-integral, and Bessel values. The functional
foundation includes every depth of the Bessel primitive tower, at all
nonzero algebraic scales modulo sign, jointly with the entire
exponential-integral towers. Continued phase embeddings and local pole
charts classify full rational fields even when an exponential-integral
argument and a Bessel scale have dependent phases: the integral
intersection is the polynomial ring in the entire integral, with Laurent
polynomial coefficients in the matching exponential. Arbitrarily many
positive derivative-ratio scales can also be retained with one harmonic
Bessel coordinate. Complex independent-scale intersections yield
nonintegrality for Mahler combinations. Further consequences include exact
unit and radical criteria, a conjugation half-defect theorem, and a common
density-one set of integer arguments on which six coordinate families
are jointly independent; every integer-order second-kind Bessel value
at those arguments is transcendental. Uniform critical asymptotics and
explicit reconstruction results complete the synthesis. The arithmetic
nature of Euler's and Euler--Gompertz's constants, and full independence
of the original four constants, remain unresolved.
\end{abstract}
\clearpage
\begingroup\small
\makeatletter
\renewcommand{\l@section}{\@dottedtocline{1}{0em}{2.2em}}
\renewcommand{\l@part}{\@dottedtocline{0}{0em}{2.2em}}
\makeatother
\tableofcontents
\endgroup
\clearpage

\section{Introduction and common notation}
The starting point is the four-constant field
\begin{equation}\label{eq:four}
 H_4=k(\gamma,\delta,\kappa,C),\qquad
 3\le\trdeg_k H_4\le4,\qquad k=\overline{\mathbb Q}.
\end{equation}
Here $\gamma$ is Euler's constant and
\[
 \delta=\int_0^\infty\frac{e^{-u}}{1+u}\,du,\qquad
 I=I_0(2),\quad K=K_0(2),\quad
 \kappa=K/I,\quad C=I_1(2)/I.
\]
The lower bound in \eqref{eq:four} is proved below, not assumed as a
consequence of separate transcendence statements. Algebraic independence
always means the absence of a nonzero polynomial relation over the stated
coefficient field. In particular, a lower bound of three does not identify
a prescribed independent triple and does not decide the individual
arithmetic nature of $\gamma$ or $\delta$.

\subsection{Functions, normalizations, and coefficient rings}
For integers $r\ge1$, define the entire functions
\[
 \Ein_r(z)=\sum_{n\ge1}\frac{(-1)^{n-1}z^n}{n^r n!},
 \qquad \Ein=\Ein_1,\qquad \eta=\Ein(1)=\gamma+\delta/e.
\]
The classical tail integral is denoted by $E_1^{\rm tail}(z)$;
on compatible branches,
$\Ein(z)=\gamma+\log z+E_1^{\rm tail}(z)$.
Thus $\Ein_1$ and $E_1^{\rm tail}$ are different functions.
We use the standard modified Bessel functions $I_\nu,K_\nu$ and the
classical Bessel functions $J_\nu,Y_\nu$. Define the harmonic series
\begin{equation}\label{eq:Y}
 Y(t)=\sum_{n\ge0}\frac{H_n(t/2)^{2n}}{(n!)^2},
 \qquad H_n=\sum_{j=1}^n\frac1j,\quad H_0=0.
\end{equation}
This $Y$ is not the classical second-kind Bessel function $Y_0$.
For nonzero algebraic $t$, put
\[
 \beta(t)=\frac{Y(t)}{I_0(t)},\qquad C(t)=\frac{I_1(t)}{I_0(t)}.
\]
The nonvanishing needed in these definitions is proved before the complex
joint-value theorem. At positive real $t$, the connection and Wronskian
identities are
\begin{align}
 \beta(t)&=\gamma+\mu(t),&
 \mu(t)&=\log(t/2)+\frac{K_0(t)}{I_0(t)},\label{eq:beta}\\
 I_0(t)K_1(t)+I_1(t)K_0(t)&=\frac1t.&&\label{eq:wronskian}
\end{align}
We abbreviate $\beta=\beta(2)=\gamma+\kappa$. The following three
combinations have an exact relation that will be retained throughout:
\begin{equation}\label{eq:coordinates}
 A=\eta-\beta=\delta/e-\kappa,\qquad
 U=3\eta-2\beta,\qquad V=\beta,\qquad U=3A+V.
\end{equation}
At the imaginary scale $2i$, set
\[
 \nu=-\frac{\pi Y_0(2)}{2J_0(2)},\qquad
 C_J=\frac{J_1(2)}{J_0(2)};
 \qquad \beta(2i)=\gamma+\nu,\quad C(2i)=iC_J.
\]
The branch-free series definitions make these signs unambiguous.
The primitive coordinate and three distinct limits are
\begin{align*}
 B&=\sum_{n\ge1}\frac{(-1)^{n-1}}{n(n!)^2},&
 \mathcal L&=B-\gamma,\\
 L_+&=\gamma+2\kappa,& L_-&=\ell=\gamma-2\kappa.
\end{align*}
In particular $\mathcal L\ne L_+$; their separate critical approximation
formulas are proved in Section~\ref{sec:limits}.

Write $\Kexp=k(e^\alpha:\alpha\in k)$. Every assertion over this
infinite field means its corresponding assertion for each finite set of
coefficients. The ring $\E$ consists of values of $E$-functions at algebraic
points, and
\[
 \Ebar=\{z\in\mathbb C:z\text{ satisfies a monic polynomial over }\E\}.
\]
This is an integral closure in $\mathbb C$, not the algebraic closure
of $\operatorname{Frac}(\E)$. Since $k\subset\E$, exclusion from
$\Ebar$ implies transcendence over $k$. The proof of exclusion is
stronger than such an individual transcendence assertion.

\subsection{The closures proved by the synthesis}
The five predecessor notes are consolidated by subject, with the
functional inputs proved first and repeated arguments replaced by internal
references. The additional results proved in this integrated manuscript
are the following.
\begin{enumerate}
\item At every depth and every nonzero algebraic scale modulo sign,
Bessel primitives can be adjoined jointly to the full exponential-integral
and Bessel family over $\Kexp$ (Theorem~\ref{depth:values}).
The conjugation principle also bounds the independence defect after a
real coefficient-field extension of transcendence degree $d$ by
$\lfloor d/2\rfloor$ (Theorem~\ref{rdef:main}).
\item At arbitrary positive algebraic $a,t$,
\[
 \Ebar\cap k(e^a,\Ein(a),\beta(t),C(t))
   =k[e^a,e^{-a},\Ein(a)]
\]
(Theorem~\ref{fre:main}). In particular, the previously open full
matched-phase intersection at $(a,t)=(1,2)$ is closed.
Theorem~\ref{gde:main} strengthens this to the entire derivative family:
\[
 \Ebar\cap k\bigl(e^a,\Ein(a),\beta(s),C(t):
                    t>0\text{ algebraic}\bigr)
 =k[e^a,e^{-a},\Ein(a)]
\]
for arbitrary positive algebraic $a,s$. This includes dependent scales
such as $2$ and $4$ and arbitrary rational expressions in the generators.
\item At complex $\mathbb Q$-linearly independent algebraic scales,
the full rational field in all $\beta(t_i),C(t_i)$ has integral
intersection exactly $k$ (Theorem~\ref{cxi:main}). This includes
the pair of scales $a,ia$ and its gamma-free combinations.
\item At nonzero $\mathbb Q$-linearly independent algebraic arguments,
the rational field in $e^{a_i},\Ein(a_i)$ and additional independent
exponential units, whose exponents are jointly $\mathbb Q$-linearly
independent with the $a_i$, has exactly the Laurent--polynomial intersection
of Theorem~\ref{ner:main}. The one-point consequence is
\[
 \Ebar\cap k(e,\gamma+\delta/e)
       =k[e,e^{-1},\gamma+\delta/e].
\]
\item A single set of at least $N-\pi(N)-1$ integers in $[1,N]$
supports joint independence of all six coordinates in
Theorem~\ref{jden:density}. At every retained integer, all integer-order
$K_r$ and $\pi Y_r$ values are transcendental over $\Kexp$.
Every positive algebraic argument has a six-coordinate rank bound of
at least five, without an exceptional argument.
\item If $\trdeg_kH_4=3$, for algebraic $r>0$ with $r\ne1$, the geometric family $K_0(2r^n)$ loses
at most one coordinate over $\Kexp H_4$, and integer arguments have
the improved lower bound $N-\pi(N)$. These are conditional strengthenings
of the unconditional density-one theorem.
\item The exact intersections give unit and radical classifications
(Section~\ref{unit:section}), including explicit integral units that
are not themselves E-values. The upper-half-plane $\Ein_r$ family remains independent after
adjoining $\gamma$, all real-argument $\Ein_r$ values, and all complex
Bessel and primitive coordinates to the exponential coefficient field.
\end{enumerate}
The distinction between pairwise distinct scales and
$\mathbb Q$-linearly independent scales remains important. The joint-value
theorem needs only distinctness modulo sign. Independent phase variables
classify the general multiscale Bessel rational field; synchronized local
charts remove that restriction for the one-harmonic-coordinate families
identified above. Several harmonic coordinates at dependent scales
remain a separate question.

\subsection{External inputs and status}
The analytic and differential-algebraic foundations are the classical
Bessel differential Galois group, Picard--Vessiot theory, ordinary-point
Siegel--Shidlovskii specialization \cite{Beukers,Kolchin}, and the
phase-embedding construction of Fischler--Rivoal \cite{FR26}.
The latter construction and the continuation step used here are stated
and justified explicitly in the integrality part. Standard connection
identities are recorded in \cite{DLMF}. Mahler's classical one-point
transcendence theorem \cite{Mahler} is part of the prior context.
The predecessor notes and supplied corpus are research manuscripts;
the present text proves extensions of their identified arguments and
makes no claim of a comprehensive historical priority determination.
The source concordance at the end records which results were merged,
strengthened, or retained as separate applications.


\part{Foundations}
% PART 01: proved foundations and simultaneous primitive extension.
% No preamble or bibliography. Labels pj:joint and cf:joint alias one theorem.
\section{Polyexponential towers and their solvable closure}
\label{pj:section}

For integers $r\geq1$, set
\[
 \Ein_r(z):=\sum_{n\geq1}\frac{(-1)^{n-1}z^n}{n^r n!},
 \qquad \Ein_0(z)=1-e^{-z}.
\]
Thus $\Ein_1$ in this section denotes the entire complementary exponential
integral, not the classical tail integral $\int_z^\infty e^{-u}du/u$.
Write $D=d/dz$ and $\vartheta=zD$; then
$\vartheta \Ein_r(\lambda z)=\Ein_{r-1}(\lambda z)$.
For every fixed $r$ and algebraic $\lambda$, this is an E-function:
a power of a denominator of $\lambda$ times
$\operatorname{lcm}(1,\ldots,N)^r$ clears its first $N$ normalized Taylor
coefficients, with exponential growth in $N$.

We first prove the primitive-tower lemma and identify its solvable
Picard--Vessiot closure. We then establish the full Bessel product
group at complex scales, before applying the classical specialization
theorem. The classical inputs are Picard--Vessiot theory and
Siegel--Shidlovskii specialization \cite{Kolchin,Beukers}.

\begin{lemma}[Primitive tower and its logarithm]\label{pj:primitive}
Let $\lambda_1,\ldots,\lambda_s$ be distinct nonzero algebraic numbers,
let $\nu_1,\ldots,\nu_d$ be a $\mathbb Z$-basis of their generated
additive group, and put
\[
 K_0=\mathbb C(z)(e^{-\nu_1z},\ldots,e^{-\nu_dz}),\qquad
 K_q=K_0\bigl(\Ein_r(\lambda_i z):1\le i\le s,\ 1\le r\le q\bigr).
\]
For every $q\ge0$ the constant field is $\mathbb C$,
$1/z\notin D(K_q)$, and
\[
 \operatorname{trdeg}_{\mathbb C(z)}K_q=d+qs.
\]
Moreover $\log z$ is transcendental over every $K_q$.
\end{lemma}
\begin{proof}
Write $X_j=e^{-\nu_jz}$ and $\lambda_i=m_i\cdot\nu$.
The rationally independent exponential rates identify $K_0$ with the
fraction field of the Laurent polynomial ring
$\mathbb C(z)[X_1^{\pm1},\ldots,X_d^{\pm1}]$,
with $DX_j=-\nu_jX_j$.
There is no irreducible nonunit $P$ in this ring with $P\mid DP$.
Indeed, if $DP=AP$, comparison of Newton polytopes gives
$A\in\mathbb C(z)$; comparison of two nonzero coefficients $a_u,a_v$
of $P$ would give
\[
 \frac{(a_u/a_v)'}{a_u/a_v}=(u-v)\cdot\nu.
\]
A rational function has no nonzero constant logarithmic derivative,
so $P$ would have a single Laurent monomial and be a unit.
Consequently differentiation of a genuine nonunit denominator raises
its pole order and cannot yield a Laurent polynomial. If
\[
 Dg=z^{-1}\sum_u c_uX^u,\qquad g\in K_0,\quad c_u\in\mathbb C,
\]
then $g$ is a Laurent polynomial. Coefficient comparison gives
$a_u'-(u\cdot\nu)a_u=c_u/z$ with $a_u\in\mathbb C(z)$.
For $u=0$ a rational derivative has zero residue. For $u\ne0$ a
pole of $a_u$ would create a higher-order pole, whereas absence of a
pole at zero makes the left side regular there. Hence all $c_u=0$.
In particular $1/z$ is not exact in $K_0$, and no nonzero constant
combination of $(1-X^{m_i})/z$ is exact there.

We use the following elementary primitive-field observation. If
$a_1,\ldots,a_s\in F$ have no nonzero constant combination in $D(F)$,
and algebraically independent primitives $u_i$ satisfy $Du_i=a_i$,
then
\[
 g\in F(u_1,\ldots,u_s),\quad Dg\in F
 \quad\Longrightarrow\quad
 g=f+\sum_i b_i u_i\quad(f\in F,\ b_i\in\mathbb C).
\]
To see this, a normalized nonconstant Darboux polynomial in the $u_i$
would have derivative zero; its highest homogeneous part has constant
coefficients, and comparison of the next degree would make a nonzero
constant combination of the $a_i$ exact in $F$. Thus no such polynomial
exists and differentiation excludes nonconstant denominators in $g$.
The same comparison of the top two homogeneous degrees of a polynomial
$g$ leaves degree at most one and proves the displayed assertion.
The observation also shows that no new constants are introduced.

Kolchin--Ostrowski applied to
$D\Ein_1(\lambda_i z)=(1-X^{m_i})/z$ gives independence of the first
primitive layer. The displayed observation and the Laurent argument
show that $1/z$ is still not exact in $K_1$.
Inductively, the last layer of $K_q/K_{q-1}$ admits every translation
$\Ein_q(\lambda_i z)\mapsto \Ein_q(\lambda_i z)+h_i$.
If $Dg=\sum_i c_i\Ein_q(\lambda_i z)/z$ with $c\ne0$ and $g\in K_q$,
subtracting a translate of this identity would make a nonzero
constant multiple of $1/z$ exact in $K_q$, a contradiction.
Kolchin--Ostrowski therefore adds $s$ independent primitives at the
next layer. If $Dg=1/z$ in $K_{q+1}$, the primitive-field observation
writes $g=f+\sum_i b_i\Ein_{q+1}(\lambda_i z)$. Translating the preceding
layer gives the same contradiction when $b\ne0$; when $b=0$ it
contradicts the preceding layer directly. This proves the induction.
All fields are realized as meromorphic functions on a fixed slit
plane, so their constants are indeed $\mathbb C$.

Finally, a primitive algebraic over a differential field of
characteristic zero lies in that field up to a constant: take the
normalized field trace of the primitive. Applied to $D\log z=1/z$,
this would contradict the nonexactness just proved. Thus adjoining
$\log z$ raises the transcendence degree by one.
\end{proof}

\begin{lemma}[Explicit solvable Picard--Vessiot closure]\label{pj:pv}
For $R\ge1$, the field $L=K_R(\log z)$ is a Picard--Vessiot extension
of $\mathbb C(z)$ with solvable differential Galois group and
\[
 \operatorname{trdeg}_{\mathbb C(z)}L=d+Rs+1.
\]
\end{lemma}
\begin{proof}
For each $i$, consider the order-$R+2$ operator
\[
 \mathcal L_i=(D+\lambda_i)D\vartheta^R.
\]
A fundamental solution basis is
\[
 1,\log z,\ldots,(\log z)^R,\Ein_R(\lambda_i z).
\]
Indeed, $D\vartheta^R \Ein_R(\lambda_i z)=\lambda_i e^{-\lambda_i z}$,
and $\Ein_R(\lambda_i z)$ cannot be a constant-coefficient polynomial in
$\log z$, as its $R$th $\vartheta$-derivative is $1-e^{-\lambda_i z}$.
The entries of a fundamental matrix and their derivatives generate
exactly
\[
 \mathbb C(z)\bigl(e^{-\lambda_i z},\log z,
                   \Ein_1(\lambda_i z),\ldots,\Ein_R(\lambda_i z)\bigr).
\]
Their compositum for all $i$ is $L$, so it is the PV field of the
direct sum of these operators.

Its Galois group preserves the flag of polynomials in $\log z$ and
acts triangularly on each displayed solution basis. More explicitly,
if $e^{-\lambda_i z}\mapsto c_i e^{-\lambda_i z}$, then
\[
 \Ein_R(\lambda_i z)\mapsto
 c_i \Ein_R(\lambda_i z)
 +(1-c_i)\frac{(\log z)^R}{R!}
 +\text{a polynomial in $\log z$ of degree below $R$}.
\]
Also $\log z\mapsto\log z+c$. Hence the group is a subgroup of a
product of triangular groups and is solvable. Its dimension is the
transcendence degree given by Lemma~\ref{pj:primitive}.
\end{proof}


\section{Complex Bessel scales and joint specialization}
\label{cf:section}

Throughout this section $K=\mathbb C(z)$ and $D=d/dz$. All functions are
realized as meromorphic functions on a fixed slit plane, with one choice of
$\log z$; the constant field is $\mathbb C$. For $t\ne0$ write
\[
 f_t(z)=I_0(tz),\qquad
 y_t(z)=Y(tz),\qquad
 Y(w)=\sum_{n\ge0}\frac{H_n(w/2)^{2n}}{(n!)^2},
 \quad H_0=0.
\]
The scale $t$ may be complex. Coefficient comparison gives
\[
 (D^2+z^{-1}D-t^2)f_t=0,\qquad
 (D^2+z^{-1}D-t^2)y_t=2f_t'/z.
\]
Consequently
\begin{equation}\label{cf:second}
 g_t(z)=y_t(z)-f_t(z)\log z
\end{equation}
is a second homogeneous solution. Its Wronskian with $f_t$ is
\begin{equation}\label{cf:W}
 f_tg_t'-f_t'g_t=-1/z.
\end{equation}
Indeed, Abel's identity makes the Wronskian a constant times $1/z$,
and $f_t=1+O(z^2)$, $g_t=-\log z+O(z^2\log z)$ determine the constant.
This normalization avoids any branch convention for $K_0(tz)$ or $\log t$.

\begin{lemma}[The complex-scale product group]\label{cf:product}
Let $t_1,\ldots,t_m\in\mathbb C^*$ satisfy $t_i\ne\pm t_j$ for $i\ne j$.
The compositum of the solution fields of
$D^2+z^{-1}D-t_j^2$ has differential Galois group
$\mathrm{SL}_2(\mathbb C)^m$ over $K$.
\end{lemma}
\begin{proof}
The individual order-zero Bessel group is $\mathrm{SL}_2$ by the classical
Bessel theorem \cite{Kolchin}, unchanged by the nonzero linear scale change.
The full group $G$ is therefore a subdirect algebraic subgroup of
$\mathrm{SL}_2^m$. Its identity component projects onto every factor:
the image is connected and normal, with finite quotient in the connected
group $\mathrm{SL}_2$.

The Lie algebra of $G^\circ$ is a subdirect subalgebra of
$\mathfrak{sl}_2^m$. Its solvable radical maps to a solvable ideal of each
factor and is therefore zero. Each simple ideal maps either to zero or
isomorphically onto a factor, and exactly one simple ideal maps onto each
factor, since distinct ideals commute. If $G$ is proper, a simple ideal
must map onto two different factors. Their adjoint representations are
then isomorphic as representations of the connected group $G^\circ$.
By the PV correspondence, the two adjoint differential modules become
isomorphic over the finite algebraic fixed field of $G^\circ$.

This is impossible for two scales $a\ne\pm b$. The rank-two system matrix
is
\[
 A_t(z)=\begin{pmatrix}0&1\\t^2&-1/z\end{pmatrix}.
\]
The adjoint connection has distinct leading eigenvalues
$0,2t,-2t$ at infinity. It can be formally diagonalized over
$\mathbb C((z^{-1}))$: at each successive order, the off-diagonal
coefficients are determined by division by nonzero eigenvalue differences.
After a finite algebraic extension, choose a place over infinity and embed
its completion in $\mathbb C((z^{-1/e}))$. An entry $h$ of a differential
intertwiner between the two formally diagonalized modules satisfies
\[
 h'=(\mu-\lambda+O(z^{-1}))h,
 \quad \lambda\in\{0,2a,-2a\},\quad\mu\in\{0,2b,-2b\}.
\]
For a nonzero Puiseux Laurent series $h$, its logarithmic derivative is
$O(z^{-1})$. Thus a nonzero entry requires $\lambda=\mu$. An invertible
intertwiner requires equality of the two eigenvalue multisets, forcing
$a=\pm b$, a contradiction. Hence the group is the full product.
\end{proof}

\begin{corollary}[Distinguished columns and nonvanishing]\label{cf:nonzero}
For nonzero algebraic $t_1,\ldots,t_m$ distinct modulo sign, the $2m$ values
$I_0(t_j),I_1(t_j)$ are algebraically independent over
$\overline{\mathbb Q}$. In particular $I_0(t)\ne0$ at every nonzero
algebraic $t$, so $C(t)=I_1(t)/I_0(t)$ and $\beta(t)=Y(t)/I_0(t)$ are defined.
\end{corollary}
\begin{proof}
The product group moves the first solution column of each rank-two
fundamental matrix independently through all nonzero linear combinations
of its two columns. A polynomial relation among $f_{t_j},f_{t_j}'$
would therefore vanish on a product of Zariski-dense punctured affine
planes. The nonzero Wronskians make the changes of coordinates invertible,
so no such relation exists. These $2m$ functions form a first-order
E-function system with only the finite singularity $z=0$.
Siegel--Shidlovskii at $z=1$ proves independence of
$f_{t_j}(1),f_{t_j}'(1)$; the latter derivatives are
$t_jI_1(t_j)$, with nonzero algebraic factors $t_j$.
Nonvanishing follows from independence.
\end{proof}

\begin{theorem}[Joint theorem at arbitrary complex algebraic scales]
\label{cf:joint}\label{pj:joint}
Let $\lambda_1,\ldots,\lambda_s$ be distinct nonzero algebraic numbers,
$R\ge1$, and let $t_1,\ldots,t_m$ be nonzero algebraic numbers distinct
modulo sign. Put $d=\dim_{\mathbb Q}\langle\lambda_1,\ldots,\lambda_s\rangle$.
Then
\[
 \operatorname{trdeg}_{\overline{\mathbb Q}}
 \overline{\mathbb Q}\bigl(e^{\lambda_i},
 \Ein_r(\lambda_i),I_0(t_j),I_1(t_j),Y(t_j):
 i\le s,r\le R,j\le m\bigr)=d+Rs+3m.
\]
The $Rs+3m$ nonexponential coordinates are algebraically independent over
the field of the displayed exponential coordinates. Moreover, after
choosing one representative from each nonzero algebraic sign class,
all the coordinates
\[
 \Ein_r(\lambda)\quad
 (r\ge1,\ \lambda\in\overline{\mathbb Q}^{\times}),
 \qquad I_0(t),C(t),\beta(t)
\]
are jointly algebraically independent over
\[
 K_{\exp}=\overline{\mathbb Q}
             (e^\alpha:\alpha\in\overline{\mathbb Q}),
\]
in the finite-subset sense. The exponential coordinates themselves are
not part of an independent family over $K_{\exp}$.
\end{theorem}
\begin{proof}
The primitive-tower and solvable-closure lemmas
(Lemmas~\ref{pj:primitive} and~\ref{pj:pv}) give a solvable PV field
\[
 S=K\bigl(e^{\lambda_i z},\Ein_r(\lambda_i z),\log z:
                         i\le s,r\le R\bigr)
\]
of transcendence degree $d+Rs+1$ over $K$. Those proofs allow arbitrary
complex algebraic $\lambda_i$. Let $B$ be the full Bessel PV field from
Lemma~\ref{cf:product}. A quotient common to its connected perfect group
$\mathrm{SL}_2^m$ and the solvable group of $S$ is trivial. Goursat's lemma
therefore makes the group of $BS$ the direct product, and
\[
 \operatorname{trdeg}_K BS=3m+d+Rs+1.
\]
More explicitly, before specialization put
\[
 F=K\bigl(e^{\lambda_i z},\Ein_r(\lambda_i z),
 I_0(t_jz),I_1(t_jz),Y(t_jz):i\le s,r\le R,j\le m\bigr).
\] Equations~\eqref{cf:second} and~\eqref{cf:W} recover the
second Bessel solutions and their derivatives after adjoining $\log z$;
conversely all displayed functions lie in $BS$. Thus $BS=F(\log z)$.
It follows that $\operatorname{trdeg}_K F\ge3m+d+Rs$; counting generators
with the exponential character lattice gives the reverse inequality.

All these entire functions are E-functions. Include $1$, both exponential
signs, the complete polyexponential chains, and
$f_{t_j},f_{t_j}',y_{t_j},y_{t_j}'$ in one first-order system. Its matrix has
no finite pole except $z=0$. The Wronskian calculation also gives
\[
 I_1(w)Y(w)-I_0(w)Y'(w)=\frac{1-I_0(w)^2}{w}.
\]
Corollary~\ref{cf:nonzero}, proved independently before forming quotients,
ensures that this identity expresses $Y'(t_j)$ in the other value
coordinates. Hence the generated value field is exactly the one in the
statement. Siegel--Shidlovskii at the ordinary point $z=1$ gives its
transcendence degree \cite{Beukers}. The exponential value field has transcendence degree $d$ by
Lindemann--Weierstrass: distinct integer characters of a rationally
independent algebraic exponent basis give distinct algebraic exponents.
This proves the asserted relative independence. The transformations
$(I_0,I_1,Y)\leftrightarrow(I_0,C,\beta)$ are birational by nonvanishing.

Finally, any relation over $K_{\exp}$ uses only finitely many algebraic
exponents in its coefficients. Add them to the finite $\lambda$ list
(omitting zero), enlarge $R$ if necessary, and apply the finite theorem.
This proves the infinite-base statement. Distinctness modulo sign is
necessary: $I_0(-t)=I_0(t)$, $Y(-t)=Y(t)$, and $C(-t)=-C(t)$.
\end{proof}

\section{Simultaneous primitives at every Bessel scale}
\label{sec:msprim}
Write $k=\overline{\mathbb Q}$ and
$K_{\exp}=k(e^\alpha:\alpha\in k)$. For $t\ne0$, let
\begin{equation}\label{msprim:H}
 H_t(z)=-\sum_{n\ge1}\frac{(tz/2)^{2n}}{2n(n!)^2},
 \qquad H_t'(z)=\frac{1-I_0(tz)}z.
\end{equation}
Thus $H_{-t}=H_t$. The conditions on distinct scales below are
therefore necessary even before considering algebraic independence.

\begin{lemma}[No exact linear combination of Bessel derivatives]
\label{msprim:nonexact}
Let $K=\mathbb C(z)$. Let $P/K$ be the full Bessel PV field at
nonzero scales $s_1,\ldots,s_n$ distinct up to sign, with group
$G_B=\mathrm{SL}_2^n$. Let $S/K$ be a solvable PV field containing
$\log z$, and put $L=PS$. Suppose
$\operatorname{Gal}(L/K)=G_B\times G_S$.
Choose any subset of scales $t_1,\ldots,t_m$ among the $s_j$.
If $c_1,\ldots,c_m\in\mathbb C$ are not all zero, no $u\in L$
satisfies
\begin{equation}\label{msprim:derivative}
 u'=\sum_{j=1}^m c_j\frac{I_0(t_jz)}z.
\end{equation}
\end{lemma}
\begin{proof}
For each active index $j$ with $c_j\ne0$, let $V_j$ be the solution
space of $y''+z^{-1}y'-t_j^2y=0$. These are pairwise nonisomorphic
irreducible $G_B$-modules: the $j$th $\mathrm{SL}_2$ factor acts by
its standard representation, and all the other factors act trivially.
The orbit of the right side of \eqref{msprim:derivative} spans
\[
 V=\bigoplus_{c_j\ne0} V_j/z.
\]
Indeed each nonzero component generates its irreducible summand,
and the summands are nonisomorphic.

If $u$ existed, the vector space
\[
 U=\operatorname{span}_{\mathbb C}\{1,\sigma(u):
                   \sigma\in G_B\times G_S\}
\]
would be finite-dimensional, with differentiation giving the exact
sequence $0\to\mathbb C\to U\to V\to0$. Any finite-dimensional
Galois-stable subspace of a PV field is a rational representation
of its differential Galois group.
Complete reducibility for $G_B$ supplies a complement $U_0$ to the
constants. It is unique since
$\operatorname{Hom}_{G_B}(V,\mathbb C)=0$.
The commuting group $G_S$ preserves $U_0$ and fixes it pointwise:
it fixes $V$, and differentiation is injective on $U_0$.
Thus $U_0\subset L^{G_S}=P$.

Projecting $U_0$ onto any active standard summand yields a row
$(u_f,u_g)$ of primitives of $(f/z,g/z)$ for a fundamental pair
$f,g$ of the selected Bessel equation. This row transforms in the
same way as $(f,g)$. Hence
\[
 (a,b)=(u_f,u_g)
       \begin{pmatrix}f&g\\f'&g'\end{pmatrix}^{-1}
\]
is invariant under the entire group $G_B$, so $a,b\in K$.
Differentiation gives
\[
 a'+t_j^2b=1/z,\qquad a+b'-b/z=0,
\]
and therefore
\begin{equation}\label{msprim:rational}
 -b''+b'/z-b/z^2+t_j^2b=1/z.
\end{equation}
A rational $b$ cannot have a nonzero finite pole: the highest-order
pole from $-b''$ could not cancel. Thus it is a Laurent polynomial.
If its lowest nonzero term is $dz^r$ with $r\le0$, the left side
has leading term $-(r-1)^2dz^{r-2}$, impossible on the right side.
If instead $r\ge1$, the left side is holomorphic at zero, since
its potential $z^{-1}$ term at $r=1$ has coefficient zero.
This also contradicts \eqref{msprim:rational}. The identically zero
choice of $b$ is plainly impossible. This proves the lemma.
\end{proof}

\begin{theorem}[Joint independence of all selected Bessel primitives]
\label{msprim:functional}
Under the hypotheses of Lemma~\ref{msprim:nonexact}, the functions
$H_{t_1}(z),\ldots,H_{t_m}(z)$ are algebraically independent over $L$.
\end{theorem}
\begin{proof}
Put $u_j=\log z-H_{t_j}(z)$, so $u_j'=I_0(t_jz)/z\in L$.
The field $M=L(u_1,\ldots,u_m)$ is a PV extension of $L$ for the
triangular system whose fundamental solution matrix is
\[
 \begin{pmatrix}1&u_1&\cdots&u_m\\
 0&1&\cdots&0\\[-2pt]
 \vdots&\vdots&\ddots&\vdots\\
 0&0&\cdots&1\end{pmatrix}.
\]
All fields are realized as fields of analytic germs on a common
simply connected domain, so their constant field is $\mathbb C$.
The PV group is an algebraic subgroup of $\mathbb G_a^m$, acting
by translations of the $u_j$.

In characteristic zero, algebraic subgroups of $\mathbb G_a^m$ are
linear subspaces. If this group were proper, some nonzero vector
$(c_1,\ldots,c_m)$ would annihilate all its translations. The
Galois correspondence would then give
$\sum_jc_ju_j\in L$, contradicting
Lemma~\ref{msprim:nonexact} after differentiation.
Consequently the group is $\mathbb G_a^m$, and
$\operatorname{trdeg}_L M=m$. Since $\log z\in L$, the same joint
independence holds for the $H_{t_j}$.
\end{proof}

\begin{theorem}[All-scale arithmetic primitive extension]
\label{msprim:values}
Let $s_1,\ldots,s_n\in k^\times$ be distinct up to sign, and select
any subset $t_1,\ldots,t_m$ of these scales. Let
$\lambda_1,\ldots,\lambda_q\in k^\times$ be distinct, with positive
integers $R_1,\ldots,R_q$. Then the
$\sum_iR_i+3n+m$ numbers
\[
 \Ein_r(\lambda_i)\quad(1\le i\le q,\ 1\le r\le R_i),
\]
\[
 I_0(s_j),\quad C(s_j)=\frac{I_1(s_j)}{I_0(s_j)},\quad
 \beta(s_j)=\frac{Y(s_j)}{I_0(s_j)}\quad(1\le j\le n),
 \qquad H_{t_l}(1)\quad(1\le l\le m)
\]
are algebraically independent over $K_{\exp}$.
Here $Y(s)=\sum_{n\ge0}H_n(s/2)^{2n}/(n!)^2$ is the harmonic
Bessel series, with $H_0=0$.
\end{theorem}
\begin{proof}
Any relation over $K_{\exp}$ uses only finitely many exponential
values. Include their corresponding functions in the solvable PV
field of the all-order polyexponential construction, together with
$\log z$. The Bessel product field and this solvable field have a
direct-product group: a connected product of $\mathrm{SL}_2$ groups
has no nontrivial solvable quotient. The previously established
E-function coordinates lie in their composite $L$.
Theorem~\ref{msprim:functional} adds exactly $m$ independent
functions to these coordinates.

For algebraic $t$, the coefficient of $z^{2n}/(2n)!$ in $H_t$ is
$-\binom{2n}{n}(t/2)^{2n}/(2n)$; its conjugate-growth and
common-denominator bounds are exponential. Together with
\eqref{msprim:H}, this verifies that $H_t$ is an E-function.
Appending its first-order equation to the old system introduces no
finite singularity other than zero. Siegel--Shidlovskii specialization
at $z=1$ therefore preserves the functional transcendence-degree
increase. Theorem~\ref{cf:joint} supplies the independence of all the other
coordinates, and passage to the displayed Bessel quotients is
birational because $I_0(s_j)\ne0$ at nonzero algebraic $s_j$
by Corollary~\ref{cf:nonzero}.
Since the finite set of exponential coefficients was arbitrary,
the assertion holds over $K_{\exp}$.
\end{proof}

Equivalently, choose one representative for every sign class in
$k^\times/\{\pm1\}$. The entire union of the polyexponential values,
Bessel triples, and primitive values indexed as above is independent
over $K_{\exp}$, in the sense of every finite subfamily. One must not
include both $H_t(1)$ and $H_{-t}(1)$ as distinct coordinates.

\subsection{Critical limits and matched ordinary-Bessel combinations}
For $\lambda>0$ put
\[
 B(\lambda)=\sum_{n\ge1}\frac{(-1)^{n-1}\lambda^n}{n(n!)^2}
           =2H_{2i\sqrt\lambda}(1).
\]
The critical Newton--Laguerre formula is
\[
 \mathcal L(\lambda)=B(\lambda)-\gamma-\log\lambda,
 \qquad
 R_N(\lambda/N)-\log(\lambda/N)\longrightarrow\mathcal L(\lambda).
\]
Here $R_N(x)=-\sum_{n=1}^N L_n(x)/n$, with
$L_n(x)=\sum_{j=0}^n(-1)^j\binom njx^j/j!$.
The exact finite-sum identity and dominated-convergence proof are given
in Proposition~\ref{prim:limit}.
For positive algebraic $\lambda$ define
\[
 j_\lambda=J_0(2\sqrt\lambda),\qquad
 c_\lambda=\frac{J_1(2\sqrt\lambda)}{J_0(2\sqrt\lambda)},\qquad
 \nu_\lambda=-\frac{\pi Y_0(2\sqrt\lambda)}{2J_0(2\sqrt\lambda)},
\]
\begin{equation}\label{msprim:betaordinary}
 \beta_\lambda:=\beta(2i\sqrt\lambda)
 =\gamma+\frac12\log\lambda+\nu_\lambda
 =\frac{\displaystyle\sum_{n\ge0}
              (-1)^nH_n\lambda^n/(n!)^2}{j_\lambda}.
\end{equation}
The factor $1/2$ in front of $\log\lambda$ is essential.

\begin{corollary}[Four independent coordinates per critical scale]
\label{msprim:matched}
For distinct positive algebraic $\lambda_1,\ldots,\lambda_m$, the
$4m$ numbers
\[
 j_{\lambda_j},\quad c_{\lambda_j},\quad
 \beta_{\lambda_j},\quad B(\lambda_j)
 \qquad(1\le j\le m)
\]
are jointly independent over $K_{\exp}$. One may replace
$B(\lambda_j)$ by
\begin{equation}\label{msprim:Q}
 Q_{\lambda_j}:=B(\lambda_j)-\beta_{\lambda_j}
 =\mathcal L(\lambda_j)+\tfrac12\log\lambda_j-\nu_{\lambda_j}.
\end{equation}
In particular, the $Q_{\lambda_j}$ are jointly independent over the
field generated by $K_{\exp}$, all the old polyexponential values,
and all the old Bessel coordinates, including those at their own
scales. Their series-and-quotient expression $B(\lambda)-\beta_\lambda$
contains no explicit $\gamma$ or logarithm.
\end{corollary}
\begin{proof}
The scales $2i\sqrt{\lambda_j}$ are nonzero algebraic and distinct
up to sign. Apply Theorem~\ref{msprim:values}, using
$I_0(2i\sqrt\lambda)=j_\lambda$ and
$C(2i\sqrt\lambda)=i c_\lambda$. The substitution
$B(\lambda_j)\mapsto B(\lambda_j)-\beta_{\lambda_j}$ is an invertible
linear change over the old coordinate field. Extra old coordinates
in any proposed relation are finite in number and may be included
in the same theorem.
\end{proof}

\begin{corollary}[A logarithm-cancelled one-defect family]
\label{msprim:W}
Let
\[
 W_\lambda:=\mathcal L(\lambda)-2\nu_\lambda
 =\gamma+B(\lambda)-2\beta_\lambda
 \qquad(\lambda\in k\cap\mathbb R_{>0}).
\]
Among all the numbers
\[
 \gamma,\quad\delta,\quad\kappa,\quad\gamma-2\kappa,\quad
 \nu_1,\quad\gamma-2\nu_1,\quad
 W_\lambda\quad(\lambda\in k\cap\mathbb R_{>0}),
\]
at most one is algebraic over $K_{\exp}$, and hence at most one is
algebraic over $k$.
\end{corollary}
\begin{proof}
Put $\eta=\Ein(1)$, $\beta_K=\gamma+\kappa$, and
$\beta_J=\gamma+\nu_1$. The previously known part follows from the
joint independence of $\eta,\beta_K,\beta_J$ over $K_{\exp}$:
normalize each displayed form as a nonzero $K_{\exp}$-multiple of
$\gamma$ plus a linear combination of these three independent
coordinates. No two normalized forms have the same coefficient
vector. Two algebraic values would give a nonzero linear combination
of independent coordinates algebraic over $K_{\exp}$.

For a pair involving $W_\lambda$, use its expression
$\gamma+B(\lambda)-2\beta_\lambda$ and eliminate $\gamma$ in the
same way. Against any old form, the coefficient of $B(\lambda)$
is nonzero. For two distinct $W_\lambda,W_\mu$, their difference
has nonzero coefficients of the independent primitive coordinates
$B(\lambda)$ and $B(\mu)$. Theorem~\ref{msprim:values} excludes
algebraicity in every case. Possible equality of $\beta_\lambda$
with an already selected Bessel coordinate causes no difficulty:
the new primitive coefficient still does not vanish.
\end{proof}

The unshifted limits $\mathcal L(\lambda)$ for $\lambda\ne1$ should
not be inserted into a one-defect count without accounting for their
$\log\lambda$ terms. The combinations in \eqref{msprim:Q} and
Corollary~\ref{msprim:W} make that normalization explicit.

\begin{corollary}[A relation over the enlarged old field forces rank four]
\label{msprim:rankcriterion}
Let $H_4=k(\gamma,\delta,\kappa,C)$, and let
\[
 F_{\mathrm{old}}=K_{\exp}\bigl(
 \Ein_r(\lambda),\ I_0(t),\ C(t),\ \beta(t):
 \lambda,t\in k^\times,\ r\ge1\bigr),
\]
with duplicate sign-related Bessel coordinates omitted.
If $\operatorname{trdeg}_kH_4=3$, the entire family
$\{H_t(1):t\in k^\times/\{\pm1\}\}$ is algebraically independent
 over $H_4F_{\mathrm{old}}$.
Consequently any nonzero polynomial relation among finitely many
of these primitive values, with coefficients in $H_4F_{\mathrm{old}}$,
would imply $\operatorname{trdeg}_kH_4=4$.
The same criterion applies to finitely many distinct $Q_\lambda$
or $W_\lambda$ at positive algebraic scales.
\end{corollary}
\begin{proof}
The field
\[
 F_0=k(e,\eta,I_0(2),C,\beta(2))
 \subset H_4(e,I_0(2))
\]
has transcendence degree five by the basic joint theorem.
If $\operatorname{trdeg}_kH_4=3$, the field on the right has
transcendence degree at most five. Therefore
$H_4(e,I_0(2))/F_0$ is algebraic, so
$H_4F_{\mathrm{old}}/F_{\mathrm{old}}$ is algebraic as well.
Theorem~\ref{msprim:values} gives joint independence of the primitive
family over $F_{\mathrm{old}}$; algebraic extension of the coefficient
field preserves this independence. The known bound
$3\le\operatorname{trdeg}_kH_4\le4$ proves the criterion.
Finally $Q_\lambda$ and $W_\lambda$ are invertible affine transforms
of $B(\lambda)=2H_{2i\sqrt\lambda}(1)$ over the indicated enlarged
coefficient field.
\end{proof}
No such relation is asserted to exist. The criterion identifies what
an additional relation would accomplish; it does not furnish one.

% All-depth extension of the simultaneous Bessel primitive theorem.
% This fragment uses cf:joint, cf:nonzero, and msprim:nonexact.
\section{Bessel primitive towers of every depth}
\label{depth:section}

The primitive extension admits arbitrary finite depths at all selected
scales.  The extra logarithmic weights require a nonexactness argument;
independence of the first primitives alone would not justify this step.
For $t\ne0$ and integers $r\ge1$ define
\begin{equation}\label{depth:definition}
 H_{t,r}(z)=-\sum_{n\ge1}
       \frac{(tz/2)^{2n}}{(2n)^r(n!)^2},\qquad
 H_{t,0}(z)=1-I_0(tz).
\end{equation}
With $\theta=z\,d/dz$, these functions satisfy
$\theta H_{t,r}=H_{t,r-1}$.  In particular $H_{t,1}=H_t$ in
\eqref{msprim:H}, and $H_{-t,r}=H_{t,r}$.

\subsection{Nonexactness with polynomial logarithmic weights}

\begin{lemma}[Weighted Bessel nonexactness]\label{depth:nonexact}
Let $K=\mathbb C(z)$, let $P/K$ be the full Bessel Picard--Vessiot
field at nonzero scales $s_1,\ldots,s_n$ distinct modulo sign, and
let $S/K$ be a solvable Picard--Vessiot field containing
$l=\log z$.  Assume that $L=PS$ has differential Galois group
$\mathrm{SL}_2^n\times\operatorname{Gal}(S/K)$.
For any polynomials $A_1(T),\ldots,A_n(T)\in\mathbb C[T]$, not all
zero, there is no $u\in L$ satisfying
\begin{equation}\label{depth:weighted-derivative}
 u'=\sum_{i=1}^n A_i(l)\frac{I_0(s_i z)}z.
\end{equation}
\end{lemma}

\begin{proof}
Write $G_B=\mathrm{SL}_2^n$.  The group $G_B$ fixes $S$ and hence
$l$.  For each active index $i$, let $V_i$ be the standard
two-dimensional solution space of the corresponding Bessel equation.
The orbit of the right side of \eqref{depth:weighted-derivative}
under $G_B$ spans the direct sum
\[
       \bigoplus_{A_i\ne0} A_i(l)V_i/z.
\]
The active summands are nonisomorphic irreducible $G_B$-modules.
If $u$ existed, the span of $1$ and its $G_B$-orbit would be
finite-dimensional: differentiation maps this span into the displayed
finite-dimensional space, with kernel consisting of constants.
Complete reducibility supplies, for each active scale $t=s_i$, a
pair of primitives $(u_f,u_g)$ transforming like a fundamental Bessel
solution pair $(f,g)$ and satisfying
\[
             u_f'=A_i(l)f/z,\qquad u_g'=A_i(l)g/z.
\]
Here finite-dimensional Galois-stable subspaces of a Picard--Vessiot
field carry rational representations.  The invariant row
\[
 (a,b)=(u_f,u_g)
       \begin{pmatrix}f&g\\f'&g'\end{pmatrix}^{-1}
\]
belongs to $L^{G_B}=S$.  Differentiating gives
\[
 a'+t^2b=A_i(l)/z,\qquad a+b'-b/z=0,
\]
and therefore
\begin{equation}\label{depth:b-equation}
 \mathcal D_t b:=-b''+b'/z-b/z^2+t^2b=A_i(l)/z.
\end{equation}

We first reduce $b$ from $S$ to $K(l)$.  Every automorphism of
$S/K(l)$ changes $b$ by a solution of $\mathcal D_t h=0$.
These homogeneous solutions are precisely $zV_i$, since
$\mathcal D_t(zf)=-z(f''+z^{-1}f'-t^2f)$.
They lie in $P$, while the difference also lies in $S$.
The product-group hypothesis gives $P\cap S=K$; no nonzero member
of $zV_i$ belongs to $K$, because the standard $\mathrm{SL}_2$
module has no invariant vector.  Thus every such difference is zero.
The fixed-field theorem, applied to the Picard--Vessiot extension
$S/K(l)$, implies $b\in K(l)$.

In fact $b\in K[l]$.  To prove this, suppose that an irreducible
monic $q(l)\in K[l]$ occurs in its denominator.  The derivation is
$D=d/dz$ with $Dl=1/z$.  If $q\mid Dq$, monicity implies
$\deg_l Dq<\deg_l q$ and hence $Dq=0$.  If $d=\deg_l q>0$,
the coefficient of $l^{d-1}$ would then give
$q_{d-1}'=-d/z$, which has no rational solution.  Therefore
$q\nmid Dq$.  If $b$ has a pole of order $h>0$ along $q$, then
$b''$ has a pole of order $h+2$ there, with nonzero leading
coefficient proportional to $h(h+1)(Dq)^2$.  The other terms on
the left of \eqref{depth:b-equation} have pole orders at most
$h+1$, whereas its right side is polynomial in $l$.  This is
impossible, so there is no such denominator.

Write $b=\sum_{j=0}^M b_j(z)l^j$, with $b_M\ne0$.
The operator $\mathcal D_t$ does not increase logarithmic degree.
If $M>\deg A_i$, its highest coefficient gives
$\mathcal D_t b_M=0$, contradicting the absence of a nonzero
rational homogeneous solution just proved.  If $M<\deg A_i$,
degrees already contradict \eqref{depth:b-equation}.  In the remaining
case the highest coefficient gives
\[
                  \mathcal D_t b_M=c/z,\qquad c\ne0.
\]
The rational obstruction in the proof of
Lemma~\ref{msprim:nonexact} excludes this equation: a rational
solution has no nonzero finite poles and is thus a Laurent
polynomial; its lowest power either gives an impossible pole of
order at least two at zero, or leaves the left side holomorphic
there.  Scaling by $c$ gives exactly the same obstruction as
\eqref{msprim:rational}.  This contradiction proves the lemma.
\end{proof}

\begin{theorem}[Functional independence of every selected depth]
\label{depth:functional}
Under the hypotheses of Lemma~\ref{depth:nonexact}, choose
nonnegative integers $R_1,\ldots,R_n$.  The
$\sum_iR_i$ functions
\[
               H_{s_i,r}(z),\qquad 1\le r\le R_i,
\]
are jointly algebraically independent over $L$.
\end{theorem}

\begin{proof}
For each active scale $t=s_i$, put
\[
 u_{t,0}=I_0(tz),\qquad
 u_{t,r}=\frac{l^r}{r!}-H_{t,r}(z)\quad(r\ge1).
\]
Then $\theta u_{t,r}=u_{t,r-1}$.  Define
\begin{equation}\label{depth:triangular-forward}
 v_{t,j}=\sum_{h=0}^j
       (-1)^h\frac{j!}{(j-h)!}l^{j-h}u_{t,h+1},
       \qquad 0\le j<R_i.
\end{equation}
Direct differentiation telescopes and gives
$v_{t,j}'=l^j I_0(tz)/z$.  Conversely,
\begin{equation}\label{depth:triangular-inverse}
 u_{t,r}=\sum_{j=0}^{r-1}
       \frac{(-1)^j l^{r-1-j}}{j!(r-1-j)!}v_{t,j}.
\end{equation}
Thus the two displayed families generate the same field over $L$.

Adjoining the $v_{t,j}$ is a Picard--Vessiot primitive extension
whose group is an algebraic subgroup of the additive group
$\mathbb G_a^N$, where $N=\sum_iR_i$.
All functions may be realized as analytic germs on a common simply
connected domain, so no constants beyond $\mathbb C$ are introduced.
If this subgroup were proper, a nonzero constant linear form in the
$v_{t,j}$ would be fixed and hence would belong to $L$.
Its derivative would contradict Lemma~\ref{depth:nonexact}.
Consequently the group is $\mathbb G_a^N$ and the $v_{t,j}$ are
independent over $L$.  The invertible triangular transformations
\eqref{depth:triangular-forward}--\eqref{depth:triangular-inverse},
followed by subtraction of $l^r/r!\in L$, prove the theorem.
\end{proof}

\subsection{Arithmetic specialization and the infinite old field}

Choose one representative of every nonzero algebraic sign class.
Define the coefficient field
\begin{equation}\label{depth:old-field}
\begin{split}
 F_{\rm old}=\Kexp\bigl(&\Ein_r(\lambda),\ I_0(t),C(t),\beta(t):\\
             &r\ge1,\ \lambda\in k^\times,
             t\in k^\times/\{\pm1\}\bigr).
\end{split}
\end{equation}
The notation for $t$ means the chosen representatives.
This field contains the old exponential-integral and Bessel
coordinates, but does not include any of the new primitive
coordinates as additional generators.

\begin{theorem}[All-scale, all-depth arithmetic extension]
\label{depth:values}
The values
\[
                 H_{t,r}(1),\qquad
 t\in k^\times/\{\pm1\},\quad r\ge1,
\]
are jointly algebraically independent over $F_{\rm old}$.
Equivalently, every finite set of old nonexponential coordinates
from Theorem~\ref{cf:joint}, together with any finite set of these
primitive coordinates, is algebraically independent over $\Kexp$.
\end{theorem}

\begin{proof}
First choose finitely many old coordinates, primitive scales and
depths, and exponential values.  Enlarge the Bessel list to contain
every selected primitive scale.  Include the finite set of
exponential functions and the required polyexponential chains,
together with $\log z$, in the solvable Picard--Vessiot field $S$
from the proof of Theorem~\ref{cf:joint}.  The Bessel product field
$P$ and $S$ have direct-product group, since a connected product
of $\mathrm{SL}_2$ groups has no nontrivial solvable quotient.
Theorem~\ref{depth:functional} therefore adds exactly
$\sum_iR_i$ independent functions to the entire old-coordinate
field, which is a subfield of $PS$.

Each $H_{t,r}$ is an E-function.  Its coefficient at degree $2n$
in exponential-series normalization is
\[
       -\binom{2n}{n}\frac{(t/2)^{2n}}{(2n)^r}.
\]
For fixed algebraic $t$ and fixed $r$, all conjugates have
exponential growth, and common denominators through degree $N$
are bounded exponentially: apart from the fixed algebraic
denominators of powers of $t/2$, it is enough to use
$\operatorname{lcm}(1,\ldots,N)^r$.
The equations
\[
 H_{t,1}'=(1-I_0(tz))/z,
 \qquad H_{t,r}'=H_{t,r-1}/z\quad(r\ge2)
\]
append the full primitive chains to the old first-order E-system
without introducing any finite singularity other than zero.
Siegel--Shidlovskii specialization at the ordinary point $z=1$
preserves the resulting functional transcendence degree
\cite{Beukers}.  As in Theorem~\ref{cf:joint}, the old derivative
coordinates are recovered from its displayed coordinates, and
the passage to $C(t),\beta(t)$ is birational because
$I_0(t)\ne0$ by Corollary~\ref{cf:nonzero}.  This proves the
finite arithmetic assertion, including relative independence
over the finite exponential field.

A relation over $F_{\rm old}$ would use only finitely many of
its generators and finitely many exponential coefficients.
Include all of them in the finite construction just given and
clear denominators.  The resulting relation contradicts that
finite independence statement.  This proves the theorem over
the entire field \eqref{depth:old-field}.
\end{proof}

\begin{corollary}[Higher ordinary-Bessel primitive values]
\label{depth:B}
For positive algebraic $\lambda$ and integers $r\ge1$, put
\[
 B_r(\lambda)=\sum_{n\ge1}
        \frac{(-1)^{n-1}\lambda^n}{n^r(n!)^2}
        =2^rH_{2i\sqrt\lambda,r}(1).
\]
All the $B_r(\lambda)$ are jointly independent over $F_{\rm old}$.
The same holds after replacing each selected value by
$B_r(\lambda)-\beta(2i\sqrt\lambda)$.
\end{corollary}
\begin{proof}
Distinct positive $\lambda$ give distinct scales modulo sign.
Apply Theorem~\ref{depth:values}; multiplication by $2^r$ and
subtraction of an old coordinate are invertible affine changes
over $F_{\rm old}$.
\end{proof}

\begin{corollary}[A one-defect family of every depth]
\label{depth:one-defect}
Put $T_{r,\lambda}=B_r(\lambda)-\gamma$.  Among the numbers
\[
\begin{gathered}
 \gamma,\ \delta,\ \kappa,\ \gamma-2\kappa,
 \nu,\ \gamma-2\nu,\\
 T_{r,\lambda}\quad(r\ge1,\ \lambda\in k\cap\mathbb R_{>0}),
\end{gathered}
\]
at most one is algebraic over $\Kexp$.
This does not establish the individual transcendence of a
prescribed $T_{r,\lambda}$.
\end{corollary}
\begin{proof}
Write $\eta=\Ein(1)$, $\beta_K=\beta(2)$ and
$\beta_J=\beta(2i)$.  Each of the six old expressions is a
nonzero $\Kexp$-multiple of $\gamma$ plus a linear combination
of $\eta,\beta_K,\beta_J$.  After normalizing the coefficient
of $\gamma$ to one, their coefficient vectors are distinct.
If two values were algebraic over $\Kexp$, eliminating $\gamma$
would make a nonzero linear combination of these independent
old coordinates algebraic over $\Kexp$, which is impossible.
For any pair involving a new $T_{r,\lambda}$, the same elimination
has a nonzero coefficient of at least one primitive coordinate.
Corollary~\ref{depth:B} excludes algebraicity also in this case.
\end{proof}

\begin{corollary}[Primitive towers under the original rank-three alternative]
\label{depth:conditional}
If $\trdeg_kH_4=3$, then $H_4F_{\rm old}/F_{\rm old}$ is
algebraic.  Over $H_4F_{\rm old}$ the entire family
$\{B_r(\lambda)\}$ is jointly independent, and so is the entire
family $\{T_{r,\lambda}\}$.  More generally, one may choose
either $B_r(\lambda)$ or $T_{r,\lambda}$ at each index and retain
joint independence.

If, instead, any one $T_{r,\lambda}$ is algebraic over
$F_{\rm old}$, then $\gamma$ is transcendental over
$F_{\rm old}$ and the six values
$\gamma,\delta,\kappa,C,\nu,C_J$ are algebraically independent
over $\Kexp$.
\end{corollary}
\begin{proof}
The field $H_4(e)$ contains
$k(e,\eta,\beta_K,C)$, whose four generators are independent
over $k$ by Theorem~\ref{cf:joint} and Lindemann--Weierstrass.
Under the rank-three hypothesis, $H_4(e)$ has transcendence
degree at most four, so it is algebraic over that subfield.
In particular $\gamma$ is algebraic over $F_{\rm old}$.
Moreover $H_4F_{\rm old}=F_{\rm old}(\gamma)$, because
$\delta=e(\eta-\gamma)$ and $\kappa=\beta_K-\gamma$.
Algebraic extension of a coefficient field preserves algebraic
independence; Corollary~\ref{depth:B} and translation by
$\gamma$ now prove the first assertions.  Both versions of a
single coordinate are not to be included simultaneously.

If $T_{r,\lambda}$ is algebraic over $F_{\rm old}$ while
$\gamma$ were also algebraic there, their sum $B_r(\lambda)$
would be algebraic, contrary to Corollary~\ref{depth:B}.
Thus $\gamma$ is transcendental over $F_{\rm old}$.
Adjoin it to the five independent old coordinates
$\eta,\beta_K,C,\beta_J,C_J$ over $\Kexp$ and use
\[
 \delta=e(\eta-\gamma),\qquad
 \kappa=\beta_K-\gamma,\qquad
 \nu=\beta_J-\gamma.
\]
This is an invertible affine change over $\Kexp$ and proves
the six-value assertion.  No algebraicity premise for such a
witness is established here.
\end{proof}

\part{Arithmetic consequences and reconstruction}
% Integration part 02: arithmetic consequences and relative independence.
% Coverage: Extensions §§2--3 and algebraic mirror-limit corollary;
% Joint Closures §§3--5 (excluding differential/E-ring/rank-classification material);
% AUV Addendum §2; Primitive Multiscale primitive arithmetic and root criteria;
% Density Monodromy independent-block worlds, conditional six, mirror algebraicity;
% new full-complex-anchor conjugation enhancement.
% The master paper supplies the central definitions, theorem environments, and
% the all-order/multiscale-primitive joint theorem labelled cf:joint.



\paragraph{Arithmetic notation.}
In this part $k=\overline{\mathbb Q}$, $H=H_4=k(\gamma,\delta,\kappa,C)$,
$\eta=\Ein(1)=\gamma+\delta/e$, and
$\mu(t)=\beta(t)-\gamma$ for positive algebraic $t$.
Thus $\mu(t)=\log(t/2)+K_0(t)/I_0(t)$ and $\mu(2)=\kappa$.
All appearances of $\Ein_r$ refer to the entire
polyexponential tower, not the classical exponential-integral tail.

% SOURCE COVERAGE: Extensions, uniform finite/infinite relation ideals and affine triples.
\section{Uniform codimension one and a global relation ideal}
\begin{theorem}\label{thm:many}
For distinct positive algebraic $t_1,\ldots,t_k$, put
\[
H_T=\Qa\bigl(\gamma,\delta,\mu(t_1),C(t_1),\ldots,
                       \mu(t_k),C(t_k)\bigr).
\]
Then
\begin{equation}\label{eq:many}
\trdeg_{\Qa}H_T\ge2k+1.
\end{equation}
Thus the deficit from algebraic independence of its $2k+2$ displayed
generators is at most one, independently of $k$.
\end{theorem}
\begin{proof}
The rational change of coordinates
$(I_0(t_j),I_1(t_j))\mapsto(I_0(t_j),C(t_j))$ is invertible because
$I_0(t_j)>0$. Theorem~\ref{cf:joint} therefore makes
$e,\eta,I_0(t_j),C(t_j),\beta(t_j)$ algebraically independent.
Discard the $k$ values $I_0(t_j)$. The remaining $2k+2$ independent
numbers belong to $H_T(e)$, by the connection identity $\beta(t)=\gamma+\mu(t)$ and
$\eta=\gamma+\delta/e$. Hence
$\trdeg H_T+1\ge\trdeg H_T(e)\ge2k+2$.
\end{proof}

\begin{corollary}\label{cor:finiteideal}
The ideal of polynomial relations over $\Qa$ among the displayed
generators of $H_T$ is either zero or generated by one irreducible
polynomial. In particular, two relatively prime nonzero polynomials
cannot both vanish at that tuple.
\end{corollary}
\begin{proof}
The evaluation kernel in the polynomial ring in $2k+2$ variables over
$\Qa$ is prime. Its height is the number of variables minus the
transcendence degree of the image's fraction field, and is at most one
by Theorem~\ref{thm:many}. A nonzero height-one prime in a polynomial
ring over a field, a unique factorization domain, is principal and is
generated by an irreducible polynomial.
\end{proof}

\begin{theorem}[The entire positive algebraic family]\label{thm:infinite}
Consider the indexed family
\begin{equation}\label{eq:infinite}
\gamma,\ \delta,\quad \mu(t),\ C(t)
\qquad(t\in\Qa\cap\mathbb R_{>0}).
\end{equation}
Either it is algebraically independent, or every polynomial relation
among its members is a multiple of one fixed irreducible polynomial $P$.
Deleting any variable that genuinely occurs in $P$ leaves an algebraically
independent family. Moreover one can choose the deleted member among
$\gamma,\delta,$ and the $\mu(t)$, retaining every $C(t)$.
\end{theorem}
\begin{proof}
Every polynomial involves only finitely many variables. If there is a
nonzero relation, factor it in a finite polynomial ring and choose an
irreducible factor $P$ that vanishes at the corresponding tuple.
Irreducibility is preserved on adjoining extra indeterminates. Given
another relation $Q$, put all arguments occurring in $P$ or $Q$ in a
finite set $T$, and include both coordinates for each argument.
Corollary~\ref{cor:finiteideal} applies in this larger ring. Its prime
kernel has height one and contains the irreducible element $P$, so
it equals $(P)$. Therefore $P$ divides $Q$.

If $X$ occurs in $P$ with positive degree, no nonzero polynomial
independent of $X$ can be a multiple of $P$, since degrees in $X$ add
under multiplication in a polynomial ring over a domain. Thus deleting
$X$ leaves no relations. Finally, all the $C(t)$ form an algebraically
independent family by Theorem~\ref{cf:joint} and the same coordinate change
as above. Consequently $P$ cannot involve only $C$-coordinates.
\end{proof}

\begin{remark}
The exceptional member in Theorem~\ref{thm:infinite} is not identified.
In particular, the theorem is consistent with $\gamma$ being algebraic:
in that case it is $\gamma$ that must be deleted. ``One relation'' means
one generator of the relation ideal; if $P=0$ is a relation, its infinitely
many polynomial multiples are also relations.
\end{remark}

\begin{corollary}[If the original four are dependent]\label{cor:relative}
Let $H=\Qa(\gamma,\delta,\kappa,C)$. If $\trdeg_{\Qa}H=3$, then:
\begin{enumerate}
\item for every finite set of distinct positive algebraic $t_j\ne2$,
all the numbers $\mu(t_j),C(t_j)$ are algebraically independent over $H$;
\item $e$ and $I_0(2)$ are algebraically independent over $H$.
\end{enumerate}
\end{corollary}
\begin{proof}
Adjoining $k$ new pairs to the original pair gives transcendence degree
at least $2(k+1)+1=2k+3$ by Theorem~\ref{thm:many}; this is the largest
possible increase $2k$ over a field of transcendence degree three.
For the second assertion, Theorem~\ref{cf:joint} shows that
$H(e,I_0(2))$ contains five independent values. Adjoining two elements
to $H$ can increase its transcendence degree by at most two, so the
relative increase is exactly two.
\end{proof}

\section{An affine family of independent triples}
\begin{theorem}\label{thm:affine}
For all but at most one $\lambda\in\Qa$, the three numbers
\begin{equation}\label{eq:triple}
C,\qquad \gamma+\kappa,\qquad \gamma+\lambda\delta
\end{equation}
are algebraically independent over $\Qa$.
\end{theorem}
\begin{proof}
Set $\mathcal F=\Qa(C,\gamma+\kappa)$. Theorem~\ref{cf:joint} gives
$\trdeg_{\Qa}\mathcal F=2$, while
$\mathcal F(\gamma,\delta)=H$ has transcendence degree at least three by
the bound $\operatorname{trdeg}_kH_4\ge3$. An exception in \eqref{eq:triple} means that
$\gamma+\lambda\delta$ is algebraic over $\mathcal F$. If distinct
$\lambda,\lambda'$ were exceptions, their difference would make
$\delta$ algebraic over $\mathcal F$, and then $\gamma$ would also be algebraic
over $\mathcal F$. This contradicts $\trdeg H\ge3$.
\end{proof}

For example, at least two of the following three triples are independent:
\[
(C,\gamma+\kappa,\gamma),\quad
(C,\gamma+\kappa,\gamma+\delta),\quad
(C,\gamma+\kappa,\gamma+2\delta).
\]
The theorem does not determine whether $\lambda=0$ is exceptional, so it
does not prove the transcendence of $\gamma$.

\begin{proposition}[A sufficient new input for full independence]\label{prop:missing}
An algebraic dependence of $e,I_0(2)$ over
$H=\Qa(\gamma,\delta,\kappa,C)$ would imply $\trdeg_{\Qa}H=4$.
Alternatively, transcendence of $\gamma$ over
$\Qa(e,\eta,C,\gamma+\kappa)$ would imply the same conclusion.
\end{proposition}
\begin{proof}
The first assertion is the contrapositive of
Corollary~\ref{cor:relative}(2), since $3\le\trdeg H\le4$.
For the second, Theorem~\ref{cf:joint} makes the four generators of
$\mathcal F_0=\Qa(e,\eta,C,\gamma+\kappa)$ independent. If $\gamma$ is
transcendental over $\mathcal F_0$, then
$\mathcal F_0(\gamma)=H(e)$ has transcendence degree five, forcing
$\trdeg H=4$.
\end{proof}
Neither extra input in Proposition~\ref{prop:missing} is proved here.
They specify what would suffice, not an available route around the
remaining arithmetic problem.



% SOURCE COVERAGE: Joint Closures, upper-half-plane relative family; finite reconstruction and mixed Hankel values.
\section{Upper-half-plane values independent over the target field}
\begin{lemma}[One real adjunction]
\label{newupper:real}
Let $\mathcal F\subset\mathbb C$ be stable under complex conjugation.
Suppose $u_1,\ldots,u_m,\overline{u_1},\ldots,\overline{u_m}$
are algebraically independent over $\mathcal F$. For every real number $c$,
the $u_1,\ldots,u_m$ are algebraically independent over $\mathcal F(c)$.
\end{lemma}
\begin{proof}
Otherwise their transcendence degree over $\mathcal F(c)$ is at most $m-1$.
Conjugation preserves $\mathcal F(c)$, so the conjugate tuple has the same
upper bound. The union then has transcendence degree at most
$2m-2$ over $\mathcal F(c)$ and at most $2m-1$ over $\mathcal F$, contradicting the
assumed independence of its $2m$ members.
\end{proof}

\begin{theorem}[An unconditional relative family]
\label{newupper:main}
The family
\[
 \{\Ein_r(\lambda):r\ge1,\
       \lambda\in\Qa,\ \operatorname{Im}\lambda>0\}
\]
is jointly algebraically independent over
$H=\Qa(\gamma,\delta,\kappa,C)$.
\end{theorem}
\begin{proof}
Let $\mathcal F$ be the field generated over $\Kexp$ by all
$\Ein_r(x)$ with nonzero real algebraic $x$ and all
$I_0(t),C(t),\beta(t)$ with positive algebraic $t$. It is stable under
complex conjugation. For any finite subfamily in the upper half-plane,
the conjugate arguments lie in the lower half-plane and do not meet the upper-half-plane or real arguments. The indexed
pairs $(r,\lambda)$ and their conjugate pairs are distinct. Theorem~\ref{cf:joint} shows that the values and
their conjugates are independent over $\mathcal F$. By
Lemma~\ref{newupper:real}, the upper-half values are independent
over $\mathcal F(\gamma)$.

Now $e,\eta=\Ein_1(1),C,\beta(2)\in \mathcal F$ and
\[
 \delta=e(\eta-\gamma),\qquad \kappa=\beta(2)-\gamma.
\]
Therefore $H\subset \mathcal F(\gamma)$, which gives the conclusion.
The argument is finite at every stage; it makes no assumption
about the arithmetic nature of $\gamma$ or about
$\operatorname{trdeg}_{\Qa}H$ beyond the established setup.
\end{proof}

\begin{corollary}[A concrete infinite sequence]
\label{newupper:i}
All the numbers $\Ein_r(i)$, $r=1,2,\ldots$, are
jointly algebraically independent over $H$. Consequently, for every $N$,
\[
 \trdeg_{\Qa}H(\Ein_1(i),\ldots,\Ein_N(i))=\trdeg_{\Qa}H+N
 \in\{N+3,N+4\}.
\] In particular
$\operatorname{Ci}(1)-i\operatorname{Si}(1)$ is transcendental over $H$.
\end{corollary}
\begin{proof}
Take $\lambda=i$. With the standard real sine and cosine integrals \cite[\S6.6]{DLMF},
\[
 \Ein_1(i)
 =\int_0^1\frac{1-\cos u}{u}\,du
       +i\int_0^1\frac{\sin u}{u}\,du
 =\gamma-\operatorname{Ci}(1)+i\operatorname{Si}(1).
\]
Since $\gamma\in H$, translation and sign change give the last assertion.
\end{proof}

\begin{remark}
The last corollary does not assert separate algebraic independence
of $\operatorname{Ci}(1)$ and $\operatorname{Si}(1)$ over $H$.
Likewise the proof does not allow an arbitrary additional real
normalization such as $\sqrt\pi$: its one-real-adjunction allowance
has already been used for $\gamma$.
\end{remark}
\section{Conjugation transfer to Euler--Gompertz integrals}

Write $k=\overline{\mathbb Q}$, $\eta=\Ein(1)=\gamma+\delta/e$,
$\kappa=K_0(2)/I_0(2)$, and $C=I_1(2)/I_0(2)$. For positive algebraic $t$, put
\[
 C(t)=\frac{I_1(t)}{I_0(t)},\qquad
 \beta(t)=\gamma+\log(t/2)+\frac{K_0(t)}{I_0(t)}.
\]
For any finite collection of distinct positive algebraic scales,
Theorem~\ref{cf:joint} gives the joint independence needed below.
Set
\[
 F=k\bigl(e,I_0(t_j),C(t_j),\beta(t_j):1\le j\le s\bigr).
\]
This field is stable under complex conjugation: its displayed generators
are real, and conjugation preserves $k$ setwise. It does not fix $k$
pointwise, which is not needed. The input makes $\eta$ transcendental over $F$.

\begin{lemma}[Finite reconstruction and conjugation]\label{transfer:lemma}
Let $u,v$ be real numbers in $k(\gamma,\delta)$ such that
$k(\gamma,\delta)/k(u,v)$ is algebraic. For every
$\lambda\in k\setminus\mathbb R$, the number $u+\lambda v$ is transcendental
over $F$. Consequently
\[
 e,\ u+\lambda v,\ I_0(t_j),C(t_j),\beta(t_j)\quad(1\le j\le s)
\]
are algebraically independent over $k$.
\end{lemma}
\begin{proof}
Suppose that $z=u+\lambda v$ is algebraic over $F$. Conjugating an algebraic
equation shows that $\bar z=u+\bar\lambda v$ is also algebraic over $F$,
since $\overline F=F$. The invertible linear system gives
\[
 v=\frac{z-\bar z}{\lambda-\bar\lambda},\qquad
 u=\frac{\lambda\bar z-\bar\lambda z}{\lambda-\bar\lambda}.
\]
Thus $u,v$ are algebraic over $F$. Finite reconstruction then makes
$\gamma,\delta$ algebraic over $F$, and $e\in F$ makes
$\eta=\gamma+\delta/e$ algebraic over $F$. This contradicts
Theorem~\ref{cf:joint}. The other $3s+1$ displayed numbers are already
algebraically independent, so adjoining a number transcendental over their
field gives the final assertion.
\end{proof}

\begin{corollary}[Every nonreal algebraic slope works]\label{transfer:slope}
For every $\lambda\in\overline{\mathbb Q}\setminus\mathbb R$, the numbers
\[
 e,\ \gamma+\lambda\delta,\ I_0(t_j),C(t_j),\beta(t_j)
 \quad(1\le j\le s)
\]
are algebraically independent. In particular,
\[
 e,\ \gamma+i\delta,\ I_0(2),\ C,\ \gamma+\kappa
\]
are algebraically independent.
\end{corollary}
\begin{proof}
Apply Lemma~\ref{transfer:lemma} with $(u,v)=(\gamma,\delta)$.
\end{proof}

This strengthens an ``at most one exceptional slope'' formulation: any
possible exceptional algebraic slope is real. It does not prove that either
real component $\gamma$ or $\delta$ is individually transcendental.

\subsection{Two positive mixed Hankel integrals}

The mixed-Hankel construction of \cite{Periods} defines
\[
 G_m=\int_0^\infty x^m e^{-x}(x-\log x)\,dx
     =m!(\gamma+m+1-H_m),\qquad
 \mathcal I_{r+1}=\int_0^\infty\frac{e^{-x}}{(1+x)^{r+1}}\,dx.
\]
Here $H_0=0$, $\mathcal I_1=\delta$, $\mathcal I_2=1-\delta$, and $\mathcal I_3=\delta/2$.
Define the first two mixed determinants by
\begin{align*}
 \mathcal H_1&=G_0\mathcal I_1=(\gamma+1)\delta,\\
 \mathcal H_2&=\det\begin{pmatrix}G_1\mathcal I_1&G_2\mathcal I_2\\G_2\mathcal I_2&G_3\mathcal I_3\end{pmatrix}.
\end{align*}
Both are strictly positive. For $\mathcal H_2$, the entries are moments $M_r$,
$r=0,1,2$, of the product variable $y=xt$ under the positive measure
\[
 x e^{-x}(x-\log x)\,\frac{e^{1-1/t}}t\,dx\,dt
 \quad\text{on }(0,\infty)\times(0,1).
\]
Indeed, $\mathcal I_{r+1}=\int_0^1 t^r e^{1-1/t}\,dt/t$. Since $x-\log x\ge1$ and
$y$ is not almost everywhere constant, strict Cauchy--Schwarz gives
$\mathcal H_2=M_0M_2-M_1^2>0$.

\begin{proposition}[Explicit degree-four reconstruction]\label{transfer:quartic}
One has
\begin{align*}
 2\mathcal H_2={}&-18+36\delta-5\delta^2-24\gamma+48\gamma\delta
       -5\gamma\delta^2-8\gamma^2+16\gamma^2\delta
       -2\gamma^2\delta^2.
\end{align*}
The number $\delta$ satisfies
\begin{equation}\label{transfer:quartic-eq}
 2T^4+(\mathcal H_1-4)T^3+(2\mathcal H_1^2-16\mathcal H_1+2\mathcal H_2+2)T^2
       +(-16\mathcal H_1^2+8\mathcal H_1)T+8\mathcal H_1^2=0,
\end{equation}
and $\gamma=\mathcal H_1/\delta-1$. Hence
\[
 [k(\gamma,\delta):k(\mathcal H_1,\mathcal H_2)]\le4.
\]
\end{proposition}
\begin{proof}
Substitute $G_1=\gamma+1$, $G_2=2\gamma+3$, and
$G_3=6\gamma+13$ into the determinant, using the three stated $I$-values,
and expand. Since $\delta>0$, the relation $\mathcal H_1=(\gamma+1)\delta$ gives
$\gamma=\mathcal H_1/\delta-1$. Substitute it in the expanded expression for $2\mathcal H_2$
and clear $\delta^2$ to obtain \eqref{transfer:quartic-eq}. Its leading
coefficient is $2$, so it is a nonzero polynomial, proving the degree bound.
\end{proof}

\begin{theorem}[Joint independence of a complex positive-integral value]
\label{transfer:hankel}
Put $Z=\mathcal H_1+i\mathcal H_2$. For every finite collection of distinct positive algebraic
$t_1,\ldots,t_s$, the numbers
\[
 e,\ Z,\ I_0(t_j),C(t_j),\beta(t_j)\quad(1\le j\le s)
\]
are algebraically independent over $\overline{\mathbb Q}$. In particular,
$C,\gamma+\kappa,Z$ are algebraically independent.
\end{theorem}
\begin{proof}
Use Proposition~\ref{transfer:quartic} and Lemma~\ref{transfer:lemma} with
$(u,v,\lambda)=(\mathcal H_1,\mathcal H_2,i)$.
\end{proof}

The transcendence of $Z$ alone was already proved in the source mixed-Hankel
construction. The new conclusion is its joint independence with the Bessel
coordinates and the explicit Euler--Bessel combination. Positivity alone
does not imply this: finite arithmetic reconstruction and conjugation are
both used.

There is also an exact statement involving the unshifted quotient $\kappa$.
The established bound
$\operatorname{trdeg}_k k(\gamma,\delta,\kappa,C)\ge3$ implies
\begin{equation}\label{transfer:unshifted}
 Z\text{ is transcendental over }k(\kappa,C),\qquad
 \operatorname{trdeg}_k k(Z,\kappa,C)
 =1+\operatorname{trdeg}_k k(\kappa,C).
\end{equation}
Indeed, if $Z$ were algebraic over $k(\kappa,C)$, conjugation and
Proposition~\ref{transfer:quartic} would make $\gamma,\delta$ algebraic over
that field, forcing the transcendence degree of all four constants to be at
most two. No independence of $\kappa,C$ is assumed or concluded here.



% SOURCE COVERAGE: Joint Closures, enlarged coefficient field and conditional entire polyexponential family.
\section{Stronger base fields and a criterion for full independence}

Set
\[
 \eta=\Ein_1(1)=\gamma+\delta/e,\qquad
 \kappa=K_0(2)/I_0(2),\qquad C=I_1(2)/I_0(2),\qquad
 H=\overline{\mathbb Q}(\gamma,\delta,\kappa,C).
\]
Let $\mathcal U$ denote the countable indexed family
\[
 \{\Ein_r(\lambda):\lambda\in\overline{\mathbb Q}^{\times},\ r\ge1,
                         (\lambda,r)\ne(1,1)\},
\]
and define the explicit coefficient field
\[
 \mathcal A=K_{\exp}\bigl(I_0(2),\mathcal U\bigr).
\]
\begin{corollary}[The three-dimensional bound survives base extension]
\label{pj:largebase}
One has
\[
 \operatorname{trdeg}_{\mathcal A}
 \mathcal A(\gamma,\delta,\kappa,C)\ge3.
\]
\end{corollary}
\begin{proof}
Theorem~\ref{cf:joint} makes $\eta$, $\beta(2)=\gamma+\kappa$, and
$C$ algebraically independent over $\mathcal A$. All three belong to
the field in the assertion, since $e\in\mathcal A$ and
$\eta=\gamma+\delta/e$.
\end{proof}
This is stronger than merely adjoining independent auxiliary numbers
to enlarge the total transcendence degree: the auxiliary values may
all occur in coefficients, and three independent dimensions of the
original four constants still remain. It does not establish a fourth
dimension.

\begin{theorem}[Rank three forces the entire polyexponential family to stay free]
\label{pj:rankthree}
Suppose $\operatorname{trdeg}_{\overline{\mathbb Q}}H=3$. Then the full
family
\[
 \{\Ein_r(\lambda):\lambda\in\overline{\mathbb Q}^{\times},\ r\ge1\},
\]
\emph{including $\eta=\Ein_1(1)$}, is algebraically independent over $H$.
More precisely, $e$ and $I_0(2)$ are algebraically independent over $H$,
and $\mathcal U$ is algebraically independent over $H(e,I_0(2))$.
\end{theorem}
\begin{proof}
For distinct $U_1,\ldots,U_N\in\mathcal U$, Theorem~\ref{cf:joint}
gives independence of the $N+5$ numbers
\[
 e,\eta,I_0(2),C,\gamma+\kappa,U_1,\ldots,U_N.
\]
They belong to $H(e,I_0(2),U_1,\ldots,U_N)$. With $N=0$ and
$\operatorname{trdeg}H=3$, the two newly adjoined generators $e,I_0(2)$
must contribute exactly two dimensions. For general $N$ the same
count shows that every $U_i$ contributes one further dimension,
proving the more precise assertions.

Finally $\delta\ne0$, so
\[
 e=\frac{\delta}{\eta-\gamma},\qquad H(\eta)=H(e).
\]
As $e$ is transcendental over $H$, so is $\eta$. Independence of every
finite subfamily of $\mathcal U$ over $H(e,I_0(2))$ implies its
independence over $H(\eta)$. Adjoining the already transcendental
$\eta$ proves independence of the entire family over $H$.
\end{proof}

\begin{corollary}[A relation sufficient for full independence]
\label{pj:relationcriterion}
If finitely many distinct polyexponential values $\Ein_r(\lambda)$ satisfy
a nonzero polynomial relation with coefficients in $H$, then
\[
 \operatorname{trdeg}_{\overline{\mathbb Q}}H=4;
\]
that is, $\gamma,\delta,\kappa,C$ are algebraically independent.
The same conclusion follows from a nonzero polynomial relation among
members of $\mathcal U$ with coefficients in $H(e,I_0(2))$.
\end{corollary}
\begin{proof}
The known bound gives $3\le\operatorname{trdeg}H\le4$.
Either proposed relation excludes the value three by
Theorem~\ref{pj:rankthree}.
\end{proof}
No such relation is proved here. This is a sufficient condition,
not a claim that the original four numbers are now independent.
Known identities expressing an $\Ein_r(1)$ as a gamma-jet coefficient
minus an incomplete tail do not meet this condition unless the
additional tail and gamma-jet data are themselves controlled over
$H$. Rapid convergence of auxiliary values to $\gamma$ is likewise
not a polynomial relation over $H$.


% SOURCE COVERAGE: AUV Addendum, exact relation ideal and Jacobian classification of all rational shifts.
\section{Exact relations and rationally shifted independent families}
\label{sec:auv-rational}

Write $k=\Qa$, $I=I_0(2)$, and
\[
 \eta=\gamma+\delta/e,\qquad \beta=\gamma+\kappa,\qquad
 A=\eta-\beta,\qquad U=3\eta-2\beta,\qquad V=\beta.
\]
Equivalently, if $\ell=\gamma-2\kappa$, then
$U=\ell+3\delta/e$ and $V=\ell+3\kappa$.
Theorem~\ref{cf:joint} makes $e,I,C,\eta,\beta$ algebraically independent over $k$.

\begin{proposition}[The complete relation ideal]\label{auv:ideal}
Every two of $A,U,V$ are algebraically independent over $k(e,I,C)$.
The evaluation homomorphism
\[
 k[X_e,X_I,X_C,X_A,X_U,X_V]\longrightarrow
 k[e,I,C,A,U,V]
\]
has kernel precisely $(X_U-X_V-3X_A)$.  In particular,
\[
 \trdeg_k k(e,I,C,A,U,V)=5.
\]
\end{proposition}
\begin{proof}
Since $\eta=A+V$ and $\beta=V$, the five numbers $e,I,C,A,V$ are
algebraically independent.  The only extra coordinate is $U=3A+V$.
Division by the monic polynomial $X_U-X_V-3X_A$ leaves a remainder
in $k[X_e,X_I,X_C,X_A,X_V]$, and independence forces any vanishing
remainder to be zero.  Finally, any two of
$A,3A+V,V$ recover $A,V$ by an invertible linear transformation.
\end{proof}

The single relation among $A,U,V$ does not persist after arbitrary
common nonlinear shifts by $C$.  The following result describes the
exception completely, even when the three rational functions differ.

\begin{theorem}[Rational shifts and their exact exception]\label{auv:shifts}
Let $f,g,h\in k(T)$ be nonconstant, and put
\[
 X=C+f(A),\qquad Y=C+g(U),\qquad Z=C+h(V).
\]
Then $e,I,X,Y,Z$ are algebraically independent over $k$, except when
\begin{equation}\label{auv:affine-exception}
 \begin{split}
 f(T)&=\alpha T+a_0,\quad
 g(T)=\rho T+b_0,\quad h(T)=\chi T+c_0,\\
 &\alpha\rho\chi\ne0,\qquad
 \rho(3\chi+\alpha)=\alpha\chi.
 \end{split}
\end{equation}
In this exceptional case, $e,I,X,Z$ are algebraically independent,
and every polynomial relation among $e,I,X,Y,Z$ follows from
\begin{equation}\label{auv:exception-relation}
 Y-b_0=\frac{3\rho}{\alpha}(X-a_0)
                  +\frac{\rho}{\chi}(Z-c_0).
\end{equation}
\end{theorem}
\begin{proof}
Put $\mathcal F=k(e,I)$.  Proposition~\ref{auv:ideal} makes $C,A,V$
algebraically independent over $\mathcal F$.  All rational evaluations in the
statement are therefore defined, and it suffices to study the rational
functions
\[
 c+f(x),\qquad c+g(3x+y),\qquad c+h(y)
\]
in the independent variables $c,x,y$ over $\mathcal F$.  In characteristic
zero, their algebraic independence is equivalent to their Jacobian
determinant being nonzero.  This determinant is
\[
 D=3g'(3x+y)h'(y)-f'(x)h'(y)+f'(x)g'(3x+y).
\]
The three derivatives are nonzero rational functions.  On defining
$F=1/f'$, $G=1/g'$, and $J=1/h'$, the identity $D=0$ becomes
\begin{equation}\label{auv:functional}
 G(3x+y)=3F(x)+J(y).
\end{equation}
Differentiate with respect to $x$ and $y$.  The result is
\[
 G'(3x+y)=F'(x)=J'(y).
\]
Since $x$ and $y$ are independent variables, these rational functions
are a common constant $s\in k$.  Hence, for constants $b,d\in k$,
\[
 F(T)=sT+b,\qquad J(T)=sT+d,\qquad G(T)=sT+3b+d.
\]
If $s\ne0$, the equality $f'(T)=1/(sT+b)$ is impossible: the
derivative of a rational function has zero residue at every finite
pole, whereas the right side has residue $1/s$ at $T=-b/s$.
Consequently $s=0$.  The nonzero constant reciprocals of the
derivatives imply that $f,g,h$ are affine with nonzero slopes
$\alpha,\rho,\chi$, and \eqref{auv:functional} reduces to
\[
 \frac1\rho=\frac3\alpha+\frac1\chi.
\]
This is exactly \eqref{auv:affine-exception}; conversely that condition
makes $D=0$ and gives \eqref{auv:exception-relation}.

In the exceptional case, $c+\alpha x+a_0$ and
$c+\chi y+c_0$ are algebraically independent over $\mathcal F$, for example by
their rank-two Jacobian.  Thus $e,I,X,Z$ are independent, and division
by the monic linear relation in $Y$ proves the assertion about all
relations, exactly as in Proposition~\ref{auv:ideal}.
\end{proof}

\begin{corollary}[Every common rational shift works]\label{auv:common}
For every nonconstant $f\in k(T)$, the five numbers
\[
 e,\ I,\ C+f(A),\ C+f(U),\ C+f(V)
\]
are algebraically independent over $k$.  This includes $f(T)=T$,
$f(T)=T^m$ for every integer $m\ge1$, and $f(T)=1/T$.
\end{corollary}
\begin{proof}
The only possible exception would require $f$ to be affine and its
common nonzero slope $\alpha$ to satisfy $4\alpha^2=\alpha^2$.
This cannot occur in characteristic zero.
\end{proof}





% SOURCE COVERAGE: Primitive Multiscale, all primitive rank, combination, and witness corollaries.
\section{Arithmetic consequences of the Bessel primitive}
Write
\[
 B=\sum_{n\ge1}\frac{(-1)^{n-1}}{n(n!)^2}=2H_{2i}(1),
 \qquad \mathcal L=B-\gamma.
\]
The analytic interpretation of $\mathcal L$ as a critical
Newton--Laguerre limit is established in the approximation section.
Here only its arithmetic consequences are needed.

\begin{corollary}[The new limit joins the one-defect family]
\label{prim:rank}
With $\nu=-\pi Y_0(2)/(2J_0(2))$ and $C_J=J_1(2)/J_0(2)$,
\[
 \operatorname{trdeg}_{k}
 k(\gamma,\delta,\kappa,C,\nu,C_J,\mathcal L)\ge6.
\]
Moreover, at most one of
\[
 \gamma,\ \delta,\ \kappa,\ \gamma-2\kappa,\
 \nu,\ \gamma-2\nu,\ \mathcal L
\]
is algebraic. If $\mathcal L$ is algebraic, the six original numbers
$\gamma,\delta,\kappa,C,\nu,C_J$ are all algebraically independent;
indeed they are jointly independent with $e,I_0(2),J_0(2)$.
\end{corollary}
\begin{proof}
Theorem~\ref{msprim:values} gives independence of
\[
 e,\eta,I_0(2),C,\gamma+\kappa,
 J_0(2),C_J,\gamma+\nu,B.
\]
The seven coordinates $e,\eta,C,\gamma+\kappa,C_J,\gamma+\nu,B$
belong to the field in the first assertion after adjoining $e$;
removing $e$ lowers transcendence degree by at most one.
Similarly $e,\eta,\gamma+\kappa,\gamma+\nu,B$ belong to
$k(\gamma,\delta,\kappa,\nu,\mathcal L,e)$, so
$\operatorname{trdeg} k(\gamma,\delta,\kappa,\nu,\mathcal L)\ge4$.
Any two different forms in the displayed list have linearly
independent coefficient vectors in these five variables. If both
were algebraic, the latter transcendence degree would be at most
three, a contradiction.

Finally, if $\mathcal L=c\in k$, then $\gamma=B-c$.
The identities $\delta=e(\eta-\gamma)$,
$\kappa=(\gamma+\kappa)-\gamma$, and
$\nu=(\gamma+\nu)-\gamma$ show that the nine independent input
coordinates generate exactly
$k(e,\gamma,\delta,\kappa,C,\nu,C_J,I_0(2),J_0(2))$.
Its nine displayed generators are therefore independent.
\end{proof}

\begin{corollary}[Explicit new independent combinations]
\label{prim:combinations}
The five numbers
\[
 \mathcal L-\kappa,\quad \mathcal L-\nu,\quad
 \mathcal L-\delta/e,\quad
 \mathcal L+\operatorname{Ei}(1),\quad
 \mathcal L+\operatorname{Ci}(1)
\]
are algebraically independent. They remain jointly independent after
adjoining $e,I_0(2),J_0(2),C,C_J$.
\end{corollary}
\begin{proof}
Put $M=\Ein(-1)$ and
$R=(\Ein(i)+\Ein(-i))/2$.
Then $\operatorname{Ei}(1)=\gamma-M$ and
$\operatorname{Ci}(1)=\gamma-R$. The five displayed numbers equal
\[
 B-(\gamma+\kappa),\quad B-(\gamma+\nu),\quad
 B-\eta,\quad B-M,\quad B-R.
\]
The six coordinates $B,\gamma+\kappa,\gamma+\nu,\eta,M,R$
are independent over $k(e,I_0(2),J_0(2),C,C_J)$ by
Theorem~\ref{msprim:values} and the invertible change from
$\Ein(\pm i)$ to their real and imaginary combinations.
The five differences are a full-rank linear transformation of six
independent variables, which proves the assertion.
\end{proof}

\begin{corollary}[A criterion over the original four-constant field]
Let $H_4=k(\gamma,\delta,\kappa,C)$. If $\mathcal L$ is algebraic
 over $H_4$, then $\operatorname{trdeg}_k H_4=4$.
Moreover
\[
 C,\ C_J,\ \mathcal L-\delta/e,\ \mathcal L-\kappa,\ \mathcal L-\nu
\]
are jointly independent over $K_{\exp}$, and remain so after adjoining
$B=\gamma+\mathcal L$. If $\mathcal L\in k$, then
$\gamma,\delta,\kappa,C,\nu,C_J$ are jointly independent over
$K_{\exp}$.
\end{corollary}
\begin{proof}
The six independent numbers $e,I_0(2),\eta,\gamma+\kappa,C,B$
lie in $H_4(e,I_0(2),\mathcal L)$. Under the algebraicity assumption,
its transcendence degree is at most $\operatorname{trdeg} H_4+2$,
so $\operatorname{trdeg}H_4\ge4$.
For the second assertion, use the invertible change of coordinates
\[
 (C,C_J,\eta,\gamma+\kappa,\gamma+\nu,B)
 \longmapsto
 (C,C_J,B-\eta,B-(\gamma+\kappa),B-(\gamma+\nu),B).
\]
For the last assertion, $\gamma=B-\mathcal L$ and
$\delta=e(\eta-\gamma)$ give the same invertible rational change
as above, now over $K_{\exp}$, where $e$ is already a coefficient.
\end{proof}



% SOURCE COVERAGE: Primitive Multiscale, both inverse points and all conditional arithmetic consequences.
\section{Inverse-function points and arithmetic consequences}
Let $q_*,x_*>0$ be defined by
\[
 \Ein(q_*)=\gamma,\qquad \beta(x_*)=\gamma.
\]
The second equation is equivalent to
$K_0(x_*)/I_0(x_*)=-\log(x_*/2)$.
Both roots exist uniquely. The first defining function is strictly
increasing, starts at zero, and is unbounded. For the second,
\[
 \beta'(x)=\frac{1-I_0(x)^{-2}}{x}>0\quad(x>0),
\]
while $\beta(0)=0$ by its series and $\beta(x)\to\infty$.
Moreover $0<q_*<1$ and $0<x_*<2$, since
$\eta>\gamma$ and $\beta(2)=\gamma+\kappa>\gamma$.
Numerically,
\[
 q_*=0.67635507786536\ldots,\qquad
 x_*=1.87744825105766\ldots.
\]

\begin{theorem}[Conditional full independence and E-integrality]
\label{thm:roots}
If either $q_*$ or $x_*$ is algebraic, the nine numbers
\[
 e,I_0(2),J_0(2),\gamma,\delta,\kappa,C,\nu,C_J
\]
are algebraically independent over $k$. The six numbers
$\gamma,\delta,\kappa,C,\nu,C_J$ are then jointly independent
over $\Kexp$ as well. The two roots cannot both be algebraic.
Furthermore,
\[
 q_*\in k\ \Longrightarrow\ \gamma,\delta\in\mathbf E,
 \qquad
 x_*\in k\ \Longrightarrow\ \gamma,\delta\notin\overline{\mathbf E}.
\]
\end{theorem}
\begin{proof}
If $q_*\in k$, add $\Ein(q_*)$ to the joint input at the
Ein point $1$ and Bessel scales $2,2i$. Its point is different from
$1$. Substituting $\gamma=\Ein(q_*)$,
$\delta=e(\eta-\gamma)$, $\kappa=\beta(2)-\gamma$, and
$\nu=\beta(2i)-\gamma$ is an invertible rational change of
coordinates. This proves both independence assertions. It also
shows $\gamma,\delta\in\mathbf E$.

If $x_*\in k$, use the additional Bessel scale $x_*$, which is
neither $\pm2$ nor $\pm2i$, and retain the extra coordinate
$\beta(x_*)=\gamma$. The same substitutions prove independence.
The image of $\beta(x_*)$ under $\Phi_-$ has genuine poles, by
the positive-scale phase formula; its numerator at a
zero of the denominator does not vanish. The monic-relation pole
criterion gives $\gamma\notin\overline{\mathbf E}$. If $\delta$ were integral,
then $\gamma=\eta-e^{-1}\delta$ would be integral, since
$e^{-1}\in\mathbf E$ and integral elements form a ring. Hence
$\delta\notin\overline{\mathbf E}$ too.

If both roots were algebraic, the joint independence of
$\Ein(q_*)$ and $\beta(x_*)$ would be contradicted by their equality.
The incompatible integrality conclusions provide another proof.
\end{proof}

Let $H_4=k(\gamma,\delta,\kappa,C)$.
Theorem~\ref{thm:roots} and Corollary~\ref{msprim:rankcriterion} imply
\[
 \trdeg_k H_4=3\ \Longrightarrow\
 q_*,x_*\text{ transcendental and }\mathcal L\text{ transcendental over }H_4.
\]
Thus any polynomial relation among the original four constants
would force both roots to be transcendental. Conversely, algebraicity
of either root supplies a sufficient condition for rank four;
no algebraicity assertion about either root has been proved.
In either algebraic-root case,
\[
 \dim_k\operatorname{span}_k\{1,e,\delta,e\eta\}=4.
\]
Indeed $e\eta=e\gamma+\delta$, and a linear relation would be a
polynomial relation among the independent numbers $e,\gamma,\delta$.



% SOURCE COVERAGE: Density Monodromy, finite-block lemma, 2/3/4-world equivalences and conditional-six independence.
\section{Independent-block alternatives and the stronger conditional family}
\label{gs:section}

Put $k=\overline{\mathbb Q}$ and abbreviate
\[
 I=I_0(2),\quad I'=I_1(2),\quad K=K_0(2),\quad
 J=J_0(2),\quad J'=J_1(2),\quad Y_*=\pi Y_0(2),
\]
\[
 \kappa=K/I,\quad C=I'/I,\quad C_J=J'/J,\quad
 \eta=\Ein(1)=\gamma+\delta/e,
 \quad \beta_I=\gamma+\kappa,\quad
 \beta_J=\gamma-\frac{Y_*}{2J}.
\]
In this section, $Y_0$ denotes the classical Bessel function of the
second kind. It is not the harmonic-coefficient entire function used
in the construction of the joint $E$-system. Also set
\[
 B=\sum_{n\ge1}\frac{(-1)^{n-1}}{n(n!)^2},\qquad
 \mathcal L=B-\gamma.
\]
Theorem~\ref{msprim:values} gives the independence in
\begin{equation}\label{gs:joint-nine}
 e,\eta,I,C,\beta_I,J,C_J,\beta_J,B
 \quad\hbox{are algebraically independent over }k.
\end{equation}
The assertions below that involve $B$ or $\mathcal L$ use this primitive
extension; the two- and three-block assertions need only the
corresponding subtuples. None of these assertions presupposes the
transcendence of $\gamma$.

\subsection{A finite-block lemma}

\begin{lemma}\label{gs:block-lemma}
Let $X_1,\ldots,X_r$ be finite tuples whose concatenation is
algebraically independent over an algebraically closed field
$k\subset\mathbb C$. If $g\in\mathbb C$ is transcendental over $k$,
then $g$ is algebraic over at most one of the fields $k(X_j)$.
Consequently, at least $r-1$ of the tuples $(X_j,g)$ are algebraically
independent over $k$.
For $r\ge2$, transcendence of $g$ is equivalent to this last assertion.
\end{lemma}

\begin{proof}
Suppose that $g$ is algebraic over both $k(X_i)$ and $k(X_j)$,
where $i\ne j$. Write $d_i=|X_i|$ and $d_j=|X_j|$.
The two finitely generated fields
$F_i=k(X_i,g)$ and $F_j=k(X_j,g)$ have transcendence degrees
$d_i$ and $d_j$, respectively. Their compositum contains the
independent concatenation $(X_i,X_j)$, and their intersection
contains $k(g)$. Submodularity of transcendence degree gives
\[
 d_i+d_j
 \ge \operatorname{trdeg}_k(F_iF_j)
       +\operatorname{trdeg}_k(F_i\cap F_j)
 \ge d_i+d_j+1,
\]
a contradiction. Thus $g$ is transcendental over all but at most
one $k(X_j)$, which is equivalent to the asserted independence of
$(X_j,g)$ since $X_j$ is already independent. Conversely, when
$r\ge2$, at least one independent tuple containing $g$ implies that
$g$ is transcendental.
\end{proof}

The lemma is the finite-block version of the algebraic-tail
intersection argument in manuscript [021] \cite{Source021}, Lemma
\texttt{em:tail-intersection}. It is an elementary field-theoretic
input, rather than a new transcendence theorem.

\begin{theorem}[Two, three, and four independent worlds]
\label{gs:worlds}
The following equivalences hold.
\begin{enumerate}
\item $\gamma$ is transcendental if and only if at least one of
\[
 (e,\gamma,\delta),\qquad (I,I',K,\gamma)
\]
is algebraically independent over $k$.
\item $\gamma$ is transcendental if and only if at least two of
\[
 (e,\gamma,\delta),\qquad (I,I',K,\gamma),\qquad
 (J,J',Y_*,\gamma)
\]
are each algebraically independent over $k$.
\item With the primitive extension, $\gamma$ is transcendental if
and only if at least three of the four tuples
\begin{equation}\label{gs:four-worlds}
 (e,\gamma,\delta),\qquad (I,I',K,\gamma),\qquad
 (J,J',Y_*,\gamma),\qquad (\gamma,\mathcal L)
\end{equation}
are each algebraically independent over $k$.
\end{enumerate}
These are statements about the individual tuples, not about the
independence of their concatenation.
\end{theorem}

\begin{proof}
Use the respective first two, first three, or all four blocks
\[
 X_1=(e,\eta),\quad X_2=(I,I',\beta_I),\quad
 X_3=(J,J',\beta_J),\quad X_4=(B).
\]
Their concatenation is independent by \eqref{gs:joint-nine}, since
replacing $C$ by $I'=IC$ and $C_J$ by $J'=JC_J$ gives invertible
rational changes of coordinates. The nonvanishing of $I,J$ follows
also from their transcendence (and for $I$ from positivity).
For each block, adjoining $\gamma$ and making another invertible
rational change of coordinates gives
\[
 \delta=e(\eta-\gamma),\qquad
 K=I(\beta_I-\gamma),\qquad
 Y_*=2J(\gamma-\beta_J),\qquad
 \mathcal L=B-\gamma.
\]
The displayed tuples and $(X_j,\gamma)$ therefore generate the same
fields and have the same numbers of entries. Apply
Lemma~\ref{gs:block-lemma}. If $\gamma$ were algebraic, every tuple
in question would be dependent, proving the reverse implications.
\end{proof}

The first assertion does not imply that both of its tuples are
independent whenever $\gamma$ is transcendental. At the level of
these algebraic-independence inputs, for instance, the assignment
$g=\eta-d/e$ with $0\ne d\in k$ makes the corresponding
$\delta=d$ algebraic while leaving $g$ transcendental and independent
of the Bessel block. Such a formal example does not assign a new
value to the actual Euler constant; it shows the limitation of this
abstract implication. The alternatives in
Theorem~\ref{gs:worlds} do not prove independence of
$\gamma,\delta,\kappa,C$.

\subsection{A stronger consequence if the original four are dependent}

\begin{corollary}\label{gs:conditional-six}
Let $H_4=k(\gamma,\delta,\kappa,C)$. If the four generators of
$H_4$ are algebraically dependent, then
\begin{equation}\label{gs:six-over-H4}
 e,\ K_0(2),\ J_0(2),\ J_1(2),\ \pi Y_0(2),\ \mathcal L
 \quad\hbox{are algebraically independent over }H_4.
\end{equation}
In particular, $K_0(2)$ and $\mathcal L$ are algebraically independent
over $H_4(e)$. Even without the $J$-block, the shorter conclusion
that $e,K_0(2),\mathcal L$ are independent over $H_4$ follows.
\end{corollary}

\begin{proof}
The five independent values $e,\eta,I,C,\beta_I$ lie in
$H_4(e,I)$, so $\operatorname{trdeg}_kH_4\ge3$. Under the stated
hypothesis equality holds. By \eqref{gs:joint-nine}, the field
\[
 F=H_4(e,I,J,C_J,\beta_J,B)
\]
has transcendence degree at least nine over $k$. It is obtained from
the rank-three field $H_4$ by adjoining six elements, so its degree
is exactly nine and these six elements are jointly independent over
$H_4$. Over $H_4$ one has the invertible transformations
\[
 I\longleftrightarrow K=\kappa I,\qquad
 (J,C_J,\beta_J)\longleftrightarrow
 (J,J'=JC_J,Y_*=2J(\gamma-\beta_J)),\qquad
 B\longleftrightarrow\mathcal L=B-\gamma.
\]
Here $\kappa>0$ and $J\ne0$. This proves
\eqref{gs:six-over-H4} and its subtuple conclusions. For the shorter
proof without the $J$-block, the six independent entries
$e,\eta,I,C,\beta_I,B$ lie in $H_4(e,I,B)$; the identical rank count
proves independence of $e,I,B$ over $H_4$.
\end{proof}

Thus, for example, algebraicity of either $K_0(2)$ or
$\mathcal L$ over $H_4(e)$ would force
$\operatorname{trdeg}_kH_4=4$. This is a conditional criterion, not
a proof of that algebraicity or of rank four.



% SOURCE COVERAGE: Density Monodromy generalized mirror algebraicity; Extensions cor:limits unique implications.
\section{Arithmetic of the two mirror limits}
The approximation results identify
\[
 L_-=\ell=\gamma-2\kappa,\qquad L_+=\gamma+2\kappa.
\]
Both differ from the Newton--Laguerre primitive limit
$\mathcal L=B-\gamma$. Their algebraic consequences are as follows.

\begin{corollary}\label{gs:mirror-algebraicity}
Among $\gamma,\delta,\kappa,\ell,L_+$, at most one is algebraic.
More generally, at most one of $\delta$, $\kappa$, and all the numbers
$\gamma+c\kappa$ with $c\in k$ can be algebraic.
\end{corollary}

\begin{proof}
Two distinct forms among $\kappa$ and $\gamma+c\kappa$ would,
if both algebraic, force both $\gamma$ and $\kappa$ algebraic,
contradicting transcendence of $\beta_I=\gamma+\kappa$.
If $\delta$ and $\kappa$ were algebraic, the identity
$e(\eta-\beta_I)=\delta-e\kappa$ would contradict independence of
$e,\eta,\beta_I$. If $\delta$ and
$s=\gamma+c\kappa$ were algebraic, the case $c=1$ is already
excluded by $s=\beta_I$. For $c\ne1$, write
$\gamma=(s-c\beta_I)/(1-c)$ and use
$\eta=\gamma+\delta/e$ to obtain the same forbidden algebraic
relation among $e,\eta,\beta_I$. The five-number assertion follows
by taking $c=0,-2,2$.
\end{proof}

\begin{corollary}\label{cor:limits}
One has
\[
 \operatorname{trdeg}_k k(\delta,C,L_+,L_-)\ge3.
\]
If either $L_+$ or $L_-$ is algebraic, then
$\gamma,\delta,C$ are algebraically independent.
\end{corollary}
\begin{proof}
The identities
\[
 \gamma=(L_++L_-)/2,\qquad \kappa=(L_+-L_-)/4
\]
give $k(\delta,C,L_+,L_-)=H_4$, proving the first assertion.
If either mirror limit is algebraic, then $\kappa$ is a
$k$-affine function of $\gamma$, so $H_4=k(\gamma,\delta,C)$.
The established lower bound three forces these three generators
to be independent.
\end{proof}
The mirror pair is an invertible linear reparametrization of
$\gamma,\kappa$. It resolves an asymptotic identification without
adding a new independent coordinate to $H_4$.



% NEW CROSS-NOTE RESULT: full-complex-Bessel-anchor enlargement of upper-half-plane relative independence.
\section{Upper-half-plane values over the full Bessel-anchor field}
\label{anchor:section}

The multiscale primitive extension and the complex-scale joint theorem
can be combined with conjugation without spending any additional real
adjunction beyond $\gamma$.
For each nonzero algebraic scale $t$, put
\[
 H_t(1)=-\sum_{n\ge1}\frac{(t/2)^{2n}}{2n(n!)^2}.
\]
Let $\mathcal T$ be any choice of one representative from each sign
class in $\overline{\mathbb Q}^{\times}/\{\pm1\}$. The field below is
independent of that choice: $I_0,\beta,H_t(1)$ are even in $t$, while
$C(-t)=-C(t)$.
Define
\begin{equation}\label{anchor:field}
 \begin{split}
 F_0=K_{\exp}\bigl(&\Ein_r(a)
          :a\in\overline{\mathbb Q}\cap\mathbb R^{\times},\ r\ge1;\\
                  &I_0(t),C(t),\beta(t),H_t(1):t\in\mathcal T\bigr).
 \end{split}
\end{equation}
Theorem~\ref{msprim:values} 
gives joint independence of these coordinates
and any finite family of polyexponential values at additional distinct
nonreal arguments. All assertions about the infinite fields are
interpreted through finite subsets of their generators.

\begin{theorem}[Upper-half-plane independence with all Bessel anchors]
\label{anchor:main}
The entire family
\[
 \{\Ein_r(\lambda):
       \lambda\in\overline{\mathbb Q},\ \operatorname{Im}\lambda>0,
       \ r\ge1\}
\]
is algebraically independent over $F_0(\gamma)$.
\end{theorem}

\begin{proof}
The field $F_0$ is stable under complex conjugation. The field
$K_{\exp}$ is stable because conjugation sends $e^\alpha$ to
$e^{\overline\alpha}$, though $K_{\exp}$ is not a real field.
The real-point polyexponential values are real. For each algebraic
$t$, conjugation sends
\[
 (I_0(t),C(t),\beta(t),H_t(1))
 \quad\hbox{to}\quad
 (I_0(\overline t),C(\overline t),\beta(\overline t),H_{\overline t}(1)).
\]
The conjugate scale is represented in $\mathcal T$ up to sign, and
the even/odd identities above put all these values back in $F_0$.
Thus arbitrary complex algebraic Bessel scales, not only real and
imaginary scales, are allowed as anchors.

Choose a finite tuple $u_1,\ldots,u_m$ from the upper-half-plane
family. Its conjugates correspond to the same orders at lower-half-
plane arguments. These indexed values are all distinct from one
another and from the real-point polyexponential values generating
$F_0$. The joint theorem therefore makes
\[
       u_1,\ldots,u_m,\overline{u_1},\ldots,\overline{u_m}
\]
algebraically independent over $F_0$. If the first $m$ values were
dependent over $F_0(\gamma)$, their rank there would be at most
$m-1$. Since $\gamma$ is real and $F_0$ is conjugation-stable,
the conjugate tuple would satisfy the same bound. Their union
would then have rank at most $2m-2$ over $F_0(\gamma)$, and at most
$2m-1$ over $F_0$, since adjoining one element increases
transcendence degree by at most one. This contradicts its rank
$2m$ over $F_0$. Every finite tuple is therefore independent.
\end{proof}

\begin{corollary}[The primitive and classical Bessel anchors are allowed]
\label{anchor:explicit}
Put
\[
 H_4=\overline{\mathbb Q}(\gamma,\delta,\kappa,C),\qquad
 \nu=-\frac{\pi Y_0(2)}{2J_0(2)},\qquad C_J=\frac{J_1(2)}{J_0(2)},
\]
\[
 B=2H_{2i}(1),\qquad \mathcal L=B-\gamma.
\]
The upper-half-plane family of Theorem~\ref{anchor:main} is
algebraically independent over
\begin{equation}\label{anchor:explicit-field}
 H_4\bigl(\nu,C_J,\mathcal L,e,I_0(2),J_0(2)\bigr).
\end{equation}
This coefficient field also contains $K_0(2)$, $J_1(2)$, and
$\pi Y_0(2)$. The conclusion is unconditional with respect to
whether $H_4$ has transcendence degree three or four.
\end{corollary}

\begin{proof}
All the generators in \eqref{anchor:explicit-field} lie in
$F_0(\gamma)$. Indeed,
\[
 \delta=e\bigl(\Ein(1)-\gamma\bigr),\quad
 \kappa=\beta(2)-\gamma,\quad
 \nu=\beta(2i)-\gamma,\quad C_J=-iC(2i),
\]
and $J_0(2)=I_0(2i)$ and $\mathcal L=2H_{2i}(1)-\gamma$.
The three further values are recovered by
\[
 K_0(2)=\kappa I_0(2),\qquad
 J_1(2)=C_JJ_0(2),\qquad
 \pi Y_0(2)=-2\nu J_0(2).
\]
Apply Theorem~\ref{anchor:main} and restrict its coefficient field.
\end{proof}

The proof does not permit an arbitrary extra real constant to be
adjoined. In particular it provides no permission to adjoin
$\log a$ or $\sqrt\pi$ unless that particular constant is already
shown to belong to the defined field. The new anchors above are
permitted because their required real information is generated by
the joint $E$-coordinates and the same single real constant
$\gamma$; they are not separate real adjunctions.

\section{The exact cost of finitely many real adjunctions}
\label{rdef:section}

The conjugation argument has a general finite-defect form. The result
below needs finite transcendence degree of the coefficient extension,
not a finite degree or a specified list of generators.

\begin{theorem}[Conjugation and the half-defect bound]
\label{rdef:main}
Let $F\subseteq L\subseteq\mathbb C$ be fields stable under complex
conjugation, with $d=\operatorname{trdeg}_F L<\infty$. Let
$(u_i)_{i\ge1}$ be a countable indexed family such that, for every
finite set of indices $T$, the concatenation
\[
                       (u_i,\overline{u_i}:i\in T)
\]
is algebraically independent over $F$. Put $h=\lfloor d/2\rfloor$.
Then for every finite $T$,
\begin{equation}\label{rdef:rank}
 \operatorname{trdeg}_L L(u_i:i\in T)\ge |T|-h.
\end{equation}
At most $h$ of the $u_i$ are algebraic over $L$. One can delete
at most $h$ members of the entire family so that the rest are jointly
algebraically independent over $L$.

Moreover, there is a finite set of indices $T_0$ such that all remaining
values are jointly algebraically independent over $L(u_i:i\in T_0)$.
If $\mathfrak I$ is the ideal of all polynomial relations in
$L[X_1,X_2,\ldots]$, then
\begin{equation}\label{rdef:finite-core}
 \mathfrak I=
  \bigl(\mathfrak I\cap L[X_i:i\in T_0]\bigr)
                         L[X_1,X_2,\ldots].
\end{equation}
The finite-variable prime ideal on the right has height at most $h$.
In particular, all relations admit finitely many polynomial generators.
If $d\le3$, the global relation ideal is either zero or generated by
one irreducible polynomial involving only finitely many variables.
\end{theorem}
\begin{proof}
For a finite tuple of length $m$, write
$r=\operatorname{trdeg}_L L(u_i:i\in T)$. Conjugation stability
makes the conjugate tuple have the same rank $r$ over $L$. The
independent concatenation lies in $L(u_i,\overline{u_i}:i\in T)$,
whose rank over $F$ is at most $d+2r$. Thus
\[
                              2m\le d+2r.
\]
Since $m-r$ is an integer, this proves \eqref{rdef:rank}. Applying
it to any finite tuple of members algebraic over $L$ gives their
number at most $h$.

Enumerate the family and greedily retain $u_i$ exactly when it is
transcendental over $L$ and the previously retained values. At each
finite stage the retained tuple is a transcendence basis for the field
generated by all values processed so far. Inequality
\eqref{rdef:rank} therefore permits at most $h$ rejections in total.
All retained values are jointly independent.

If there is no rejection, take $T_0$ empty. Otherwise, take $T_0$
to be the finite initial segment ending with the last rejection.
Every rejected value is algebraic over the previously retained ones,
so $L(u_i:i\in T_0)$ is algebraic over the field generated by the
retained values in that segment. All later values are jointly
independent over the latter field, and remain so after this algebraic
extension. This proves the asserted independence over the finite core.

To obtain \eqref{rdef:finite-core}, group a polynomial relation by
monomials in variables outside $T_0$. Independence of their values over
$L(u_i:i\in T_0)$ says that every coefficient vanishes at the core
tuple. Thus every coefficient lies in the finite-core relation ideal,
which proves the displayed equality. Its height is
\[
 |T_0|-\operatorname{trdeg}_L L(u_i:i\in T_0)\le h.
\]
The polynomial ring in the finitely many core variables is noetherian,
so that ideal is finitely generated. If $h\le1$, a nonzero prime
ideal of that height in this UFD is principal and generated by an
irreducible polynomial. Its extension gives the final assertion.
\end{proof}

\begin{remark}[Sharpness as a field-theoretic statement]
The factor $1/2$ cannot be improved under these hypotheses alone.
Take real numbers $x_1,y_1,\ldots,x_h,y_h$ algebraically independent
over $k$, put $u_j=x_j+iy_j$, and adjoin all $2h$ real coordinates
to $F=k$. The values and their conjugates were independent over $F$,
but all $h$ values now lie in the enlarged field. Additional independent
real generators give the corresponding examples for odd $d$.
\end{remark}

\begin{corollary}[A finite relation core after adjoining
Euler's constant, $\pi$, and $\log2$]\label{rdef:upper}
Let $F_0$ be the field in \eqref{anchor:field}, or its enlargement
obtained by replacing the depth-one primitives $H_t(1)$ by all
$H_{t,r}(1)$ with $r\ge1$. Put
\[
                         L=F_0(\gamma,\pi,\log2).
\]
For the family
\[
 \mathcal U=\{\Ein_r(\lambda):
       \lambda\in\overline{\mathbb Q},\ \operatorname{Im}\lambda>0,
       \ r\ge1\},
\]
every $m$-element subfamily has rank at least $m-1$ over $L$.
At most one member is algebraic over $L$, and deleting at most one
member makes the entire family algebraically independent over $L$.
All its polynomial relations over $L$ are generated by a single
finite-variable irreducible polynomial, if there is any relation at all.
The same statements hold over the algebraic closure of $L$ inside
$\mathbb C$.
\end{corollary}
\begin{proof}
The conjugation stability of $F_0$ was proved in
Theorem~\ref{anchor:main}. The same proof applies to all primitive
depths, because conjugation sends $H_{t,r}(1)$ to
$H_{\overline t,r}(1)$ and these functions are even in the scale.
The joint theorem, with Theorem~\ref{depth:values} for the enlarged
anchor, makes every finite upper-half-plane tuple and its conjugates
jointly independent over $F_0$. The three adjoined constants are real,
so $L$ is conjugation-stable and
$\operatorname{trdeg}_{F_0}L\le3$. Apply
Theorem~\ref{rdef:main}. Its algebraic closure is again conjugation-stable
and has the same transcendence degree over $F_0$.
\end{proof}

For example, for every $N\ge1$,
\[
 \operatorname{trdeg}_L L\bigl(\Ein(i),\Ein(2i),\ldots,\Ein(Ni)\bigr)
       \ge N-1.
\]
The corollary neither identifies a dependent member nor asserts that
any dependence exists. In particular, it gives no arithmetic
classification of $\gamma$ or of $\delta$.

\part{Integral intersections of rational fields}
% Unified Part III: E-value integrality and rational-field classifications.
% Assembled from five source notes and the newly audited closures.
% The only deliberate omissions are proofs strictly superseded by stronger
% statements; their conclusions are retained as explicit corollaries.


\section{Continued phase embeddings and their Bessel images}
\label{mbc:section}

Put $k=\overline{\mathbb Q}$.  Let $\E$ be the ring of
$E$-function values at algebraic points and let $\Ebar$
be its integral closure inside $\mathbb C$.  Thus ``integral'' means
satisfying a monic polynomial with coefficients in $\E$;
it does not mean algebraic over its fraction field.
The phase-embedding construction from manuscript [033] \cite{Source033} is reproduced here. The arithmetic analytic input in this section is the
Fischler--Rivoal phase embedding and the properties of its component
functions \cite[\S2 and Proposition~3.1]{FR26}.  In particular, the
conclusions here do not require a new algebraic-independence input.

\subsection{Continuation of the phase embedding}

Write
\[
 \mathcal R=\mathbb C[\overline{\mathbb Q}_{\mathrm{add}}]
 =\left\{\sum_{\rho\in k}c_\rho e^{\rho X}:
                 c_\rho\in\mathbb C,\ \text{finite support}\right\}.
\]
The equality with the group algebra uses linear independence of
distinct exponential functions.  Fix $-\pi/2<\theta<\pi/2$ and use
the right-lateral convention at an exceptional direction.  If $f$
is an $E$-function, let
\[
 f(z)=\sum_\rho f_\rho(z)e^{\rho z}
\]
be its canonical component decomposition.  Fischler--Rivoal state
immediately before their Theorem~2.3 that each
$e^{\rho z}f_\rho(z)$ solves an $E$-operator whose only possible
finite singularity is zero.  Consequently every component admits
holomorphic continuation along every path in $\mathbb C^\times$.
Let $M^j$ mean continuation around a based loop at $1$ of
counterclockwise winding number $j\in\mathbb Z$.

\begin{lemma}[Continued embeddings]\label{mbc:continued}
For every $j\in\mathbb Z$, the prescription
\[
 \Phi_\theta^{[j]}(f(1))(X)
       =\sum_\rho (M^j f_\rho)(1)e^{\rho X}
\]
defines an injective $k$-algebra homomorphism
$\E\longrightarrow\mathcal R$.  It satisfies
\[
 \Phi_\theta^{[j]}(\xi)(1)=\xi,
 \qquad \Phi_\theta^{[j]}(e^a)=e^{aX}\quad(a\in k).
\]
\end{lemma}
\begin{proof}
If $f(1)=g(1)$, Beukers' zero-removal theorem, as used in the proof
of Fischler--Rivoal's Proposition~3.1, gives an $E$-function $h$
with $f-g=(z-1)h$.  Choose one sufficiently small positive
perturbation of $\theta$ for all three functions.  Uniqueness of
the canonical decompositions gives
\[
 f_\rho-g_\rho=(z-1)h_\rho.
\]
Continue this identity around the same loop.  Since the continued
$h_\rho$ is holomorphic at the endpoint $1$, evaluation there
shows that the proposed map is independent of the representative.

The component identities for sums and products are
\[
 (f+g)_\rho=f_\rho+g_\rho,
 \qquad (fg)_\rho=\sum_{\alpha+\beta=\rho}f_\alpha g_\beta.
\]
These identities hold initially by uniqueness and remain true
after continuation.  They prove additivity and multiplicativity.
Constant functions and pure exponentials give the stated
normalizations.  Finally the entire function $f$ is single valued,
so continuation of its total component sum yields
\[
 \sum_\rho(M^j f_\rho)(1)e^\rho=f(1).
\]
This is evaluation compatibility and proves injectivity.
\end{proof}

The summation direction is held fixed in this construction.  The
continued components are not asserted to be the canonical
asymptotic components in a new sector.

\begin{lemma}[Normality and the pole obstruction]\label{mbc:pole}
The ring $\mathcal R$ is integrally closed.  If
$\xi\in\operatorname{Frac}(\E)\cap\Ebar$,
the extension of every $\Phi_\theta^{[j]}$ to the fraction field
sends $\xi$ into $\mathcal R$.  In particular, an image with a
finite pole proves $\xi\notin\Ebar$.
\end{lemma}
\begin{proof}
Suppose $H\in\operatorname{Frac}(\mathcal R)$ is integral over
$\mathcal R$.  The supports of a numerator and denominator for
$H$, together with the finitely many coefficients of a monic
integral equation, generate a finitely generated subgroup
$\Gamma\subset k_{\mathrm{add}}$.  This subgroup is free abelian,
and $\mathbb C[\Gamma]$ is a Laurent polynomial ring.  We have
$H\in\operatorname{Frac}(\mathbb C[\Gamma])$, and $H$ is integral
over $\mathbb C[\Gamma]$.  Normality of a Laurent polynomial ring
therefore gives $H\in\mathbb C[\Gamma]\subset\mathcal R$.

Apply a continued embedding to a monic integral equation for
$\xi$.  Its image is integral over $\mathcal R$ and belongs to
$\operatorname{Frac}(\mathcal R)$, hence lies in $\mathcal R$.
Every element of $\mathcal R$ is entire.  Alternatively, at a
finite pole of a meromorphic image, the highest power in the
monic equation has strictly larger pole order than every other
term, which is impossible.
\end{proof}

\subsection{Explicit winding formulas}

For positive algebraic $t$ define
\[
 I=I_0(t),\quad J=I_1(t),\quad K=K_0(t),\quad L=K_1(t),
 \qquad g=\gamma+\log(t/2).
\]
All four Bessel values here are positive.  Let
\[
 Y(t)=\sum_{n\ge0}\frac{H_n(t/2)^{2n}}{(n!)^2},
 \qquad \beta(t)=\frac{Y(t)}{I_0(t)}
              =\gamma+\log(t/2)+\frac{K_0(t)}{I_0(t)},
 \qquad C(t)=\frac{I_1(t)}{I_0(t)}.
\]
The function $Y$ is the harmonic Bessel $E$-function, not the
ordinary Bessel function $Y_0$.  In particular, $\beta(t)$ and
$C(t)$ belong to $\operatorname{Frac}(\E)$.
Set
\[
 \varepsilon=\frac1{\pi i},\qquad \tau=2\pi i.
\]
At direction $\theta=-\pi/4$, the two canonical components of
$I_0(tz)$, including their exponential factors, are
\[
 y_+(z)=I_0(tz)-\varepsilon K_0(tz),\qquad
 y_-(z)=\varepsilon K_0(tz).
\]
These are the standard two-phase Bessel components used in the
original pole calculation. These follow by applying the Bessel
connection formula to the recessive solution in the adjoining sector
\cite[Sections 10.27 and 10.34]{DLMF}. The local continuation laws are
\[
 M^jK_0(tz)=K_0(tz)-j\tau I_0(tz),\qquad
 M^jK_1(tz)=K_1(tz)+j\tau I_1(tz).
\]
Define
\begin{align*}
 A_j&=e^{-t}\bigl((1+2j)I-\varepsilon K\bigr),&
 B_j&=e^t\bigl(\varepsilon K-2jI\bigr),\\
 P_j&=e^{-t}\bigl((1+2j)J+\varepsilon L\bigr),&
 Q_j&=e^t\bigl(-2jJ-\varepsilon L\bigr).
\end{align*}
The imaginary parts of $\varepsilon K$ show that $A_jB_j\ne0$
for every integer $j$.  Differentiating the two component
identities gives
\begin{align}
 \Phi_{-\pi/4}^{[j]}(I_0(t))
      &=A_je^{tX}+B_je^{-tX},\label{mbc:Iimage}\\
 \Phi_{-\pi/4}^{[j]}(I_1(t))
      &=P_je^{tX}+Q_je^{-tX}.\label{mbc:Jimage}
\end{align}
The Wronskian identity $IL+JK=1/t$ implies
\begin{equation}\label{mbc:determinant}
 A_jQ_j-B_jP_j=-\frac{\varepsilon}{t}\ne0.
\end{equation}

For completeness, the canonical decomposition of $Y(tz)$ is
\[
 Y(tz)=(g+\log z)y_+(z)
       +\bigl(1+\varepsilon(g+\log z)\bigr)K_0(tz).
\]
Continue both brackets and $K_0$ around the same loop, using
$M^j\log z=\log z+j\tau$.  At $z=1$ this gives
\begin{equation}\label{mbc:Yimage}
 \Phi_{-\pi/4}^{[j]}(Y(t))
  =(g+j\tau)A_je^{tX}
   +(g+j\tau+\pi i)B_je^{-tX}.
\end{equation}
These formulas also check evaluation compatibility at $X=1$.

% Standalone proof fragment.  Requires amsmath, amssymb, amsthm.
% Uses only the winding embeddings and classical Bessel connection formulas.


\section{The full rational two-variable Bessel field}
\label{brp:section}

The winding argument closes a larger field than the two one-variable
fields separately.  Throughout this section, $k=\overline{\mathbb Q}$,
$t>0$ is algebraic, and
\[
 I=I_0(t),\quad J=I_1(t),\quad K=K_0(t),\quad L=K_1(t),\qquad
 \beta=\gamma+\log(t/2)+K/I,\quad C=J/I.
\]
The independence needed below will follow directly from the winding
embeddings; no additional joint-specialization input is required in
this section.

\begin{theorem}[Full rational Bessel-field exclusion]
\label{brp:main}
Let $\Ebar$ denote the elements of $\mathbb C$ integral
over the ring $\E$ of $E$-function values.  Then
\[
 \boxed{\ \Ebar\cap k(\beta,C)=k.\ }
\]
Equivalently, every nonconstant rational function in $\beta,C$ with
algebraic coefficients lies outside $\Ebar$.
\end{theorem}

\begin{proof}
Put $g=\gamma+\log(t/2)$ and $\tau=2\pi i$.  For each integer $j$,
write the continued phase images using
\begin{align*}
 A_j&=e^{-t}\bigl((1+2j)I+iK/\pi\bigr),&
 B_j&=e^t\bigl(-2jI-iK/\pi\bigr),\\
 P_j&=e^{-t}\bigl((1+2j)J-iL/\pi\bigr),&
 Q_j&=e^t\bigl(-2jJ+iL/\pi\bigr).
\end{align*}
All $A_j,B_j$ are nonzero, since $I,K$ are positive.  The Bessel
Wronskian $IL+JK=1/t$ gives
\[
 A_jQ_j-B_jP_j=\frac{i}{\pi t}.
\]
With $q=e^{2tX}$, the two images are
\begin{equation}\label{brp:images}
 x_j(X)=g+j\tau+\frac{i\pi B_j}{A_jq+B_j},\qquad
 y_j(X)=\frac{P_jq+Q_j}{A_jq+B_j}.
\end{equation}
They satisfy
\begin{equation}\label{brp:line}
 y_j=C+r_j(x_j-\beta),\qquad
 r_j=\frac{1}{\pi^2tA_jB_j}.
\end{equation}
Indeed the slope follows by eliminating $q$, and the intercept follows
from evaluation of both continued phase images at $X=1$.

These slopes are nonzero and take infinitely many distinct values.
More precisely,
\[
 A_jB_j=-2j(2j+1)I^2-\frac{i(4j+1)IK}{\pi}
              +\frac{K^2}{\pi^2},\qquad
 r_j\sim-\frac{1}{4\pi^2tI^2j^2}
 \quad(j\longrightarrow+\infty).
\]
For every fixed nonzero $r$, the equation $r_j=r$ has at most two
integer solutions, since it is a nonconstant quadratic equation in $j$.

This also proves that $\beta,C$ are algebraically independent over
$k$.  Indeed, if a nonzero $H\in k[U,V]$ satisfied $H(\beta,C)=0$,
applying each continued phase embedding would make $H$ vanish
identically on the line \eqref{brp:line}.  There are infinitely many
distinct such lines, whereas a nonzero polynomial has only finitely
many distinct linear factors.  This is impossible.  In particular,
both $\beta$ and $C$ are transcendental.

As $X$ ranges over $\mathbb C$, $q$ ranges over $\mathbb C^\times$.
Consequently \eqref{brp:images} covers the projective line
\eqref{brp:line} except for the two finite points
\begin{align*}
 U_j&=\bigl(g+j\tau,\ P_j/A_j\bigr),\\
 V_j&=\bigl(g+j\tau+i\pi,\ Q_j/B_j\bigr).
\end{align*}
In particular, its point at infinity is attained at
$q=-B_j/A_j\ne0$.  The omitted points satisfy
\begin{equation}\label{brp:endpoint-asymptotics}
 (U_j)_1=j\tau+O(1),\quad (V_j)_1=j\tau+O(1),\qquad
 (U_j)_2=C+O(j^{-1}),\quad (V_j)_2=C+O(j^{-1}).
\end{equation}

Take a nonconstant $F\in k(U,V)$ and write $F=P/Q$ with coprime
$P,Q\in k[U,V]$, $Q\ne0$.  Suppose first that $Q$ is nonconstant.
We choose a sufficiently large positive integer $j$ with the following
three properties.

First, the polynomial
\[
 Q_j^{\mathrm{line}}(U)=Q\bigl(U,C+r_j(U-\beta)\bigr)
\]
is nonconstant.  To see this, let $d>0$ be the total degree of $Q$
and let $Q_d$ be its top homogeneous part.  The coefficient of $U^d$
in this restriction is $Q_d(1,r_j)$.  The polynomial $Q_d(1,T)$
is nonzero, so it vanishes for only finitely many slopes, and thus for
only finitely many integers $j$.

Second, neither omitted point lies on $Q=0$.  Write
\[
 Q(U,V)=\sum_{a=0}^D q_a(V)U^a,\qquad q_D\ne0.
\]
Transcendence of $C$ implies $q_D(C)\ne0$.  The asymptotics
\eqref{brp:endpoint-asymptotics} show that, for $D>0$,
\[
 Q(U_j)=q_D(C)(j\tau)^D+O(j^{D-1}),\qquad
 Q(V_j)=q_D(C)(j\tau)^D+O(j^{D-1}).
\]
For $D=0$ the two values tend to $q_0(C)\ne0$.  Therefore both
values are nonzero for all sufficiently large $j$.

Third, the line \eqref{brp:line} contains no common zero of $P,Q$.
Since $P,Q$ are coprime in $k[U,V]$, their common affine zero set is
finite, with all its points in $k^2$.  A point $(a,b)$ in this set
can lie on the line only if
\[
 r_j=\frac{b-C}{a-\beta}.
\]
Here $a-\beta\ne0$, and the ratio is nonzero because $b-C\ne0$.
Each common zero therefore excludes at most two integers $j$.

Choose $j$ satisfying all three properties.  The nonconstant polynomial
$Q_j^{\mathrm{line}}$ has a finite complex root.  At this root the
restricted numerator is nonzero, and the corresponding point is not
one of $U_j,V_j$.  Hence it is attained by \eqref{brp:images} at a
finite $X$, and $F(x_j(X),y_j(X))$ has a pole there.

If $Q$ is constant, $P$ is nonconstant.  The same leading-homogeneous-
part argument chooses $j$ for which
$P(U,C+r_j(U-\beta))$ is a nonconstant polynomial in $U$.  It has
a pole at the point at infinity of the line, which is attained at
$q=-B_j/A_j\ne0$, again giving a pole at a finite $X$.

Finally, $\beta,C\in\operatorname{Frac}(\E)$, so every continued
phase embedding extends to their rational field.  If $F(\beta,C)$
were integral over $\E$, applying this embedding to a monic
integral equation would give a monic equation with entire coefficients
for $F(x_j(X),y_j(X))$.  At a pole, the highest power has a uniquely
largest pole order, which is impossible.  Thus no such nonconstant
$F(\beta,C)$ is integral.  Constants in $k$ are already in
$\E$, proving the equality.
\end{proof}

\begin{corollary}[Earlier reciprocal exceptions are removed]
For every positive algebraic $t$ and every $\alpha\in k$,
\[
 1/(\beta(t)-\alpha),\quad 1/(C(t)-\alpha)\notin\Ebar.
\]
Both denominators are nonzero. At $t=2$, writing $V=\gamma+\kappa$,
\[
 \Ebar\cap k(V,C)=k.
\]
In particular the original Theorem B needs no assumption
$\gamma+\log(t/2)\ne0$ for its reciprocal conclusion.
\end{corollary}
\begin{proof}
These are nonconstant rational functions of the two algebraically
independent coordinates in Theorem~\ref{brp:main}.
\end{proof}

\begin{remark}[Scope]
The theorem is a two-variable rational-field statement, not just a
statement about two separate one-variable fields.  For example, it
excludes every nonconstant quotient $P(\beta,C)/Q(\beta,C)$ with
$P,Q\in k[U,V]$, including denominators that mix both coordinates.
The independent-scale extension is proved in
Theorem~\ref{mbmulti:main}; arbitrary rationally dependent scales
are not covered by that rational-field conclusion.
\end{remark}


\section{Mixed sums outside the integral closure of E-values}
Let $\E$ be the ring of values of E-functions at algebraic points, and
let $\Ebar$ be its integral closure in $\mathbb C$. This is a ring containing
$\Qa$. In particular it is a $\Qa$-vector subspace of $\mathbb C$.
We use the zero-winding morphisms $\Phi_\pm=\Phi_{\pm\pi/4}$
from Lemma~\ref{mbc:continued}.  Their explicit images are the
$j=0$ case of the preceding winding formulas and its conjugate.
To avoid conflicting with
coefficients of a linear combination, denote the image coefficients by
$A_t,B_t,C_t,D_t$. With $I=I_0(t)$, $K=K_0(t)$ and
$g_t=\gamma+\log(t/2)$, they are
\begin{align}
A_t&=e^{-t}(I+iK/\pi),& B_t&=-ie^tK/\pi,\nonumber\\
C_t&=e^{-t}(I_1(t)-iK_1(t)/\pi),&D_t&=ie^tK_1(t)/\pi,\label{intlin:images}\\
\Phi_-(I_0(t))(x)&=A_te^{tx}+B_te^{-tx}=:D_t^*(x),\nonumber\\
\Phi_-(I_1(t))(x)&=C_te^{tx}+D_te^{-tx},\nonumber\\
\Phi_-(Y(t))(x)&=g_tD_t^*(x)+e^tK e^{-tx}.\nonumber
\end{align}
Here $\Phi_\pm=\Phi_{\pm\pi/4}$. The plus images have conjugate
coefficients. Equation \eqref{eq:wronskian} gives
\begin{equation}\label{intlin:det}
A_tD_t-B_tC_t=\frac{i}{\pi t}.
\end{equation}
The complex constants $g_t$ in these formulas are coefficients of the
images; their appearance does not assert $\gamma\in\E$.

\begin{lemma}[Pole test]\label{intlin:pole}
If $u,v\in\E$, $u\ne0$, and some $\Phi_\theta$ and $x_0\in\mathbb C$ satisfy
$\Phi_\theta(u)(x_0)=0\ne\Phi_\theta(v)(x_0)$, then $v/u\notin\Ebar$.
\end{lemma}
\begin{proof}
A monic equation for $v/u$ over $\E$, multiplied by a suitable power
of $u$, would give an identity whose image at $x_0$ is
$\Phi_\theta(v)(x_0)^d=0$. No extension of the morphism to $\Ebar$ is used.
\end{proof}

\begin{lemma}[Ordered nearest zeros]\label{intlin:zeros}
The smallest positive imaginary part of a zero of $D_t^*$ is
\[
y(t)=\frac{1}{2t}\arctan\frac{\pi I_0(t)}{K_0(t)}.
\]
The function $y$ is strictly decreasing on $(0,\infty)$.
\end{lemma}
\begin{proof}
Writing $q(t)=\arctan(\pi I/K)\in(0,\pi/2)$, the equation
$e^{2tx}=-B_t/A_t$ shows that the zero imaginary parts are
$q(t)/(2t)+\pi n/t$, $n\in\mathbb Z$. The Wronskian gives
\[
q'(t)=\frac{\pi}{t(K^2+\pi^2I^2)}.
\]
It remains to show $tq'<q$. Put $r=\pi I/K$. If $IK\ge1$, then
\[
q=\arctan r>\frac{r}{1+r^2}
 =\frac{\pi IK}{K^2+\pi^2I^2}\ge tq'.
\]
If $IK<1$, use $I\ge1$ to obtain $r>\pi I^2\ge\pi$ and
\[
q>\arctan\pi>\frac1\pi
 \ge\frac{\pi}{K^2+\pi^2I^2}=tq'.
\]
Thus $y'=(tq'-q)/(2t^2)<0$.
\end{proof}

\begin{theorem}[Arbitrary finite mixed sums]\label{thm:poles}
For distinct positive algebraic $t_1,\ldots,t_k$ and
$a_j,b_j\in\Qa$, not all zero,
\begin{equation}\label{intlin:mixed}
\sum_{j=1}^k\bigl(a_jC(t_j)+b_j\beta(t_j)\bigr)\notin\Ebar.
\end{equation}
The conclusion is unchanged on adding any element of $\Ebar$.
\end{theorem}
\begin{proof}
Keep only active indices, for which $(a_j,b_j)\ne(0,0)$, and choose
the largest active argument $T$, writing $a_T,b_T$ for its coefficients.
At the nearest upper zero of
$D_T^*$, Lemma~\ref{intlin:zeros} shows that no other active denominator
vanishes. At this zero, the numerator belonging to scale $T$ has image
equal to a nonzero exponential factor times
\begin{equation}\label{intlin:residues}
a_T\frac{i}{\pi T A_T}+b_Te^TK_0(T).
\end{equation}
For the plus morphism, use the conjugate zero in the lower half-plane.
The analogous expression is
$-a_Ti/(\pi T\overline{A_T})+b_Te^TK_0(T)$.
If both expressions vanished, subtraction would give
$a_T(1/A_T+1/\overline{A_T})=0$. Since
$\operatorname{Re}A_T=e^{-T}I_0(T)>0$, it follows that $a_T=0$
and then $b_T=0$, a contradiction. This argument does not assume
that $a_T,b_T$ are real.

Write \eqref{intlin:mixed} over the common denominator
$u=\prod I_0(t_j)$, which is a nonzero E-value. Its numerator $v$
is also an E-value. For at least one of the two morphisms the chosen
zero of its image denominator leaves precisely the nonzero active
numerator above, multiplied by nonzero images of the other denominators.
Lemma~\ref{intlin:pole} applies. The assertion concerning addition follows
because $\Ebar$ is a ring.
\end{proof}

\begin{corollary}\label{cor:integralrank}
The $2k$ classes of $C(t_j),\beta(t_j)$ in $\mathbb C/\Ebar$ are linearly
independent over $\Qa$. The space of $\Qa$-linear forms in
$\gamma,C(t_j),\mu(t_j)$ whose value belongs to $\Ebar$ has dimension
at most one.
\end{corollary}
\begin{proof}
The first assertion is Theorem~\ref{thm:poles}. In a linear form,
replace $\mu(t_j)$ by $\beta(t_j)-\gamma$. Two independent forms
with values in $\Ebar$ would have a nonzero linear combination without
$\gamma$, contradicting the first assertion.
\end{proof}

\begin{remark}
This is non-integrality over the E-value ring, not transcendence over its
fraction field. All the $C(t_j)$ and $\beta(t_j)$ already belong to
$\operatorname{Frac}(\E)$. The distinctness hypothesis is weaker than the rational independence
needed for the multiscale rational-field classifications below.
\end{remark}



% Proof-bearing insert, 22 October 2026.
% Depends on the continued phase embeddings and explicit winding formulas
% in Keio_Bessel_Density_Monodromy_2026.tex. No joint-specialization
% independence input is required for this independent-scale statement.
% Requires the macros \Qa, \E, \Ebar and standard theorem environments.

\section{The full rational field at independent Bessel scales}

The independent-scale hypothesis permits one to freeze all but one
coordinate in the exponential torus. This avoids any generic-position
assumption about the synchronized family of product lines.

\begin{lemma}[Intersection with an exponential sublattice]
\label{mbmulti:cosets}
Let $G$ be a torsion-free abelian group and $\Lambda\subseteq G$ a
subgroup. Inside the fraction field of $\mathbb C[G]$ one has
\[
 \mathbb C[G]\cap\operatorname{Frac}(\mathbb C[\Lambda])
       =\mathbb C[\Lambda].
\]
\end{lemma}
\begin{proof}
Decompose the group algebra as the direct sum of its $G/\Lambda$
coset components. If $H\in\mathbb C[G]$ and $DH=N$ with nonzero
$D\in\mathbb C[\Lambda]$ and $N\in\mathbb C[\Lambda]$, then every
nonidentity coset component $H_c$ satisfies $DH_c=0$. The group algebra
is a domain, so all these components vanish. Thus $H$ is supported on
$\Lambda$.
\end{proof}

\begin{lemma}[The one-pair pole argument over a larger coefficient field]
\label{mbmulti:onepair}
Let $F\subseteq\mathbb C$ be a field. Fix a positive algebraic $t$,
and let $\beta=\beta(t)$ and $C=C(t)$ be algebraically independent
over $F$. With the constants and winding formulas of the one-point
Bessel theorem, put
\[
 x_j(q)=g+j\tau+\frac{i\pi B_j}{A_jq+B_j},\qquad
 y_j(q)=\frac{P_jq+Q_j}{A_jq+B_j}.
\]
For every nonconstant $f\in F(U,V)$ there is an integer $j$ such that
$f(x_j(q),y_j(q))$ has a pole at a point of $\mathbb C^\times$.
This is a statement about rational functions with complex coefficients;
no phase embedding is asserted to fix $F$.
\end{lemma}
\begin{proof}
The winding formulas give
\[
 y_j=C+r_j(x_j-\beta),\qquad
 r_j=\frac1{\pi^2tA_jB_j},\qquad
 r_j\sim-\frac1{4\pi^2tI_0(t)^2j^2}.
\]
Each fixed nonzero slope occurs for at most two integer windings.
The image of $q\in\mathbb C^\times$ is the projective line displayed
above with precisely two finite points omitted. Their first coordinates
are $j\tau+O(1)$, and their second coordinates are $C+O(j^{-1})$.
The point at infinity of the line is attained at
$q=-B_j/A_j\in\mathbb C^\times$.

Write $f=P/Q$ in lowest terms in $F[U,V]$. If $Q$ is nonconstant,
the restriction $Q(U,C+r_j(U-\beta))$ is nonconstant for all but
finitely many $j$: its top coefficient is $Q_d(1,r_j)$, where $Q_d$
is the nonzero highest homogeneous part. Write
$Q(U,V)=\sum_{a=0}^Dq_a(V)U^a$, $q_D\ne0$. Since $C$ is
transcendental over $F$, $q_D(C)\ne0$. At each omitted point the
value of $Q$ is therefore $q_D(C)(j\tau)^D+O(j^{D-1})$ when
$D>0$, and tends to $q_0(C)\ne0$ when $D=0$. Neither omitted point
lies on $Q=0$ for sufficiently large $j$.

The common affine zeros of the coprime polynomials $P,Q$ form a finite
set, all of whose coordinates are algebraic over $F$. For such a zero
$(a,b)$, transcendence of $\beta,C$ over $F$ implies
$a\ne\beta$, $b\ne C$. It lies on the selected line only when
\[
 r_j=\frac{b-C}{a-\beta},
\]
which excludes at most two windings per common zero. Choose a sufficiently
large $j$ avoiding these finitely many exceptions. A root of the
nonconstant restricted denominator is then an attained point where the
numerator is nonzero, and yields the required pole.

If $Q$ is constant, the same highest-homogeneous-part argument makes
the restriction of $P$ nonconstant for some $j$. Its pole at infinity
is attained at $q=-B_j/A_j$, proving this case as well.
\end{proof}

\begin{lemma}[Independence from the frozen line pencils]
\label{mbmulti:independence}
For positive algebraic $t_1,\ldots,t_m$ linearly independent over
$\mathbb Q$, the $2m$ numbers $\beta(t_i),C(t_i)$ are algebraically
independent over $\Qa$.
\end{lemma}
\begin{proof}
Fix an index $i$ and put $F_i=\Qa(\beta(t_l),C(t_l):l\ne i)$.
Suppose a nonzero polynomial $H\in F_i[U,V]$ vanishes at
$(\beta(t_i),C(t_i))$. Express its coefficients as rational functions
of the other actual generators, choosing denominators nonzero at those
values. Apply each continued embedding to this relation. The winding
formulas involve only $q_l=e^{2t_lX}$, which are algebraically independent
as exponential functions because the $t_l$ are rationally independent.
Hence the embedded relation is a rational identity in these independent
variables. Set $q_l=e^{2t_l}$ for every $l\ne i$. Evaluation compatibility
recovers every original coefficient of $H$; its coefficient denominators
remain nonzero. Thus
\[
 H\bigl(x_{i,j}(q_i),y_{i,j}(q_i)\bigr)=0
 \qquad\text{for every }j\in\mathbb Z.
\]
The polynomial $H$ therefore vanishes identically on infinitely many
distinct affine lines, with slopes $r_{i,j}$. A nonzero polynomial in
two variables has only finitely many distinct linear factors, a
contradiction. This proves independence of each pair over the field
of all other pairs. Successively deleting pairs gives total
transcendence degree $2m$.
\end{proof}

\begin{theorem}[Full rational exclusion at independent scales]
\label{mbmulti:main}
Let $t_1,\ldots,t_m$ be positive algebraic numbers linearly independent
over $\mathbb Q$. Then
\[
 \boxed{\quad
 \Ebar\cap\Qa\bigl(\beta(t_1),C(t_1),\ldots,
                          \beta(t_m),C(t_m)\bigr)=\Qa.
 \quad}
\]
\end{theorem}
\begin{proof}
Abbreviate $\beta_i=\beta(t_i)$ and $C_i=C(t_i)$; these are jointly
independent by Lemma~\ref{mbmulti:independence}. Suppose
$\xi=R(\beta_1,C_1,\ldots,\beta_m,C_m)$ is integral over $\E$, where
$R$ is a rational function with algebraic coefficients. The continued
embedding $\Phi_{-\pi/4}^{[j]}$ sends it to an element $H_j$ of the
integrally closed exponential-polynomial ring
\[
 \mathcal R=\mathbb C[\Qa_{\mathrm{add}}].
\]
The explicit winding formulas also put $H_j$ in the fraction field
of $\mathbb C[\Lambda]$, where $\Lambda=\sum_i\mathbb Z(2t_i)$.
Lemma~\ref{mbmulti:cosets} therefore gives
\[
 H_j\in\mathbb C[q_1^{\pm1},\ldots,q_m^{\pm1}],
       \qquad q_i=e^{2t_iX}.
\]
Here rational independence of the $t_i$ identifies the displayed ring
with a Laurent polynomial ring in $m$ independent variables. This
identification is essential: the following argument takes place on
the full exponential torus, not just on its analytic one-parameter
image.

If $R$ is nonconstant, choose an index $i$ on whose pair of variables
it depends. Set
\[
 F_i=\Qa(\beta_l,C_l:l\ne i).
\]
Joint independence identifies the specialization of $R$ in that pair
with a nonconstant rational function $f\in F_i(U,V)$, and makes
$\beta_i,C_i$ algebraically independent over $F_i$.

For every winding $j$, specialize $q_l=e^{2t_l}$ for all $l\ne i$.
Evaluation compatibility of the continued embeddings says that the
two rational images for scale $t_l$ then take the values $\beta_l,C_l$,
independently of $j$. Consequently the specialization of $H_j$ is
\[
 f\bigl(x_{i,j}(q_i),y_{i,j}(q_i)\bigr).
\]
This identity is valid as a rational identity: after clearing
denominators, the specialized denominator cannot vanish identically,
since its value at $q_i=e^{2t_i}$ is the nonzero original denominator
at the jointly independent numerical tuple. The specialized $H_j$
is a Laurent polynomial in $q_i$, and hence has no pole on
$\mathbb C^\times$.

Lemma~\ref{mbmulti:onepair}, applied with coefficient field $F_i$,
provides a winding for which the same function has a pole on
$\mathbb C^\times$. This contradiction shows that $R$ is constant.
Conversely, all algebraic constants already belong to $\E$.
\end{proof}

\begin{remark}
The theorem concerns the entire rational field and allows arbitrary
mixed numerators and denominators. The windings remain synchronized
across all scales. No separate continuation of individual factors,
generic product-line intersection, or independent choice of windings
is used. Rationally dependent scales are not covered by the freezing
argument, because their exponential coordinates satisfy relations.
\end{remark}


\begin{corollary}[Mixed polynomial exclusion]
\label{newmix:main}\label{thm:polynomial}
For positive algebraic $t_1,\ldots,t_m$ linearly independent over
$\mathbb Q$,
\[
 \Ebar\cap\Qa[\beta(t_1),C(t_1),\ldots,\beta(t_m),C(t_m)]=\Qa.
\]
The same equality holds after omitting any of the displayed variables,
in particular for the polynomial ring in the $\beta(t_i)$ alone.
\end{corollary}
\begin{proof}
A polynomial is a rational function, so Theorem~\ref{mbmulti:main}
applies.  This includes the earlier polynomial conclusions without
a separate torus-density proof.
\end{proof}


% New proof-bearing section for the five-paper synthesis.
% Requires the continued embeddings (mbc:continued) and normality/pole
% lemma (mbc:pole), or equivalent results earlier in the synthesis.
% The subgroup-support argument is recalled here to avoid label coupling.
% Revision insert: complex Bessel integral intersections.
% Place after the independent-positive-scale rational exclusion theorem.
% Depends on cf:nonzero, cf:joint, mbc:continued, mbc:pole,
% mbmulti:cosets, mbmulti:onepair, and mbmulti:main.

\section{Complex scales and rational fields without Euler's constant}
\label{cxi:section}

The positivity assumption in the preceding Bessel integral-intersection
theorems can be removed.  It is convenient to use a component description
that is independent of a particular sector-dependent sign.  Throughout
this section $k=\Qa$.  The functions $\beta(t)=Y(t)/I_0(t)$ and
$C(t)=I_1(t)/I_0(t)$ retain their branch-free definitions.  Auxiliary
occurrences of $K_0(t)$ and $\log(t/2)$ use compatible branches.

\begin{lemma}[Complex Bessel components and their winding pencils]
\label{cxi:components}
Let $t\in k^\times$.  Choose a nonexceptional direction
$\theta\in(-\pi/2,\pi/2)$ for the canonical component decompositions
of $I_0(tz)$ and $Y(tz)$.  Write
\[
 I=I_0(t),\quad J=I_1(t),\quad K=K_0(t),\quad L=K_1(t),\qquad
 g=\gamma+\log(t/2),\quad \tau=2\pi i.
\]
There are constants $c,d\in\mathbb C$, with $d\ne0$, such that the
two components of $I_0(tz)$, including their exponential factors, are
\begin{equation}\label{cxi:split}
 f_+(z)=cI_0(tz)+dK_0(tz),\qquad
 f_-(z)=(1-c)I_0(tz)-dK_0(tz).
\end{equation}
Their exponential phases are $t$ and $-t$, respectively.  Put
$a=(1-c)/d$ and $b=-c/d$.  The corresponding components of $Y(tz)$ are
\begin{equation}\label{cxi:Ysplit}
 (g+\log z+a)f_+(z),\qquad (g+\log z+b)f_-(z).
\end{equation}
For integer $j$, define
\begin{align*}
 A_j&=e^{-t}\bigl((c-d\tau j)I+dK\bigr),&
 B_j&=e^t\bigl((1-c+d\tau j)I-dK\bigr),\\
 P_j&=e^{-t}\bigl((c-d\tau j)J-dL\bigr),&
 Q_j&=e^t\bigl((1-c+d\tau j)J+dL\bigr).
\end{align*}
Then $A_jB_j\ne0$ for all but finitely many $j$.  For such $j$, the
continued images of $\beta(t),C(t)$, with $q=e^{2tX}$, are
\begin{equation}\label{cxi:images}
 x_j(q)=g+\tau j+a-\frac{B_j}{d(A_jq+B_j)},\qquad
 y_j(q)=\frac{P_jq+Q_j}{A_jq+B_j}.
\end{equation}
They parametrize the line
\begin{equation}\label{cxi:line}
 y=C(t)+r_j(x-\beta(t)),\qquad
 r_j=-\frac{d^2}{tA_jB_j}
     \sim-\frac{1}{4\pi^2tI_0(t)^2j^2}
       \quad(j\longrightarrow+\infty).
\end{equation}
Each fixed nonzero slope occurs at most twice.  The image of
$q\in\mathbb C^\times$ omits precisely two finite points of this
projective line.  Their first coordinates are $j\tau+O(1)$ and their
second coordinates are $C(t)+O(j^{-1})$.  The point at infinity is
attained at $q=-B_j/A_j\in\mathbb C^\times$.
\end{lemma}

\begin{proof}
The classical two-exponential expansion of the order-zero Bessel
equation gives nonzero formal components with phases $t$ and $-t$.
Their directional sums are independent solutions of the same
second-order equation; equivalently they are the appropriate rotated
Hankel solutions.  This description follows from the Bessel connection
and large-argument formulas in \cite[Sections 10.17, 10.27, and 10.40]{DLMF}.
Both phases occur: a single phase with a power-series amplitude on a
sector of opening greater than $\pi$ would contradict the two different
dominant exponentials on rays with positive and negative
$\operatorname{Re}(tz)$.  Directional summation commutes with the
differential equation.  Consequently the two sums are independent
solutions, and their sum is $I_0(tz)$.

Since $I_0(tz),K_0(tz)$ form a fundamental pair, the two components
therefore have the form \eqref{cxi:split}.  If $d=0$, both would be
constant multiples of $I_0(tz)$, contrary to their independence.
The constants $a,b$ satisfy $af_++bf_-=K_0(tz)$.  Hence the sum of
the expressions in \eqref{cxi:Ysplit} is
\[
 (g+\log z)I_0(tz)+K_0(tz)=Y(tz).
\]
Multiplication by $\log z$ preserves the corresponding exponential
phase and the logarithmic Nilsson--Gevrey class.  Uniqueness of the
canonical decomposition \cite[Theorem 2.3]{FR26} proves
\eqref{cxi:Ysplit}.

The classical continuation formulas are
\[
 M^jK_0(tz)=K_0(tz)-j\tau I_0(tz),\qquad
 M^jK_1(tz)=K_1(tz)+j\tau I_1(tz),
 \qquad M^j\log z=\log z+j\tau;
\]
see \cite[Section 10.34]{DLMF}.  Applying
Lemma~\ref{mbc:continued} to \eqref{cxi:split}, its derivative divided
by $t$, and \eqref{cxi:Ysplit} gives the coefficients $A_j,B_j,P_j,Q_j$
and the two rational images \eqref{cxi:images}.  The Wronskian gives
\[
 A_jQ_j-B_jP_j=d(IL+JK)=d/t\ne0.
\]
Moreover $I\ne0$ by Corollary~\ref{cf:nonzero}, also included in the
coordinate system of Theorem~\ref{cf:joint}.  Thus $A_jB_j$ is a
quadratic polynomial in $j$ with leading coefficient
$-d^2\tau^2I^2\ne0$.  Only finitely many integer windings need be
discarded; no assertion that every winding is nondegenerate is required.

Eliminating $q$ now yields the slope in \eqref{cxi:line}; evaluation
at $X=1$ gives its stated intercept.  The leading coefficient of
$A_jB_j$ gives the slope asymptotic, and its degree gives the bound
of two on repeated nonzero slopes.  At $q=\infty$ and $q=0$ the
respective omitted points are
\[
 (g+\tau j+a,P_j/A_j),\qquad (g+\tau j+b,Q_j/B_j).
\]
Their second coordinates tend to $J/I$ at rate $O(j^{-1})$.
Finally the common denominator vanishes at the displayed nonzero
value of $q$, while the nonzero Wronskian prevents degeneration of
the projective parametrization.
\end{proof}

\begin{theorem}[Exact rational intersection at complex independent scales]
\label{cxi:main}
Let $t_1,\ldots,t_m\in k^\times$ be linearly independent over
$\mathbb Q$.  Then
\[
 \boxed{\quad
 \Ebar\cap k\bigl(\beta(t_j),C(t_j):1\le j\le m\bigr)=k.
 \quad}
\]
The $2m$ displayed coordinates are algebraically independent over $k$.
In particular, the one-point conclusion holds at every nonzero
algebraic complex $t$.
\end{theorem}

\begin{proof}
Joint independence follows from Theorem~\ref{cf:joint}: rational
independence excludes repeated scales modulo sign.  Choose one common
direction that is nonexceptional for the finitely many Bessel functions
under consideration.  Such directions exist because each function
excludes only finitely many directions modulo $2\pi$.  Use this one
direction and a common integer winding $j$ for every scale.

We first explain precisely why the coefficient-field pole argument
survives the complex extension.  Suppose $F\subseteq\mathbb C$ and
$\beta(t),C(t)$ are independent over $F$.  For a nonconstant
$f=P/Q\in F(U,V)$ in lowest terms, the proof of
Lemma~\ref{mbmulti:onepair} uses only: infinitely many distinct nonzero
line slopes with finitely many repetitions; the stated asymptotics of
the two omitted points; and attainment of the point at infinity.
All three properties hold by Lemma~\ref{cxi:components} after finitely
many windings have been discarded.  Indeed, if $Q$ is nonconstant,
its restriction to the line is nonconstant for all but finitely many
slopes.  Writing $Q(U,V)=\sum_{a=0}^Dq_a(V)U^a$ with $q_D\ne0$,
transcendence of $C(t)$ over $F$ gives $q_D(C(t))\ne0$; hence the
denominator does not vanish at either omitted point for sufficiently
large $j$.  The finitely many common affine zeros of $P,Q$ have
coordinates algebraic over $F$ and exclude only finitely many slopes.
A root of the restricted denominator therefore gives an attained pole.
If $Q$ is constant, a nonconstant restriction of $P$ has a pole at
the attained point at infinity.  Thus some sufficiently large $j$
produces a pole in $\mathbb C^\times$ in every case.

Now let $\xi$ belong to the intersection in the theorem.  Normality
and Lemma~\ref{mbc:pole} put every continued image of $\xi$ in
$\mathbb C[k_{\rm add}]$.  The explicit images put it also in
\[
 \operatorname{Frac}\mathbb C[q_1^{\pm1},\ldots,q_m^{\pm1}],
 \qquad q_j=e^{2t_jX}.
\]
The $q_j$ are algebraically independent as exponential functions.
Lemma~\ref{mbmulti:cosets} consequently puts the image in the
displayed Laurent polynomial ring.  If $\xi$ is nonconstant, choose
a coordinate pair on which it depends, and freeze every other $q_l$
at $e^{2t_l}$.  Evaluation compatibility recovers the original
coefficient field generated by the other Bessel pairs.  This is the
same rational specialization used in Theorem~\ref{mbmulti:main}; its
denominator is nonzero because it is nonzero at the original numerical
tuple.  The resulting Laurent polynomial has no pole on
$\mathbb C^\times$, whereas the preceding one-pair argument supplies
a winding for which it must have such a pole.  This contradiction
proves that $\xi\in k$.  The opposite inclusion is immediate.
\end{proof}

\begin{corollary}[The Mahler coordinate and its derivative ratio]
\label{cxi:Mahler}
For every positive algebraic $\alpha$, put
\[
 M_\alpha=\frac{\pi Y_0(\alpha)}{2J_0(\alpha)}
              -\log(\alpha/2)-\gamma,
 \qquad R_\alpha=\frac{J_1(\alpha)}{J_0(\alpha)}.
\]
Then $M_\alpha,R_\alpha$ are algebraically independent and
\[
 \Ebar\cap k(M_\alpha,R_\alpha)=k.
\]
In particular $M_\alpha$ is not integral over $\E$.
\end{corollary}
\begin{proof}
The Bessel connection formulas give
\[
 \beta(i\alpha)=\gamma+\log(\alpha/2)
                -\frac{\pi Y_0(\alpha)}{2J_0(\alpha)},\qquad
 C(i\alpha)=i\frac{J_1(\alpha)}{J_0(\alpha)}.
\]
The denominators are nonzero by Corollary~\ref{cf:nonzero}, since
$J_0(\alpha)=I_0(i\alpha)$.  Apply Theorem~\ref{cxi:main} at
$t=i\alpha$ and make the invertible algebraic changes of coordinates
$M_\alpha=-\beta(i\alpha)$ and $R_\alpha=C(i\alpha)/i$.
\end{proof}

\begin{corollary}[Exact intersections for three coordinates without $\gamma$]
\label{cxi:gammafree}
For positive algebraic $a$, define
\begin{align*}
 D(a)&=\frac{K_0(2a)}{I_0(2a)}
             +\frac{\pi Y_0(2a)}{2J_0(2a)},\\
 C_I(a)&=\frac{I_1(2a)}{I_0(2a)},&
 C_J(a)&=\frac{J_1(2a)}{J_0(2a)}.
\end{align*}
If $a_1,\ldots,a_m$ are linearly independent over $\mathbb Q$, then
these $3m$ coordinates are jointly algebraically independent.  For
\[
 F=k\bigl(D(a_j),C_I(a_j),C_J(a_j):1\le j\le m\bigr)
\]
one has $\Ebar\cap F=k$.  Thus every nonconstant rational expression
in these coordinates lies outside $\Ebar$.
\end{corollary}
\begin{proof}
The $2m$ complex numbers $2a_j,2ia_j$ are linearly independent over
$\mathbb Q$, by taking real and imaginary parts.  The identities
\[
 D(a)=\beta(2a)-\beta(2ia),\qquad
 C_I(a)=C(2a),\qquad C_J(a)=C(2ia)/i
\]
show that $F$ is a subfield of the field in
Theorem~\ref{cxi:main}.  They also identify the displayed $3m$-tuple
with part of an invertible algebraic linear change of its $4m$
independent generators.  Both conclusions follow.
\end{proof}

\begin{corollary}[Radical exclusion]
\label{cxi:radicals}
Let $F$ be either a field from Theorem~\ref{cxi:main} or a subfield
containing $k$, including the fields in Corollaries~\ref{cxi:Mahler}
and~\ref{cxi:gammafree}.  Define the set
\[
 \operatorname{Rad}(F)=
 \{y\in\mathbb C:y^n\in F\text{ for some integer }n\ge1\}.
\]
Then $\Ebar\cap\operatorname{Rad}(F)=k$.
The notation denotes a set and does not assert that it is a field.
\end{corollary}
\begin{proof}
If $y\in\Ebar$ and $y^n\in F$, the ring property gives
$y^n\in\Ebar\cap F=k$.  Thus $y$ is algebraic over the algebraically
closed field $k$ and belongs to $k$.  The converse is immediate.
\end{proof}

\begin{remark}
At $a=1$, Corollary~\ref{cxi:gammafree} applies to the three explicit
coordinates $\kappa+\pi Y_0(2)/(2J_0(2))$, $C$, and $C_J$.
The argument uses the independent complex scales $2$ and $2i$.
It does not settle the separate problem of rationally dependent
positive scales, such as $2$ and $4$.
\end{remark}

\section{Exact multivariate exponential-integral intersections}
\label{ner:section}

Put $k=\overline{\mathbb Q}$.  As before, $\E$ is the
ring of $E$-function values and $\Ebar$ is its
integral closure in $\mathbb C$.  The continued phase embeddings
give an exact answer for a rational field generated by finitely
many exponential-integral values and their exponential units.
All independence assumptions on the arguments in this section
are over $\mathbb Q$, not merely pairwise distinctness.

\begin{theorem}[Exact exponential-integral rational intersection]
\label{ner:main}
Let $a_1,\ldots,a_m,b_1,\ldots,b_r$ be nonzero algebraic
numbers that are linearly independent over $\mathbb Q$, where
$m\ge1$ and $r\ge0$.  Put
\[
 E_i=e^{a_i},\qquad \eta_i=\Ein(a_i),
 \qquad Z_l=e^{b_l}.
\]
Then the $2m+r$ numbers $E_i,\eta_i,Z_l$ are algebraically
independent over $k$, and
\begin{equation}\label{ner:main-intersection}
 \boxed{\quad
 \Ebar\cap k(E_1,\eta_1,\ldots,E_m,\eta_m,
                                      Z_1,\ldots,Z_r)
 =k[E_1^{\pm1},\ldots,E_m^{\pm1},
                    \eta_1,\ldots,\eta_m,Z_1^{\pm1},\ldots,Z_r^{\pm1}].
 \quad}
\end{equation}
In particular,
\begin{align}
 \Ebar\cap k(\eta_1,\ldots,\eta_m)
       &=k[\eta_1,\ldots,\eta_m],\label{ner:ein-only}\\
 \Ebar\cap k(Z_1,\ldots,Z_r,\eta_1,\ldots,\eta_m)
       &=k[Z_1^{\pm1},\ldots,Z_r^{\pm1},
                                      \eta_1,\ldots,\eta_m].
                                      \label{ner:extra-units}
\end{align}
For a single arbitrary nonzero algebraic $a$, this gives
\begin{equation}\label{ner:paired-one-point}
 \Ebar\cap k(e^a,\Ein(a))
       =k[e^a,e^{-a},\Ein(a)].
\end{equation}
\end{theorem}

\subsection{Winding coordinates and independence}

Fix one admissible right-lateral summation direction for the
finitely many functions involved, and put $\tau=2\pi i$.
With compatible logarithm branches, there are constants
\[
 c_i=\gamma+\log a_i,\qquad
 d_i=E_i(\eta_i-c_i)
\]
for which the continued phase embeddings satisfy
\begin{equation}\label{ner:windings}
 \Phi^{[j]}(E_i)=u_i=e^{a_iX},\qquad
 \Phi^{[j]}(\eta_i)
       =c_i+j\tau+\frac{d_i-j\tau E_i}{u_i},\qquad
 \Phi^{[j]}(Z_l)=z_l=e^{b_lX}.
\end{equation}
These follow by continuing the decomposition
$\Ein(a_i z)=c_i+\log z+E_1^{\mathrm{tail}}(a_i z)$:
the entire function $\Ein(a_i z)$ is single valued,
whereas $\log z$ increases by $\tau$ per turn.  Different initial
compatible branches only change the fixed constants $c_i,d_i$.
Since the arguments $a_i,b_l$ are linearly independent over
$\mathbb Q$, the phase variables $u_i,z_l$ are algebraically
independent over $\mathbb C$ as exponential functions.

It is useful to replace $\eta_i$ by the $E$-value
$V_i=E_i\eta_i$.  Formula~\eqref{ner:windings} becomes
\begin{equation}\label{ner:line-image}
 \Phi^{[j]}(V_i)
       =(c_i+j\tau)u_i+d_i-j\tau E_i.
\end{equation}
Thus the image in each $(u_i,v_i)$ coordinate plane is the line
through the fixed point $(E_i,V_i)$ with slope $c_i+j\tau$,
parameterized by $u_i\in\mathbb C^\times$.

Here the claimed joint algebraic independence can be proved
directly, without importing the joint Bessel theorem.  Suppose
that a nonzero $P\in k[U_1,\ldots,U_m,V_1,\ldots,V_m,
Z_1,\ldots,Z_r]$ vanishes at these values.  Applying every
continued embedding gives the polynomial identity
\[
 P\bigl(u,\ c_i u_i+d_i+j\tau(u_i-E_i),\ z\bigr)=0
 \qquad(j\in\mathbb Z).
\]
The identity is polynomial in the independent phase variables
$u,z$ and in $j\tau$.  Let $D$ be the largest total degree of
$P$ in its $V$ variables, and let $P_D$ be the nonzero part
homogeneous of degree $D$ in those variables.  The coefficient
of $(j\tau)^D$ yields
\[
 P_D(u,u-E,z)=0.
\]
The numbers $E_1,\ldots,E_m$ are algebraically independent over
$k$ by Lindemann--Weierstrass: distinct integer combinations of
the $a_i$ give distinct algebraic exponential exponents.
Consequently the last displayed identity also holds as a
polynomial identity with $E_1,\ldots,E_m$ replaced by independent
indeterminates.  The invertible substitution $V_i=U_i-E_i$
would then give $P_D=0$, a contradiction.  This proves the
independence of $E_i,V_i,Z_l$, and hence of $E_i,\eta_i,Z_l$.

\subsection{A coefficient-field version of the one-point argument}

\begin{lemma}[A line pencil detects every nonunit denominator]
\label{ner:pencil}
Let $K\subset\mathbb C$ be a field, let $E$ be transcendental
over $K$, and fix $c,d\in\mathbb C$ and $\tau\ne0$.
For $j\in\mathbb Z$, put
\[
 v_j(u)=(c+j\tau)u+d-j\tau E.
\]
If $F\in K(U,V)$ does not belong to $K[U,U^{-1},V]$,
then for some integer $j$ the rational function $F(u,v_j(u))$
is well defined and has a pole at a point of $\mathbb C^\times$.
\end{lemma}
\begin{proof}
Multiplication by a power of $U$ does not affect the assertion.
We may write $F=P/Q$ with coprime $P,Q\in K[U,V]$, where
$Q$ is nonconstant and is not divisible by $U$.
The omitted point of the parameterized line is
$(0,d-j\tau E)$.  Since $Q(0,V)$ is a nonzero polynomial,
\[
 Q(0,d-j\tau E)\ne0
\]
for all but finitely many $j$; here $\tau E\ne0$.
If $q_D$ is the highest nonzero homogeneous part of $Q$ in
total degree, the leading coefficient of its restriction to
the line is $q_D(1,c+j\tau)$, nonzero for all but finitely
many $j$.  Thus that restriction is nonconstant.

The common zero set of the coprime polynomials $P,Q$ is finite,
and its coordinates are algebraic over $K$.  Every line in
this pencil passes through $(E,cE+d)$.  Since $E$ is not
algebraic over $K$, every common zero belongs to at most one
line of the pencil.  Avoid the resulting finitely many indices.
The restricted $Q$ now has a finite root away from $u=0$,
and the restricted $P$ is nonzero at that root.  This is a
genuine pole of $F(u,v_j(u))$ in $\mathbb C^\times$.
\end{proof}

\subsection{Freezing the other coordinates}

Let $\mathcal R$ again be the exponential-polynomial ring with
complex coefficients and algebraic phases, and let $\Lambda$
be the additive subgroup generated by the $a_i,b_l$.
The associated group ring is
\[
 \mathcal R_\Lambda
       =\mathbb C[u_1^{\pm1},\ldots,u_m^{\pm1},
                                z_1^{\pm1},\ldots,z_r^{\pm1}].
\]
We recall the subgroup-support fact
\begin{equation}\label{ner:coset}
 \mathcal R\cap\operatorname{Frac}(\mathcal R_\Lambda)
       =\mathcal R_\Lambda.
\end{equation}
This is Lemma~\ref{mbmulti:cosets}, applied to the additive group
of algebraic phases; no saturation hypothesis is needed.

\begin{proof}[Proof of the intersection statements]
Let
\[
 \xi\in\Ebar\cap k(E_1,V_1,\ldots,E_m,V_m,Z_1,\ldots,Z_r).
\]
All generators are $E$-values, so $\xi\in\operatorname{Frac}(\E)$.
Normality of $\mathcal R$ and the continued embeddings imply
$\Phi^{[j]}(\xi)\in\mathcal R$ for every $j$.
The explicit formulas place this image in
$\operatorname{Frac}(\mathcal R_\Lambda)$, so
\eqref{ner:coset} gives
\begin{equation}\label{ner:Laurent-image}
 \Phi^{[j]}(\xi)\in
 \mathbb C[u_1^{\pm1},\ldots,u_m^{\pm1},z_1^{\pm1},\ldots,z_r^{\pm1}].
\end{equation}

Fix an index $i$ and set
\[
 K_i=k(E_l,V_l\ (l\ne i),Z_1,\ldots,Z_r).
\]
Regard $\xi$ as a rational function in $E_i,V_i$ over $K_i$.
The established joint independence implies that $E_i$ is
transcendental over $K_i$.  In the Laurent-polynomial identity
for each winding, substitute $u_l=E_l$ for $l\ne i$ and
$z_l=Z_l$ for every $l$.  Formula~\eqref{ner:line-image} gives
\[
 (c_l+j\tau)E_l+d_l-j\tau E_l=V_l,
\]
so this substitution fixes every coefficient in $K_i$ for
every winding.  The right side of \eqref{ner:Laurent-image}
becomes a Laurent polynomial in $u_i$.
If the rational function in $E_i,V_i$ did not belong to
$K_i[E_i,E_i^{-1},V_i]$, Lemma~\ref{ner:pencil} would give
a winding whose specialized image has a pole in
$\mathbb C^\times$, a contradiction.
For precision, the substitutions can always be performed on
the product identity ``denominator times image equals
numerator''; the lemma chooses a winding for which the
restricted denominator is a nonzero polynomial.  No undefined
specialization of a rational expression is required.

Likewise, fix an extra exponential variable $Z_l$ and put all
other phase variables equal to their actual values.  A rational
function in $Z_l$ over the field of the other generators must
then be a Laurent polynomial in $Z_l$: otherwise it has a
finite nonzero pole, contradicting \eqref{ner:Laurent-image}.

These coordinatewise conclusions imply the global ring
membership
\[
 \xi\in k[E_1^{\pm1},\ldots,E_m^{\pm1},V_1,\ldots,V_m,
                                      Z_1^{\pm1},\ldots,Z_r^{\pm1}].
\]
Here is the elementary algebraic justification.  This Laurent
polynomial ring is a UFD.  In a reduced fraction, a nonunit
irreducible denominator factor remains a nonunit after
localizing all variables outside at least one pair $(E_i,V_i)$
or one variable $Z_l$.  If it contains a $V_i$, use that pair;
otherwise it has at least two Laurent monomials and some
exponential variable has different exponents in those monomials.
This contradicts the corresponding coordinatewise conclusion.

Conversely, every displayed generator and every allowed
exponential inverse is an $E$-value.  The whole displayed ring
therefore lies in $\E$.  Since $V_i=E_i\eta_i$ and
$E_i$ is a unit, this ring is exactly the right side of
\eqref{ner:main-intersection}.
Finally, algebraic independence allows us to intersect it with
$k(\eta_1,\ldots,\eta_m)$, or with
$k(Z_1,\ldots,Z_r,\eta_1,\ldots,\eta_m)$, by uniqueness of
Laurent expansions in the omitted $E_i$ variables.  This gives
\eqref{ner:ein-only} and \eqref{ner:extra-units}.
\end{proof}

\begin{remark}[Scope and the point $a=1$]
For $a=1$, the new exact ring is
\[
 \Ebar\cap k(e,\gamma+\delta/e)
          =k[e,e^{-1},\gamma+\delta/e].
\]
Thus adjoining the exponential unit does not create any
additional integral rational expressions: precisely the
polynomials in the exponential-integral value with Laurent
polynomial exponential coefficients survive.
The theorem does not assume an arithmetic classification of
$\gamma$ or $\delta$.  The multiargument proof requires the
stated $\mathbb Q$-linear independence; it does not assert the
same classification for arbitrary distinct, linearly dependent
arguments, or for arbitrary additional exponential units.
\end{remark}



\begin{corollary}[The one-variable exponential-integral field]
\label{mbc:ein}
For every nonzero algebraic $a$,
\[
 \Ebar\cap k(\Ein(a))=k[\Ein(a)].
\]
\end{corollary}
\begin{proof}
This is the case $m=1$, $r=0$ of
\eqref{ner:ein-only}.  A singleton containing a nonzero number is
linearly independent over $\mathbb Q$.
\end{proof}


\section{A combined exponential-integral and Bessel intersection}
\label{ncr:section}

\begin{corollary}[The exact mixed rational intersection]
\label{ncr:main}
Let $a_1,\ldots,a_m,b_1,\ldots,b_r$ be nonzero algebraic numbers,
and let $t_1,\ldots,t_s$ be positive algebraic numbers, where
$m,s\ge1$ and $r\ge0$. Assume that the entire list
\[
 a_1,\ldots,a_m,b_1,\ldots,b_r,2t_1,\ldots,2t_s
\]
is linearly independent over $\mathbb Q$. Set
\[
 E_i=e^{a_i},\qquad \eta_i=\Ein(a_i),\qquad
 Z_l=e^{b_l},\qquad \beta_j=\beta(t_j),\qquad C_j=C(t_j).
\]
These $2m+r+2s$ numbers are algebraically independent over
$k=\overline{\mathbb Q}$.  Write
\[
 \mathcal F=k(E_i,\eta_i,Z_l,\beta_j,C_j:
                  1\le i\le m,\ 1\le l\le r,\ 1\le j\le s).
\]
Then
\begin{equation}\label{ncr:intersection}
 \boxed{\quad \Ebar\cap\mathcal F
 =k[E_i^{\pm1},\eta_i,Z_l^{\pm1}:1\le i\le m,\ 1\le l\le r].\quad}
\end{equation}
Thus the Bessel coordinates introduce no additional integral rational
expressions, even when numerators and denominators mix all the displayed
families. Omitting the generators $E_i$ on the left gives the corresponding
right side $k[\eta_i,Z_l^{\pm1}]$; omitting both $E_i$ and $Z_l$ gives
$k[\eta_i]$.
\end{corollary}
\begin{proof}
Write the independent phase variables as
\[
 u_i=e^{a_iX},\qquad z_l=e^{b_lX},\qquad q_j=e^{2t_jX}.
\]
For winding $h\in\mathbb Z$, use the exponential-integral formulas
\eqref{ner:windings} and the Bessel formulas from
Lemma~\ref{mbmulti:onepair}. All images are rational functions of
these independent variables. At the original exponential coordinates
$u_i=E_i$, $z_l=Z_l$, and $q_j=e^{2t_j}$, evaluation compatibility
recovers the original numerical generators for every winding.

First we justify the joint independence needed in the coefficient-field
argument. The exponential and exponential-integral generators are
independent by Theorem~\ref{ner:main}. For each Bessel pair let $F_j$
be the field of all other displayed generators. If a nonzero polynomial
$H\in F_j[U,V]$ vanished at $(\beta_j,C_j)$, apply every continued
embedding and freeze all phase variables except $q_j$ at their original
exponential coordinates. As in Lemma~\ref{mbmulti:independence}, this
recovers the original coefficients of $H$ and forces $H$ to vanish on
infinitely many distinct affine Bessel lines. It is therefore the zero
polynomial, a contradiction. Coefficient fractions cause no difficulty:
choose their denominators nonzero at the original generators, clear
them before freezing, and evaluate those denominators at the same
original coordinates. Each pair is consequently independent over the
field of all the others. Together with the independence of the remaining
exponential-integral block, this proves the asserted total independence.

Now let $\xi$ belong to the intersection on the left of
\eqref{ncr:intersection}. Normality of the exponential-polynomial ring
and the subgroup-support identity \eqref{ner:coset}, equivalently
Lemma~\ref{mbmulti:cosets}, imply that every continued image of $\xi$
is a Laurent polynomial in $u_i,z_l,q_j$. Fix a Bessel index $j$ and
freeze all the other phase variables at their original exponential
coordinates. The result is a Laurent polynomial in $q_j$ and is the
one-pair rational image of $\xi$, viewed as a rational function in
$\beta_j,C_j$ over $F_j$. If this rational function were nonconstant,
Lemma~\ref{mbmulti:onepair} would provide a winding for which it has
a pole in $\mathbb C^\times$. This is impossible for a Laurent
polynomial. Thus $\xi$ is independent of that pair of variables.

Repeat over all Bessel pairs. Algebraic independence identifies the
numerical field with the corresponding pure rational function field,
so being independent of each pair leaves precisely
\[
 \xi\in k(E_1,\eta_1,\ldots,E_m,\eta_m,Z_1,\ldots,Z_r).
\]
Theorem~\ref{ner:main} now gives the right side of
\eqref{ncr:intersection}. Conversely, all its generators, including
the indicated exponential inverses, are $E$-values. The two statements
obtained by omitting generators follow by intersecting with the relevant
subfields and using algebraic independence, as in
\eqref{ner:ein-only}--\eqref{ner:extra-units}.
\end{proof}

\begin{remark}[The independence of the phases is essential to this proof]
The result uses one synchronized winding for every family; independent
windings are never assumed. The free phase variables make the coordinate
freezing possible. In particular, it does not cover the resonant field
$k(\Ein(1),\beta(2),C(2))$, since its phases $1$ and $4$
are rationally dependent.
\end{remark}



\section{Rational-field exclusion at independent positive scales}
\label{mre:section}

Throughout this section, $\E$ is the ring of E-function values at
algebraic points and $\Ebar$ is its integral closure in $\mathbb C$.
For a positive algebraic number $a$, put
\begin{align}
 h_a&=E_1^{\mathrm{tail}}(a):=\int_a^\infty\frac{e^{-u}}u\,du,
 &g_a&=\gamma+\log a,\label{mre:tail}\\
 \eta_a&=\Ein(a)=g_a+h_a,
 &B_a&=\beta(2a)=g_a+\frac{K_0(2a)}{I_0(2a)},\label{mre:values}\\
 D_a&=\eta_a-B_a
     =E_1^{\mathrm{tail}}(a)-\frac{K_0(2a)}{I_0(2a)},
 &C_a&=\frac{I_1(2a)}{I_0(2a)}.\label{mre:D}
\end{align}
The logarithm in these formulas is real. The values $\eta_a$ belong to
$\E$, while $B_a,C_a,D_a$ belong to $\operatorname{Frac}(\E)$.

\subsection{The images and their surjectivity}

We use the injective $\Qa$-algebra morphisms of Fischler--Rivoal
\cite[Proposition~3.1]{FR26}, in the two directions
$\Phi_\pm=\Phi_{\pm\pi/4}$. Their target is contained in
\[
 \mathcal R
 =\left\{\sum_{\lambda\in S}c_\lambda e^{\lambda x}:
       S\subset\Qa\text{ finite},\ c_\lambda\in\mathbb C\right\}.
\]
They satisfy $\Phi_\pm(\xi)(1)=\xi$ and extend to embeddings of
$\operatorname{Frac}(\E)$ into $\operatorname{Frac}(\mathcal R)$.
The coefficients of the exponential polynomials are allowed to be
arbitrary complex numbers; the phases are algebraic.

Set
\[
 I_a=I_0(2a),\qquad K_a=K_0(2a),\qquad
 q_a=e^a h_a,\qquad \rho_a=\frac{\pi I_a}{K_a},\qquad
 s_a=e^{-4a}(1-i\rho_a).
\]
All of $I_a,K_a,h_a,q_a,\rho_a$ are positive, and $s_a\ne0$.
With $w=e^{-ax}$, the sectorial Ein and Bessel decompositions give
\begin{equation}\label{mre:images}
 \Phi_-(\eta_a)=g_a+q_aw,\qquad
 \Phi_-(B_a)=g_a+\frac{i\pi w^4}{w^4-s_a}.
\end{equation}
For example, the first formula follows by applying the decomposition
$\Ein(az)=g_a+\log z+E_1^{\mathrm{tail}}(az)$;
the coefficient of the exponential part at $z=1$ is
$e^aE_1^{\mathrm{tail}}(a)$. For $\Phi_+$, conjugate the coefficients
in \eqref{mre:images}; $g_a$ and $q_a$ remain unchanged.

\begin{lemma}[Every target value has a nonzero finite preimage]
\label{mre:surjectivity}
Let $u,v\in\Qa^\times$. The rational function
\begin{equation}\label{mre:rational-image}
 R_{a,u,v}(w)
 =(u+v)g_a+u q_aw+\frac{vi\pi w^4}{w^4-s_a}
\end{equation}
maps $\mathbb C^\times$ onto $\mathbb P^1(\mathbb C)$.
Every finite target value has a preimage at which
$w^4-s_a\ne0$.
\end{lemma}
\begin{proof}
The four roots of $w^4=s_a$ are nonzero simple poles, since $v\ne0$.
Thus $\infty$ is attained. For a finite target $\alpha$, put
$c=(u+v)g_a-\alpha$. The numerator of $R_{a,u,v}(w)-\alpha$ is
\[
 N_\alpha(w)
 =u q_aw^5+(c+vi\pi)w^4-u q_as_aw-cs_a.
\]
This has degree five. If $c\ne0$, its constant coefficient is nonzero,
so it has a nonzero root. If $c=0$, its root at zero is simple because
the coefficient $-u q_as_a$ of $w$ is nonzero; its remaining four roots,
counted with multiplicity, are nonzero. At any root of $w^4=s_a$ we have
$N_\alpha(w)=vi\pi s_a\ne0$. A nonzero root of $N_\alpha$ is therefore
a nonpole preimage of $\alpha$.
\end{proof}

\subsection{Normality and the scale subgroup}

Distinct exponential functions with distinct constant phases are
linearly independent over $\mathbb C$. Thus $\mathcal R$ identifies
with the group algebra $\mathbb C[\Qa_{\mathrm{add}}]$. For an additive
subgroup $\Lambda\subset\Qa$, write $\mathcal R_\Lambda$
for its group algebra, realized by the exponential polynomials whose
phases lie in $\Lambda$.

\begin{lemma}[Two group-algebra facts]\label{mre:group-ring}
The ring $\mathcal R$ is integrally closed in its fraction field.
Moreover, for every additive subgroup $\Lambda\subset\Qa$,
\begin{equation}\label{mre:intersection}
 \mathcal R\cap\operatorname{Frac}(\mathcal R_\Lambda)
 =\mathcal R_\Lambda.
\end{equation}
No primitivity or saturation assumption on $\Lambda$ is required.
\end{lemma}
\begin{proof}
Normality is Lemma~\ref{mbc:pole}; the subgroup intersection is
Lemma~\ref{mbmulti:cosets}.  Their proofs use respectively finite
phase support and the decomposition into additive-subgroup cosets.
\end{proof}

In particular, if $\xi\in\operatorname{Frac}(\E)$ is integral over
$\E$, applying either morphism to a monic relation and using
Lemma~\ref{mre:group-ring} gives
\begin{equation}\label{mre:integral-image}
 \Phi_\pm(\xi)\in\mathcal R.
\end{equation}
This is stronger than merely asserting that the meromorphic image has
no finite pole.

\begin{theorem}[Full rational-field exclusion]
\label{mre:main}
Let $a_1,\ldots,a_m$ be positive algebraic numbers that are linearly
independent over $\mathbb Q$. Choose
$u_j,v_j\in\Qa^\times$, and put
\[
 Z_j=u_j\eta_{a_j}+v_jB_{a_j}\qquad(1\le j\le m).
\]
Then the $Z_j$ are algebraically independent over $\Qa$, and
\begin{equation}\label{mre:main-equation}
 \Ebar\cap\Qa(Z_1,\ldots,Z_m)=\Qa.
\end{equation}
Equivalently, for every
$F\in\Qa(T_1,\ldots,T_m)\setminus\Qa$, the value
$F(Z_1,\ldots,Z_m)$ is not integral over $\E$.
\end{theorem}
\begin{proof}
Put $w_j=e^{-a_jx}$ and
$\Lambda=\sum_{j=1}^m\mathbb Za_j$. Linear independence of the $a_j$
identifies $\mathcal R_\Lambda$ with
$\mathbb C[w_1^{\pm1},\ldots,w_m^{\pm1}]$ and its fraction field with
$\mathbb C(w_1,\ldots,w_m)$. Here the $w_j$ may therefore be treated as
independent variables. Formula~\eqref{mre:images} gives
\[
 \Phi_-(Z_j)=R_j(w_j),\qquad R_j=R_{a_j,u_j,v_j}.
\]
The substitution $T_j\mapsto R_j(w_j)$ is an injection of rational
function fields over $\mathbb C$: each $R_j$ is a nonconstant function
of its own independent variable. More explicitly,
Lemma~\ref{mre:surjectivity} supplies a nonpole point above every finite
tuple of target values, so no nonzero polynomial in the $T_j$ can vanish
after this substitution. Since $\Phi_-$ fixes $\Qa$, it follows that
the original values $Z_1,\ldots,Z_m$ are algebraically independent.

Suppose that $F(Z_1,\ldots,Z_m)$ is integral over $\E$. By
\eqref{mre:integral-image}, its image belongs to $\mathcal R$.
It also belongs to $\operatorname{Frac}(\mathcal R_\Lambda)$.
Lemma~\ref{mre:group-ring} therefore gives the rational identity
\begin{equation}\label{mre:laurent-identity}
 F(R_1(w_1),\ldots,R_m(w_m))
 =H(w_1,\ldots,w_m),
 \qquad
 H\in\mathbb C[w_1^{\pm1},\ldots,w_m^{\pm1}].
\end{equation}

Assume $F$ is nonconstant, and choose an index $j$ on which it depends.
Fix the other target variables at generic finite complex values
$\alpha_\ell$ ($\ell\ne j$), so that the specialized denominator of
$F$ is not the zero polynomial and the specialized rational function
\[
 f(T_j)=F(\alpha_1,\ldots,\alpha_{j-1},T_j,
                  \alpha_{j+1},\ldots,\alpha_m)
\]
is nonconstant. Such choices form a nonempty Zariski open set: if
$F=P/Q$ with coprime polynomials, retain a nonzero coefficient of $Q$
as a polynomial in $T_j$ and a nonzero coefficient of
$P_{T_j}Q-PQ_{T_j}$.
For each $\ell\ne j$, Lemma~\ref{mre:surjectivity} supplies
$c_\ell\in\mathbb C^\times$ with
$R_\ell(c_\ell)=\alpha_\ell$ and with the denominator of
$R_\ell$ nonzero at $c_\ell$.

Substitute these $c_\ell$ in \eqref{mre:laurent-identity}, after
clearing the rational denominators. This is legitimate because the
chosen coordinates are nonzero and are not poles of the $R_\ell$,
and the remaining denominator $Q(\alpha_1,\ldots,R_j(w_j),\ldots,
\alpha_m)$ is not identically zero. We obtain
\[
 f(R_j(w_j))=h(w_j),\qquad h\in\mathbb C[w_j^{\pm1}].
\]
The nonconstant rational function $f$ has a pole at some
$\alpha\in\mathbb P^1(\mathbb C)$. By
Lemma~\ref{mre:surjectivity}, $R_j$ takes the value $\alpha$ at a finite
nonzero point. A pole pulls back to a pole under a nonconstant rational
map, including when $\alpha=\infty$. Thus $f\circ R_j$ has a pole in
$\mathbb C^\times$, whereas the Laurent polynomial $h$ has none.
This contradiction proves that $F$ is constant, necessarily in $\Qa$.
Finally, $\Qa\subset\E$ proves the reverse inclusion in
\eqref{mre:main-equation}.
\end{proof}

\begin{corollary}[Tail--Bessel differences]\label{mre:differences}
For linearly independent positive algebraic $a_1,\ldots,a_m$,
\[
 \Ebar\cap\Qa(D_{a_1},\ldots,D_{a_m})=\Qa.
\]
In particular, every nonconstant rational function over $\Qa$ of
$D_1,D_{\sqrt2},D_{\sqrt3}$ is nonintegral over $\E$.
For example,
\[
 \frac{D_1D_{\sqrt2}+1}{D_{\sqrt3}-D_1}\notin\Ebar;
\]
the denominator is nonzero by algebraic independence.
For one point, $\Ebar\cap\Qa(D_a)=\Qa$ for every positive algebraic $a$.
\end{corollary}
\begin{proof}
Take $u_j=1$ and $v_j=-1$ in Theorem~\ref{mre:main}.
The three positive algebraic numbers $1,\sqrt2,\sqrt3$ are linearly
independent over $\mathbb Q$, as is seen by applying the independent
sign-changing automorphisms of $\mathbb Q(\sqrt2,\sqrt3)$.
\end{proof}

\begin{remark}[Two different hypotheses on the scales]
The joint independence theorem, Theorem~\ref{cf:joint}, already makes
$D_{a_1},\ldots,D_{a_m}$ algebraically independent whenever the positive
algebraic $a_j$ are merely \emph{distinct}. Indeed, it gives joint
independence of $\eta_{a_j},B_{a_j}$, and the differences are part of an
invertible linear change of coordinates.
The stronger rational-field exclusion proved here requires
\emph{linear independence over $\mathbb Q$}. Distinctness, or even
pairwise irrational ratios for three or more scales, does not justify
the independent Laurent variables used in its proof. No rational-field
exclusion at arbitrary distinct scales is asserted here.
\end{remark}

% INTERNAL PROOF CHECKPOINT, 22 October 2026.
% Not yet inserted into the authoritative integrated manuscript.
% Newly derived proof, independently checked internally; external review pending.
% Requires the manuscript's continued phase embeddings, pole obstruction,
% and joint algebraic independence of Ein(a), beta(t), C(t).
% No rational-independence or commensurability hypothesis is used below.

\section{The full Ein--Bessel rational intersection at arbitrary positive scales}
\label{fri:section}

Write $k=\overline{\mathbb Q}$, let $\mathbf E$ be the ring of
$E$-function values at algebraic arguments, and let
$\overline{\mathbf E}$ be the set of complex numbers integral over
$\mathbf E$.  In particular, $\overline{\mathbf E}$ is an integral
closure of a ring, not the algebraic closure of its fraction field.
This section uses two previously established inputs: the continued
phase embeddings $\Phi_j$ and their finite-pole obstruction, and the
joint algebraic independence over $k$ of
$\Ein(a),\beta(t),C(t)$ for positive algebraic $a,t$.

\begin{theorem}[Full rational intersection, including resonant scales]
\label{fri:main}
For every pair of positive algebraic numbers $a,t$, put
\[
 \eta=\Ein(a),\qquad
 \beta=\beta(t)=\gamma+\log(t/2)+\frac{K_0(t)}{I_0(t)},
 \qquad C=C(t)=\frac{I_1(t)}{I_0(t)}.
\]
Then
\begin{equation}\label{fri:main-equation}
 \boxed{\quad
 \overline{\mathbf E}\cap k(\eta,\beta,C)=k[\eta].
 \quad}
\end{equation}
There is no restriction on the ratio $t/a$.
\end{theorem}

\subsection{Normalized phase images and three local charts}

All logarithms of positive real numbers in this section are real.
Set
\[
 p=i\pi,\quad \tau=2\pi i=2p,\quad r=\frac{2t}{a},
 \quad I=I_0(t),\quad h=\frac1{tI^2},
\]
and, for positive integers $j$, set
\begin{align*}
 c_{a,j}&=\gamma+\log a+j\tau,&
 c_{t,j}&=\gamma+\log(t/2)+j\tau,\\
 d_j&=\eta-c_{a,j},&v_j&=\beta-c_{t,j},&
 s_j&=\frac{h}{v_j(v_j-p)}.
\end{align*}
Thus $d_j=-\tau j+O(1)$, $v_j=-\tau j+O(1)$,
$d_j/v_j\to1$, and $s_j\ne0$ for all sufficiently large $j$, with
$s_j=O(j^{-2})$ and $|s_j|\asymp j^{-2}$.

Use the exponential coordinate
\[
 q=e^{2t(1-X)}.
\]
The normalized continued images are
\begin{align}
 E_j(X):=\Phi_j(\eta)(X)
   &=c_{a,j}+d_j e^{a(1-X)},\label{fri:Eimage}\\
 B_j(X):=\Phi_j(\beta)(X)
   &=\beta+
     \frac{v_j(p-v_j)(q-1)}{p+v_j(q-1)},\label{fri:Bimage}\\
 D_j(X):=\Phi_j(C)(X)
   &=C-\frac{h(q-1)}{p+v_j(q-1)}
     =C+s_j\bigl(B_j(X)-\beta\bigr).
     \label{fri:Cimage}
\end{align}
Here $D_j$ denotes the image of the ratio $C$, not a Bessel
denominator.  These identities follow directly from the winding
coefficients in the one-point theorem.  For example, writing
$g_t=\gamma+\log(t/2)$ and $K/I=\beta-g_t$, the two normalized
coefficients of the $I_0$ image are $1-v_j/p$ and $v_j/p$.
Consequently
\[
 \Phi_j(I_0(t))(X)
 =I q^{-1/2}\frac{p+v_j(q-1)}p.
\]
The $Y$ image gives \eqref{fri:Bimage}; the Bessel Wronskian gives
\eqref{fri:Cimage}.  The Ein decomposition and its logarithmic
continuation give \eqref{fri:Eimage}.

Whenever $q$ is sufficiently close to $1$, use the holomorphic
logarithm with $\log 1=0$, and put
\[
 X=1-\frac{\log q}{2t},\qquad
 E_j=c_{a,j}+d_j q^{1/r}.
\]
This is only a local inverse coordinate: the original functions in
$X$ remain single valued.  In particular, no rationality of $r$ is
needed.

\begin{lemma}[Local degenerations]\label{fri:charts}
Put
\begin{equation}\label{fri:critical}
 u_* = \eta-\frac{p}{r}.
\end{equation}
The following statements hold as $j\to+\infty$.
\begin{enumerate}
\item If the image $B_j$ is prescribed to be a variable $x$ in a
fixed bounded disk, its local inverse near $X=1$ is given by
\begin{equation}\label{fri:beta-chart}
 q_j(x)=1+
 \frac{p(x-\beta)}{v_j(p-v_j-(x-\beta))}.
\end{equation}
On every such disk, uniformly,
\[
 E_j(x)\longrightarrow\eta,
 \qquad D_j(x)=C+s_j(x-\beta)\longrightarrow C.
\]
\item If the image $D_j$ is prescribed to be a variable $w$ in a
fixed closed disk not containing $C$, its local inverse near $X=1$
is given by
\begin{equation}\label{fri:C-chart}
 q_j(w)=1-\frac{p(w-C)}{h+v_j(w-C)}.
\end{equation}
Uniformly on that disk,
\[
 E_j(w)\longrightarrow u_*,
 \qquad B_j(w)=\beta+\frac{w-C}{s_j}.
\]
In particular, $|B_j(w)|\asymp j^2$ uniformly there.
\item If $\alpha\in k$, there is a local point near $X=1$ where
$E_j=\alpha$, and at that point
\begin{equation}\label{fri:eta-chart}
 B_j=
 \frac{r(\alpha-\eta)\tau}{p+r(\alpha-\eta)}j+O(1),
 \qquad D_j=C+O(j^{-1}).
\end{equation}
The displayed coefficient of $j$ is nonzero.
\item The Bessel images have a common simple pole at the local
point with
\begin{equation}\label{fri:pole-coordinate}
 q_j^{\rm pole}=1-\frac{p}{v_j}.
\end{equation}
The Ein image is holomorphic there, and its value $e_j^{\rm pole}$
satisfies $e_j^{\rm pole}\to u_*$.
\end{enumerate}
All local inverse maps used above are holomorphic and have nonzero
derivative for sufficiently large $j$.
\end{lemma}

\begin{proof}
Solving \eqref{fri:Bimage} for $q$ gives
\eqref{fri:beta-chart}.  Its denominator has size comparable to
$j^2$ on a fixed bounded disk, so $q_j(x)-1=O(j^{-2})$ uniformly.
The local expansion $q^{1/r}=1+(q-1)/r+O((q-1)^2)$, together with
$d_j=O(j)$, gives $E_j(x)=\eta+O(j^{-1})$.
The formula for $D_j(x)$ is exact.  The derivative of the nonconstant
M\"obius map \eqref{fri:beta-chart} does not vanish; explicitly, with
$b=x-\beta$, it is
\[
 q_j'(x)=\frac{p(p-v_j)}{v_j(p-v_j-b)^2}\ne0
\]
for all sufficiently large $j$.

Solving \eqref{fri:Cimage} gives \eqref{fri:C-chart}.  On a disk
avoiding $C$, its denominator does not vanish for large $j$, and
\[
 q_j(w)-1=-\frac{p}{v_j}+O(j^{-2})
\]
uniformly.  Hence $E_j(w)=\eta-pd_j/(rv_j)+O(j^{-1})\to u_*$.
The exact expression for $B_j(w)$ follows from the line identity in
\eqref{fri:Cimage}.  The inverse derivative is
\[
 q_j'(w)=-\frac{ph}{(h+v_j(w-C))^2}\ne0.
\]
In both charts $q_j\to1$, so the chosen logarithm defines a finite
holomorphic local coordinate $X$ with nonzero derivative.

For the third assertion prescribe
\[
 z_j=e^{a(1-X_j)}
     =\frac{\alpha-c_{a,j}}{d_j}
     =1+\frac{\alpha-\eta}{d_j}.
\]
This is nonzero and close to $1$.  Locally $q_j=z_j^r$, and thus
\[
 q_j-1=\frac{r(\alpha-\eta)}{d_j}+O(j^{-2}),
 \qquad
 p+v_j(q_j-1)\longrightarrow p+r(\alpha-\eta).
\]
The last limit is nonzero.  Otherwise, since $\eta$ and $r$ are real,
one would have $\pi=-r\operatorname{Im}\alpha$, contrary to the
transcendence of $\pi$.  Also $\alpha\ne\eta$, since $\eta$ is
transcendental.  Substitution into \eqref{fri:Bimage}--\eqref{fri:Cimage}
gives \eqref{fri:eta-chart}.  The local inverse of $E_j$ has derivative
$-1/(a(\alpha-c_{a,j}))$, which is nonzero.

Finally, \eqref{fri:pole-coordinate} is a nonzero simple zero of
$p+v_j(q-1)$ for large $j$.  The numerators of the Bessel polar terms
do not vanish there.  Moreover
\[
 e_j^{\rm pole}
 =c_{a,j}+d_j(1-p/v_j)^{1/r}
 =\eta-\frac{pd_j}{rv_j}+O(j^{-1})\longrightarrow u_*.
\]
The map from $X$ to $q$ has nonzero derivative, so these are also
simple poles in the original variable $X$.
\end{proof}

\begin{lemma}[Independence at the critical limit]\label{fri:critical-independent}
The numbers $u_*$ and $C$ are algebraically independent over $k$.
In particular, $u_*$ is transcendental.
\end{lemma}
\begin{proof}
The already established joint theorem makes $\eta,C$ algebraically
independent.  If $u_*$ were algebraic over $k(C)$, complex conjugation
would make $\overline{u_*}$ algebraic over the same field, because
$C$ is real and conjugation preserves $k$.  But
$(u_*+\overline{u_*})/2=\eta$, a contradiction.  Together with the
transcendence of $C$, this proves the assertion.
\end{proof}

\subsection{Eliminating every denominator}

\begin{proof}[Proof of Theorem~\ref{fri:main}]
The reverse inclusion in \eqref{fri:main-equation} follows from
$\eta\in\mathbf E$.  For the other inclusion, take
\[
 \xi=\frac{P(\eta,\beta,C)}{Q(\eta,\beta,C)}
       \in\overline{\mathbf E},
 \qquad P,Q\in k[U,V,W],\quad (P,Q)=1,
\]
with $Q\ne0$.  If $P=0$, the conclusion is immediate.  Otherwise,
joint independence of $\eta,\beta,C$ makes all nonzero polynomial
denominators nonzero at the actual tuple.  Each continued embedding
extends to this rational field.  By the pole obstruction, none of
its images of $\xi$ can have a finite pole.

Suppose first that $Q$ has an irreducible factor $q_0$ of positive
$V$-degree.  Gauss's lemma makes $q_0$ irreducible over $k(U,W)$,
and $q_0\nmid P$ makes $q_0$ and $P$ coprime there.  Their resultant
with respect to $V$ is a nonzero element of $k[U,W]$.  Joint
independence of $\eta,C$ ensures that this resultant, as well as
the leading $V$-coefficient of $q_0$, remains nonzero at $(\eta,C)$.
Consequently
\[
 q_0(\eta,V,C),\quad P(\eta,V,C)
\]
are coprime polynomials, the first of positive degree.
Choose one of its roots $x_0$ and a sufficiently small closed disk
about $x_0$ on which the second polynomial does not vanish and whose
boundary contains no zero of the first.

Use the first chart of Lemma~\ref{fri:charts}.  On this disk the
functions
\[
 q_0(E_j(x),x,D_j(x)),\qquad P(E_j(x),x,D_j(x))
\]
converge uniformly to the two displayed polynomials.  Rouch\'e's
theorem supplies a zero of the first function in the disk for every
sufficiently large $j$, while the second stays nonzero throughout
the disk.  At that zero the full denominator $Q$ vanishes and its
numerator does not.  The local inverse chart transfers the resulting
pole to a finite $X$.  This contradicts integrality.  Therefore
\begin{equation}\label{fri:denom-no-beta}
 Q\in k[U,W].
\end{equation}

Next let $q_0(U,W)$ be an irreducible factor of $Q$ of positive
$W$-degree.  Expand the numerator as
\[
 P(U,V,W)=\sum_{\ell=0}^M P_\ell(U,W)V^\ell.
\]
Because $q_0\nmid P$, there is a largest index $m$ for which
$q_0\nmid P_m$.  Every $P_\ell$ with $\ell>m$ is divisible by
$q_0$.  Irreducibility and Gauss's lemma give
\[
 R(U)=\operatorname{Res}_W(q_0,P_m)\in k[U]\setminus\{0\}.
\]
Lemma~\ref{fri:critical-independent} implies that $R(u_*)\ne0$ and
that the leading $W$-coefficient of $q_0$ is nonzero at $u_*$.
Thus $q_0(u_*,W)$ is nonconstant and coprime to $P_m(u_*,W)$.
The same lemma gives $q_0(u_*,C)\ne0$.

Choose a root $w_0$ of $q_0(u_*,W)$.  It differs from $C$.  Choose
a small closed disk about $w_0$, avoiding $C$, on whose boundary
$q_0(u_*,W)$ does not vanish and throughout which $P_m(u_*,W)$
does not vanish.  In the second chart of Lemma~\ref{fri:charts},
Rouch\'e's theorem gives a zero $w_j$ of $q_0(E_j(w),w)$ in this
disk for all large $j$.  On the disk $P_m(E_j(w),w)$ is uniformly
bounded away from zero, all the other coefficient functions are
bounded, and $|B_j(w)|\asymp j^2$.

At $w=w_j$, every term with $\ell>m$ vanishes exactly, because
$q_0$ divides its coefficient.  Therefore
\[
 P(E_j(w_j),B_j(w_j),w_j)
 =B_j(w_j)^m\left(P_m(E_j(w_j),w_j)+O(|B_j(w_j)|^{-1})\right)
 \ne0
\]
for large $j$; when $m=0$, the same assertion holds with no lower
terms.  The denominator vanishes there.  This is again a finite
pole, contradicting integrality.  Combining this with
\eqref{fri:denom-no-beta}, we obtain $Q\in k[U]$.

If $Q$ were nonconstant, choose a root $\alpha\in k$.  Coprimality
gives $P(\alpha,V,W)\ne0$.  Let $m$ be its degree in $V$, and
write its leading coefficient as $b_m(W)\in k[W]\setminus\{0\}$.
Transcendence of $C$ gives $b_m(C)\ne0$.  The third chart of
Lemma~\ref{fri:charts} sets $E_j=\alpha$, makes $B_j$ diverge
linearly with nonzero coefficient, and makes $D_j\to C$.
Consequently $P(\alpha,B_j,D_j)\ne0$ for all large $j$.
The denominator vanishes at these points, giving the same
contradiction.  Thus $Q$ is constant.

\subsection*{Eliminating polynomial dependence on the Bessel coordinates}

It remains to treat $P(\eta,\beta,C)$.  Suppose that $P$ has
positive total degree $D$ in $(V,W)$, and write its highest
homogeneous part as
\[
 P_D(U,V,W)=\sum_{b=0}^{D}p_b(U)V^{D-b}W^b,
 \qquad
 H_D(U,S)=\sum_{b=0}^{D}p_b(U)S^b\ne0.
\]
At the simple Bessel pole from Lemma~\ref{fri:charts}, the exact
line identity reads
\[
 D_j=C+s_j(B_j-\beta).
\]
The coefficient of $B_j^D$ in the polar expansion is therefore
$H_D(e_j^{\rm pole},s_j)$.  Let $b_0$ be the smallest index with
$p_{b_0}\ne0$.  Since $s_j\to0$ and
$e_j^{\rm pole}\to u_*$, we have
\[
 \frac{H_D(e_j^{\rm pole},s_j)}{s_j^{b_0}}
 \longrightarrow p_{b_0}(u_*)\ne0.
\]
The last inequality uses the transcendence of $u_*$.  Hence, for
all sufficiently large $j$, an order-$D$ pole survives.  Terms of
lower Bessel total degree cannot cancel it, and the Ein image is
holomorphic at the point.  This contradicts the pole obstruction.
Thus $D=0$, so $P\in k[U]$, completing the proof.
\end{proof}

\begin{corollary}[The matched-scale question]\label{fri:matched}
For every positive algebraic $a$,
\[
 \overline{\mathbf E}\cap
 k\bigl(\Ein(a),\beta(2a),C(2a)\bigr)=k[\Ein(a)].
\]
In particular, with $\eta=\Ein(1)$, the previously unresolved
intersection $\overline{\mathbf E}\cap k(\eta,\beta(2),C(2))$
equals $k[\eta]$.
\end{corollary}

\begin{remark}[Scope and review status]
This argument concerns one Bessel block together with one Ein value;
it does not assert the corresponding full rational-field statement
for several rationally dependent Bessel scales.  It also establishes
neither the irrationality of $\gamma$ nor full algebraic independence
of the original four constants.  The proof above is a newly derived,
internally audited extension of the manuscript's established phase
embeddings and joint independence theorem.  Historical priority and
external expert verification remain separate questions.
\end{remark}

% Proof-bearing revision insert, 22 October 2026.
% Owned by resonant_closure.  Requires new_full_resonant_intersection.tex.
% Newly derived and internally audited; external review and priority pending.

\section{Adjoining the matching exponential}
\label{fre:section}

The full rational intersection remains exact after adjoining the
exponential at the Ein argument.  The inverse exponential is necessarily
present: both $e^a$ and $e^{-a}$ are units of the $E$-value ring.

\begin{theorem}[Exact exponential--Ein--Bessel intersection]
\label{fre:main}
Let $a,t$ be arbitrary positive algebraic numbers, and put
\[
 E=e^a,\qquad \eta=\Ein(a),\qquad
 \beta=\beta(t),\qquad C=C(t),\qquad k=\overline{\mathbb Q}.
\]
Then
\begin{equation}\label{fre:main-equation}
 \boxed{\quad
 \overline{\mathbf E}\cap k(E,\eta,\beta,C)
       =k[E,E^{-1},\eta].
 \quad}
\end{equation}
No rational-independence or commensurability condition on $a,t$
is required.
\end{theorem}

\subsection{Limits with the exponential coordinate}

Retain the notation $p=i\pi$, $\tau=2\pi i$, $r=2t/a$, $h$,
$c_{a,j}$, $c_{t,j}$, $d_j$, $v_j$, $s_j$ from
Section~\ref{fri:section}.  In particular,
\[
 u_* = \eta-p/r,\qquad q=e^{2t(1-X)}.
\]
Write $E_j,B_j,D_j$ for the continued images of $\eta,\beta,C$
as in \eqref{fri:Eimage}--\eqref{fri:Cimage}; the letter $E_j$
continues to denote the Ein image, not the exponential image.
For clarity, denote the latter by
\begin{equation}\label{fre:exponential-image}
 \chi_j(X):=\Phi_j(e^a)(X)=e^{aX}.
\end{equation}
This expression is independent of the winding index.  Near $X=1$,
using the same logarithmic branch as before,
\begin{equation}\label{fre:exponential-q}
 \chi_j=E q^{-1/r}.
\end{equation}

The joint independence theorem over the field generated by algebraic
exponentials, together with the transcendence of $E=e^a$, makes
\begin{equation}\label{fre:independent-tuples}
 E,\eta,\beta,C
 \quad\hbox{algebraically independent over }k.
\end{equation}
We also have algebraic independence of $E,u_*,C$.  Indeed, $E,C$
are real and algebraically independent.  If $u_*$ were algebraic
over $k(E,C)$, complex conjugation would make $\overline{u_*}$
algebraic over that same field, and their average $\eta$ would be
algebraic over $k(E,C)$, contradicting
\eqref{fre:independent-tuples}.  In particular, $E,u_*$ are
algebraically independent.

The first two charts of Lemma~\ref{fri:charts} now give, uniformly
on the fixed closed disks specified there,
\begin{align}
 B_j=x\text{ bounded}:\quad
   &(\chi_j,E_j,D_j)\longrightarrow(E,\eta,C),
     \label{fre:beta-limit}\\
 D_j=w\text{ away from }C:\quad
   &(\chi_j,E_j)\longrightarrow(E,u_*),
   \qquad B_j=\beta+(w-C)/s_j.
     \label{fre:C-limit}
\end{align}
The assertion for $\chi_j$ follows from $q_j\to1$.

We also need the Ein chart with a variable target, rather than only
an algebraic target.  Prescribe $E_j=u$ and set
\begin{equation}\label{fre:Ein-chart}
 z_j(u)=\frac{u-c_{a,j}}{d_j}
       =1+\frac{u-\eta}{d_j},\qquad
 q_j(u)=z_j(u)^r,
 \qquad \chi_j(u)=\frac{E}{z_j(u)}.
\end{equation}
On every fixed closed disk avoiding $u_*$, the chosen branches are
holomorphic for all large $j$, and, uniformly,
\begin{align}
 \chi_j(u)&\longrightarrow E,&D_j(u)&\longrightarrow C,
       \label{fre:Ein-limits}\\
 \frac{B_j(u)}j&\longrightarrow
 F(u):=\frac{\tau r(u-\eta)}{p+r(u-\eta)}.
       \label{fre:Ein-beta-growth}
\end{align}
This follows by the same expansion used for
\eqref{fri:eta-chart}.  On a disk also avoiding $\eta$, the limiting
function $F$ is bounded away from zero, so $|B_j(u)|\asymp j$.
The local inverse has derivative
$-1/(a(u-c_{a,j}))\ne0$.

Finally, at the Bessel pole \eqref{fri:pole-coordinate},
\begin{equation}\label{fre:pole-limits}
 \chi_j\longrightarrow E,
 \qquad E_j\longrightarrow u_*,
 \qquad s_j\longrightarrow0.
\end{equation}
The exponential image is holomorphic and nonzero there.

\subsection{Proof of the exact intersection}

\begin{proof}[Proof of Theorem~\ref{fre:main}]
The ring on the right of \eqref{fre:main-equation} is contained in
$\mathbf E$, proving one inclusion.  For the reverse inclusion, let
\[
 \xi=\frac{P(E,\eta,\beta,C)}{Q(E,\eta,\beta,C)}
 \in\overline{\mathbf E},
 \qquad P,Q\in k[S,U,V,W],\quad (P,Q)=1,
\]
with $Q\ne0$.  The case $P=0$ is immediate.  Joint independence
makes every nonzero denominator at the actual tuple nonzero.  The
continued embeddings extend to the resulting fraction field, and
integrality excludes every finite pole of an image of $\xi$.

\paragraph{Denominators involving $V$.}
Let $q_0$ be an irreducible factor of $Q$ of positive $V$-degree.
Gauss's lemma and $q_0\nmid P$ imply that
\[
 \operatorname{Res}_V(q_0,P)
      \in k[S,U,W]\setminus\{0\}.
\]
Independence of $E,\eta,C$ makes this resultant and the leading
$V$-coefficient of $q_0$ nonzero at $(E,\eta,C)$.
Hence $q_0(E,\eta,V,C)$ is a nonconstant polynomial coprime to
$P(E,\eta,V,C)$.

Choose a root and a small closed disk around it, with no boundary
root and with the numerator nonzero throughout the disk.  In the
finite-$B_j$ chart \eqref{fre:beta-limit}, the corresponding image
polynomials converge uniformly to these two polynomials.  Rouch\'e's
theorem gives a denominator zero while the numerator stays nonzero.
The holomorphic local inverse transfers this uncancelled pole to a
finite $X$, a contradiction.  Thus $Q\in k[S,U,W]$.

\paragraph{Denominators involving $W$.}
Let $q_0(S,U,W)$ be an irreducible factor of $Q$ of positive
$W$-degree.  Expand
\[
 P(S,U,V,W)=\sum_\ell P_\ell(S,U,W)V^\ell,
\]
and choose the largest $m$ with $q_0\nmid P_m$.  Every higher
coefficient is divisible by $q_0$.  The resultant
\[
 R(S,U)=\operatorname{Res}_W(q_0,P_m)
          \in k[S,U]\setminus\{0\}
\]
remains nonzero at $(E,u_*)$, by their algebraic independence.
The leading $W$-coefficient remains nonzero as well.  Thus
$q_0(E,u_*,W)$ is nonconstant and coprime to $P_m(E,u_*,W)$.
Independence of $E,u_*,C$ also gives $q_0(E,u_*,C)\ne0$.

Choose a root $w_0\ne C$ and a sufficiently small closed disk
avoiding $C$ on which $P_m(E,u_*,w)$ is nonzero and whose boundary
contains no root of $q_0(E,u_*,w)$.  Use the finite-$D_j$ chart
\eqref{fre:C-limit}.  Rouch\'e's theorem supplies roots $w_j$ of
$q_0(\chi_j(w),E_j(w),w)$ in the disk.  At these roots, all
coefficients with $\ell>m$ vanish exactly.  The coefficient of
$B_j^m$ is bounded away from zero, while lower coefficients remain
bounded and $|B_j|\asymp j^2$.  The numerator is therefore nonzero
for all large $j$.  This is an uncancelled pole and a contradiction.
It follows that $Q\in k[S,U]$.

\paragraph{Denominators involving $U$.}
Let $q_0(S,U)$ be an irreducible factor of $Q$ of positive
$U$-degree.  Using the same expansion in $V$, choose the largest
$m$ for which $q_0\nmid P_m(S,U,W)$.  Irreducibility persists after
adjoining the independent variable $W$, so Gauss's lemma gives
\[
 R(S,W)=\operatorname{Res}_U(q_0,P_m)
          \in k[S,W]\setminus\{0\}.
\]
This resultant remains nonzero at $(E,C)$.  Also the leading
$U$-coefficient of $q_0$ remains nonzero at $E$.  Thus
$q_0(E,U)$ is nonconstant and coprime to $P_m(E,U,C)$.

Every root $u_0$ of $q_0(E,U)$ differs from both $\eta$ and
$u_*$.  Indeed, either equality would contradict respectively
the independence of $E,\eta$ or of $E,u_*$.  Choose a small closed
disk around one such root, avoiding both exceptional values, on
which $P_m(E,u,C)$ stays nonzero and whose boundary contains no
root of $q_0(E,u)$.  Apply the Ein chart
\eqref{fre:Ein-chart}--\eqref{fre:Ein-beta-growth}.
Rouch\'e's theorem gives zeros of $q_0(\chi_j(u),u)$ in the disk.
At those zeros, all higher $V$-coefficients vanish exactly,
the coefficient of $B_j^m$ stays nonzero, and $|B_j|\asymp j$.
The same leading-term argument gives a nonzero numerator at a
denominator zero.  This contradiction proves $Q\in k[S]$.

\paragraph{The nonzero exponential roots.}
Suppose $\alpha\in k^\times$ is a root of $Q$.  Coprimality gives
\[
 P(\alpha,U,V,W)\ne0.
\]
Choose any finite $X_0$ satisfying $e^{aX_0}=\alpha$.  At this point
put
\[
 z_0=e^{a(1-X_0)}=\frac{E}{\alpha},\qquad
 q_0=e^{2t(1-X_0)}.
\]
Then $z_0\ne1$, since $E$ is transcendental.  Moreover $q_0\ne1$:
the positive real number $|\alpha|$ is algebraic, whereas $E$ is
transcendental, and
\[
 |q_0|=\left(\frac{E}{|\alpha|}\right)^{2t/a}\ne1.
\]
For this fixed $X_0$, the exact phase formulas give
\begin{equation}\label{fre:fixed-exp-growth}
 E_j(X_0)=\tau j(1-E/\alpha)+O(1),\qquad
 B_j(X_0)=\tau j+O(1),\qquad
 D_j(X_0)\longrightarrow C.
\end{equation}
For large $j$ their Bessel denominators are nonzero, since
$p+v_j(q_0-1)$ has a nonzero linear term in $j$.

Let $D$ be the total degree in $(U,V)$ of
$P(\alpha,U,V,W)$, and let $H_D(U,V,W)$ be its nonzero highest
homogeneous part in those two variables.  The substitution
$H_D(T,1,W)$ is nonzero.  Algebraic independence of $E,C$ implies
\[
 H_D(1-E/\alpha,1,C)\ne0.
\]
Consequently \eqref{fre:fixed-exp-growth} gives
\[
 \frac{P(\alpha,E_j(X_0),B_j(X_0),D_j(X_0))}{(\tau j)^D}
 \longrightarrow H_D(1-E/\alpha,1,C)\ne0.
\]
The numerator is nonzero for large $j$, while its denominator
$Q(\chi_j(X_0))=Q(\alpha)$ vanishes.  Since
$\chi_j'(X_0)=a\alpha\ne0$, this produces a finite pole.
Thus $Q$ has no nonzero root.  As $k$ is algebraically closed,
$Q(S)$ is a nonzero constant times a nonnegative power of $S$.
We have proved
\[
 \xi\in k[E,E^{-1},\eta,\beta,C].
\]

\paragraph{The remaining Laurent polynomial.}
Write $\xi$ as a polynomial in $\beta,C$ with coefficients in
$k[E,E^{-1},\eta]$.  Suppose its Bessel total degree is $D>0$.
Its highest homogeneous part can be encoded as
\[
 H_D(S,U,Z)=\sum_{b=0}^{D}p_b(S,U)Z^b,
 \qquad p_b\in k[S,S^{-1},U],
\]
where $Z$ records the slope in the image line
$W=s_jV+(C-s_j\beta)$, with $\beta,C$ denoting the fixed original values.
Let $b_0$ be the smallest index with $p_{b_0}\ne0$.
At the Bessel pole, \eqref{fre:pole-limits} gives
\[
 \frac{H_D(\chi_j^{\rm pole},E_j^{\rm pole},s_j)}{s_j^{b_0}}
 \longrightarrow p_{b_0}(E,u_*)\ne0.
\]
The last inequality follows from independence of $E,u_*$; negative
powers of $E$ cause no difficulty because $E\ne0$ and the image
$\chi_j$ is holomorphic and nonzero.  Therefore an order-$D$
pole survives for large $j$, exactly as in the polynomial part of
Theorem~\ref{fri:main}.  Lower Bessel total degrees cannot cancel it.
This contradicts integrality and forces $D=0$.
Thus $\xi\in k[E,E^{-1},\eta]$, proving
\eqref{fre:main-equation}.
\end{proof}

\begin{corollary}[Matched exponential and Bessel scales]
\label{fre:matched}
For every positive algebraic $a$,
\[
 \overline{\mathbf E}\cap
 k\bigl(e^a,\Ein(a),\beta(2a),C(2a)\bigr)
 =k[e^a,e^{-a},\Ein(a)].
\]
In particular, at $a=1$ the full resonant field obtained by adjoining
$e$ to $k(\eta,\beta(2),C(2))$ has exactly the Laurent-polynomial
intersection $k[e,e^{-1},\eta]$.
\end{corollary}

\begin{remark}[What is and is not included]
The exponential adjoined here is $e^a$, matching the Ein argument.
The theorem permits an arbitrary positive algebraic Bessel argument
$t$.  It does not assert the analogous full rational-field theorem
for several dependent Bessel blocks or for an arbitrary collection
of additional exponential coordinates.  Its proof uses the
continued phase embeddings and the previously established joint
independence theorem; it does not prove irrationality of $\gamma$.
\end{remark}

% New closure: arbitrary distinct positive logarithmic-derivative scales.
% Dependencies: cf:joint, mbc:continued, mbc:pole, ner:windings.
% This fragment deliberately makes no full multiscale (beta,C) assertion.
\section{Rational intersections at dependent logarithmic-derivative scales}
\label{dcr:section}

Rational dependence between Bessel phases does not obstruct the exact
intersection for their logarithmic derivatives. The next result allows,
for example, the two scales $2$ and $4$, together with an exponential
integral at an arbitrary positive algebraic point.

\begin{theorem}[Arbitrary distinct positive logarithmic-derivative scales]
\label{dcr:main}
Let $a,t_1,\ldots,t_m$ be positive algebraic numbers, with the $t_i$
pairwise distinct, and put
\[
 \eta_a=\Ein(a),\qquad C_i=\frac{I_1(t_i)}{I_0(t_i)},
 \qquad k=\overline{\mathbb Q}.
\]
Then
\begin{equation}\label{dcr:intersection}
 \boxed{\quad
 \Ebar\cap k(\eta_a,C_1,\ldots,C_m)=k[\eta_a].
 \quad}
\end{equation}
No rational independence among $a,t_1,\ldots,t_m$ is required.
\end{theorem}

We first record three local forms of the continued phase images. They
use the same winding for all coordinates, so they do not presume
independent exponential variables.

\begin{lemma}[Local charts for synchronized windings]
\label{dcr:charts}
Write $\eta=\eta_a$, $p=i\pi$, $\tau=2\pi i$, and
\[
 g_a=\gamma+\log a,\quad d_j=\eta-g_a-j\tau,\quad
 g_i=\gamma+\log(t_i/2),\quad
 v_{i,j}=\beta(t_i)-g_i-j\tau,
 \quad h_i=\frac1{t_i I_0(t_i)^2}.
\]
For the continued embeddings $\Phi^{[j]}=\Phi_{-\pi/4}^{[j]}$, let
$\eta_j(X)=\Phi^{[j]}(\eta)(X)$ and
$\mathscr C_{i,j}(X)=\Phi^{[j]}(C_i)(X)$. Their exact formulas are
\begin{align}
 \eta_j(X)&=g_a+j\tau+d_j e^{a(1-X)},\label{dcr:eta-image}\\
 \mathscr C_{i,j}(X)
 &=C_i-h_i\frac{q_i-1}{p+v_{i,j}(q_i-1)},
 \qquad q_i=e^{2t_i(1-X)}.\label{dcr:C-image}
\end{align}
Put $r_i=2t_i/a$ and $u_i^*=\eta-p/r_i$. As $j\to+\infty$:
\begin{enumerate}
\item On each compact set avoiding $u_1^*,\ldots,u_m^*$, there is
an analytic local parametrization with $\eta_j(X_j(u))=u$ and
\[
 \mathscr C_{i,j}(X_j(u))\longrightarrow C_i\qquad(1\le i\le m)
\]
uniformly on that compact set.
\item Fix $i$. On each compact set avoiding $C_i$, there is an
analytic local parametrization with
$\mathscr C_{i,j}(X_{i,j}(w))=w$ and
\[
 \eta_j(X_{i,j}(w))\longrightarrow u_i^*,\qquad
 \mathscr C_{l,j}(X_{i,j}(w))\longrightarrow C_l\quad(l\ne i)
\]
uniformly on that compact set. Both parametrizations have nonzero
derivative for sufficiently large $j$.
\item There is a simple pole $X_{i,j}^{\rm pole}$ of
$\mathscr C_{i,j}$, tending to $1$, at which all the other displayed
images are holomorphic and satisfy
\[
 \eta_j(X_{i,j}^{\rm pole})\longrightarrow u_i^*,\qquad
 \mathscr C_{l,j}(X_{i,j}^{\rm pole})\longrightarrow C_l
 \quad(l\ne i).
\]
\end{enumerate}
\end{lemma}
\begin{proof}
Formula~\eqref{dcr:eta-image} is the positive-point case of
\eqref{ner:windings}. Formula~\eqref{dcr:C-image} follows by substituting
the coefficients in \eqref{mbc:Iimage}--\eqref{mbc:Jimage} and using
the Bessel Wronskian. In particular, its constants are coefficients of
the phase images, not elements asserted to be fixed by the embedding.

For the first chart, choose the logarithm near $1$ and set
\[
 X_j(u)=1-\frac1a\log\left(1+\frac{u-\eta}{d_j}\right).
\]
Then
\[
 q_i=\left(1+\frac{u-\eta}{d_j}\right)^{r_i},\qquad
 p+v_{i,j}(q_i-1)\longrightarrow p+r_i(u-\eta).
\]
The last expression is nonzero away from $u_i^*$, whereas $q_i-1$
tends uniformly to zero. This proves the first assertion.

For the second chart solve \eqref{dcr:C-image} with value $w$:
\[
 q_i(w)=1+\frac{p(C_i-w)}{h_i-v_{i,j}(C_i-w)},\qquad
 X_{i,j}(w)=1-\frac1{2t_i}\log q_i(w).
\]
Uniformly away from $C_i$,
$q_i(w)-1=-p/v_{i,j}+O(j^{-2})$. Consequently
\[
 \eta_j(X_{i,j}(w))\longrightarrow\eta-\frac{pa}{2t_i}=u_i^*,
 \qquad
 p+v_{l,j}\bigl(q_i(w)^{t_l/t_i}-1\bigr)
       \longrightarrow p(1-t_l/t_i).
\]
The latter limit is nonzero for $l\ne i$, which proves the claimed
limits for the other logarithmic derivatives. The first chart has
derivative $-1/(a(d_j+u-\eta))$. In the second chart,
$q_i'(w)=-ph_i/(h_i-v_{i,j}(C_i-w))^2\ne0$ and $q_i(w)\ne0$
for large $j$. Thus the inverse changes of parameter preserve poles.

Finally take $q_i=1-p/v_{i,j}$ and the logarithm near $1$ to define
$X_{i,j}^{\rm pole}$. The numerator and the derivative of the denominator
in \eqref{dcr:C-image} are nonzero there for sufficiently large $j$,
so this pole is simple. The same limits just computed show that the
other images are holomorphic there and have the asserted limits.
\end{proof}

\begin{proof}[Proof of Theorem~\ref{dcr:main}]
By Theorem~\ref{cf:joint}, the real numbers $\eta,C_1,\ldots,C_m$
are jointly algebraically independent over $k$. Let
$F=k(C_1,\ldots,C_m)$. This field is stable under complex conjugation.
For every $i$, the number $u_i^*$ is transcendental over $F$: otherwise
its conjugate would also be algebraic over $F$, and their average
$\eta$ would be algebraic over $F$, a contradiction. In particular,
$u_i^*$ and any subtuple of the $C_l$ are jointly independent over $k$.

Take a nonzero rational function
\[
 R=\frac{P(U,V_1,\ldots,V_m)}{Q(U,V_1,\ldots,V_m)},\qquad
 P,Q\in k[U,V_1,\ldots,V_m],\quad \gcd(P,Q)=1.
\]
All numerical denominators are nonzero by joint independence. We show
that its value can be integral over $\E$ only when it is a polynomial
in $\eta$.

Suppose first that $D=\deg_U Q>0$. Over the rational function field
$k(V_1,\ldots,V_m)$, the two polynomials $P,Q$ remain coprime as
polynomials in $U$. Their resultant in $U$ is therefore nonzero.
Specializing the $V_i$ to the independent $C_i$ preserves that
nonvanishing. Hence $Q(U,C_1,\ldots,C_m)$ has degree $D$ and shares
no root with $P(U,C_1,\ldots,C_m)$.

Every root $u_0$ of this specialized denominator differs from all the
$u_i^*$: a root is algebraic over $F$, whereas every $u_i^*$ is
transcendental over $F$. Choose a small disk about such a root $u_0$
avoiding every critical point, with no zero of the limiting denominator
on its boundary. In the first chart of Lemma~\ref{dcr:charts}, the
embedded denominator converges uniformly on this disk to
$Q(u,C_1,\ldots,C_m)$. Rouch\'e's theorem supplies a nearby denominator
zero for all large $j$. The numerator is nonzero near $u_0$, by the
resultant argument and uniform convergence. Thus the embedded rational
function has a pole at a finite $X$.

Next suppose that $Q$ is independent of $U$ but is nonconstant. Choose
an index $i$ on which it depends and set $F_i=k(C_l:l\ne i)$. After
specializing the other coordinates, $Q$ is a nonconstant polynomial
$Q_i(V)\in F_i[V]$, and $P$ is a polynomial $P_i(U,V)\in F_i[U,V]$
coprime to $Q_i$. These statements follow by localization in the other
independent variables. Let $\alpha$ be a root of $Q_i$. Then
$\alpha\ne C_i$ because $C_i$ is transcendental over $F_i$.
Moreover,
\[
                       P_i(u_i^*,\alpha)\ne0.
\]
Otherwise transcendence of $u_i^*$ over the algebraic extension
$F_i(\alpha)$ would give the polynomial identity $P_i(U,\alpha)=0$.
The minimal polynomial of $\alpha$ over $F_i$ would divide both $P_i$
and $Q_i$, contradicting their coprimality.
Use the second chart of Lemma~\ref{dcr:charts} on a small disk about
$\alpha$ avoiding $C_i$. The denominator converges to $Q_i(w)$ and
the numerator to $P_i(u_i^*,w)$. Rouch\'e's theorem again gives an
uncancelled denominator zero, hence a finite pole of the phase image.

It remains to treat a polynomial $P$ depending on at least one $V_i$.
Write its positive degree in this variable as $d$, with leading
coefficient $P_d(U,(V_l)_{l\ne i})\ne0$. At the pole in the third
chart of Lemma~\ref{dcr:charts}, this coefficient tends to
\[
                  P_d(u_i^*,(C_l)_{l\ne i})\ne0,
\]
by the independence established at the start of the proof. All the
other coordinates are holomorphic there. Consequently the polynomial
phase image has a pole of order $d$ for sufficiently large $j$.

Every pole contradicts integrality by Lemma~\ref{mbc:pole}.
Thus only polynomials in $\eta$ can belong to the intersection.
Conversely $\eta$ is an $E$-value, so $k[\eta]\subseteq\E$.
This proves \eqref{dcr:intersection}.
\end{proof}

\begin{corollary}[The full logarithmic-derivative field at dependent scales]
\label{dcr:C-only}
For pairwise distinct positive algebraic $t_1,\ldots,t_m$,
\[
 \Ebar\cap k\bigl(C(t_1),\ldots,C(t_m)\bigr)=k.
\]
In particular
\[
 \Ebar\cap k\bigl(C(2),C(4)\bigr)=k,
 \qquad
 \Ebar\cap k\bigl(\Ein(1),C(2),C(4)\bigr)=k[\Ein(1)].
\]
The same exclusion holds for the field generated by all $C(t)$ with
positive algebraic $t$, since every rational expression uses only
finitely many arguments.
\end{corollary}
\begin{proof}
Choose any positive algebraic $a$ and apply the theorem. Independence
of $\Ein(a)$ from the $C(t_i)$ gives
$k[\Ein(a)]\cap k(C(t_1),\ldots,C(t_m))=k$. The infinite-family
assertion follows by restricting any proposed element to its finite
set of arguments.
\end{proof}

\subsection{One harmonic coordinate and arbitrarily many derivative scales}

The same local charts permit one harmonic Bessel coordinate to be
adjoined, without imposing rational independence on any of the scales.

\begin{theorem}[One harmonic coordinate at dependent derivative scales]
\label{dcr:one-beta}
Let $a,s,t_1,\ldots,t_m$ be positive algebraic numbers, with the $t_i$
pairwise distinct. Then
\begin{equation}\label{dcr:one-beta-intersection}
 \boxed{\quad
 \Ebar\cap k\bigl(\Ein(a),\beta(s),C(t_1),\ldots,C(t_m)\bigr)
       =k[\Ein(a)].
 \quad}
\end{equation}
In particular, the equality remains true when all positive algebraic
derivative scales $C(t)$ are adjoined simultaneously.
\end{theorem}
\begin{proof}
Adjoining $C(s)$ if necessary only enlarges the field, so we may
relabel the distinct scales such that $t_1=s$. Abbreviate
$\beta_s=\beta(s)$ and $\eta=\Ein(a)$, and retain the notation of
Lemma~\ref{dcr:charts}. Write
$\mathscr B_j(X)=\Phi^{[j]}(\beta_s)(X)$ and $b_j=v_{1,j}$. The
exact formulas and the line identity \eqref{brp:line} give
\begin{align}
 \mathscr B_j(X)
 &=g_1+j\tau+
       \frac{p b_jq_1}{p+b_j(q_1-1)},\label{dcr:beta-image}\\
 \mathscr C_{1,j}(X)
 &=C_1+\sigma_j\bigl(\mathscr B_j(X)-\beta_s\bigr),
 \qquad
 \sigma_j\sim-\frac1{4\pi^2sI_0(s)^2j^2}.
 \label{dcr:beta-line}
\end{align}
For large $j$ the slope $\sigma_j$ is nonzero.

We shall use one additional local chart. Prescribe a bounded value $v$
of $\mathscr B_j$ and set
\[
 q_1(v)=1+\frac{p(v-\beta_s)}{b_j(p-v+g_1+j\tau)},
 \qquad
 X_j^{\beta}(v)=1-\frac1{2s}\log q_1(v),
\]
with the logarithm near $1$. These formulas give
$\mathscr B_j(X_j^{\beta}(v))=v$ exactly. Since $q_1(v)-1=O(j^{-2})$,
uniformly on bounded sets,
\begin{equation}\label{dcr:beta-chart}
 \eta_j(X_j^{\beta}(v))\longrightarrow\eta,
 \qquad
 \mathscr C_{l,j}(X_j^{\beta}(v))\longrightarrow C_l
 \quad(1\le l\le m).
\end{equation}
The change of parameter has nonzero derivative for all sufficiently
large $j$, as follows either from its displayed formula or from the
nonconstant M\"obius map \eqref{dcr:beta-image}.

Two growth estimates on the earlier charts will also be needed.
On the first chart of Lemma~\ref{dcr:charts}, uniformly away from
the critical points,
\begin{equation}\label{dcr:beta-growth-eta}
 \frac{\mathscr B_j(X_j(u))}{j}
   \longrightarrow
   \frac{\tau r_1(u-\eta)}{p+r_1(u-\eta)}.
\end{equation}
On the chart with $\mathscr C_{i,j}=w$, uniformly away from $C_i$,
\begin{align}
 \frac{\mathscr B_j(X_{i,j}(w))}{j}
 &\longrightarrow \frac{\tau(s/t_i)}{s/t_i-1}
       &&(i\ne1),\label{dcr:beta-growth-C-other}\\
 \frac{\mathscr B_j(X_{1,j}(w))}{j^2}
 &\longrightarrow-4\pi^2sI_0(s)^2(w-C_1)
       &&(i=1).\label{dcr:beta-growth-C-own}
\end{align}
To verify \eqref{dcr:beta-growth-eta}, insert
$q_1-1=r_1(u-\eta)/d_j+O(j^{-2})$ into
\eqref{dcr:beta-image}. For \eqref{dcr:beta-growth-C-other}, use
$q_i-1=-p/v_{i,j}+O(j^{-2})$ and $q_1=q_i^{s/t_i}$.
The denominator in \eqref{dcr:beta-image} then tends to
$p(1-s/t_i)$. Formula~\eqref{dcr:beta-growth-C-own} follows directly
from \eqref{dcr:beta-line}.

The numbers $\eta,\beta_s,C_1,\ldots,C_m$ are jointly independent
by Theorem~\ref{cf:joint}. Let $P/Q$ be a reduced nonzero rational
function in the formal variables $U,B,V_1,\ldots,V_m$, with coefficients
in $k$, evaluated at that tuple.

If $\deg_B Q>0$, coprimality over $k(U,V_1,\ldots,V_m)$ gives a
nonzero resultant in $B$. Independence of $\eta,C_1,\ldots,C_m$
preserves it after specialization. Thus the nonconstant polynomial
$Q(\eta,B,C_1,\ldots,C_m)$ has a root where the specialized numerator
is nonzero. The chart \eqref{dcr:beta-chart} and Rouch\'e's theorem
give a nearby uncancelled denominator zero of the phase image.
This contradicts integrality by Lemma~\ref{mbc:pole}.

We may therefore suppose that $Q$ is independent of $B$. If it is
nonconstant, choose an irreducible factor
$q(U,V_1,\ldots,V_m)$ of $Q$. Write
$P=\sum_b p_b(U,V_1,\ldots,V_m)B^b$, and let $d$ be the largest
index for which $q$ does not divide $p_d$. Such an index exists because
$P,Q$ are coprime. At an exact zero of $q$ all terms with $b>d$ vanish.

If $\deg_U q>0$, its resultant in $U$ with $p_d$ is nonzero over
$k(V_1,\ldots,V_m)$. After specializing the independent $C_l$, choose
a root $u_0$ of $q(U,C_1,\ldots,C_m)$. Then
$p_d(u_0,C_1,\ldots,C_m)\ne0$. Furthermore $u_0$ differs from
$\eta$ and from every $u_i^*$, since it is algebraic over
$F=k(C_1,\ldots,C_m)$ whereas these numbers are transcendental over
$F$. In the first chart of Lemma~\ref{dcr:charts}, Rouch\'e's theorem
gives zeros $u_j\to u_0$ of the embedded factor $q$. By
\eqref{dcr:beta-growth-eta}, the corresponding $\beta_s$ images grow
linearly in $j$ with a nonzero limiting coefficient. At these zeros
the terms of $P$ above degree $d$ vanish exactly, its degree-$d$
coefficient tends to a nonzero number, and all lower coefficients
remain bounded. Hence the numerator is nonzero for large $j$.

If $q$ is independent of $U$, choose a variable $V_i$ on which it
depends and put $F_i=k(C_l:l\ne i)$. Over $F_i$, the specialization
$q_i(V)$ is a nonconstant polynomial coprime to the specialized
polynomial $p_d(U,V)$. Let $\alpha$ be a root of $q_i$. As in the
proof of Theorem~\ref{dcr:main}, one has $\alpha\ne C_i$ and
$p_d(u_i^*,\alpha,(C_l)_{l\ne i})\ne0$. Use the chart with
$\mathscr C_{i,j}=w$ near $\alpha$. Rouch\'e's theorem gives nearby
zeros of the embedded $q$, while
\eqref{dcr:beta-growth-C-other}--\eqref{dcr:beta-growth-C-own} show
that the $\beta_s$ image tends to infinity, with order $j$ or $j^2$
and nonzero leading coefficient. The same argument with the highest
coefficient not divisible by $q$ shows that the numerator is nonzero there.
In either case an uncancelled pole contradicts integrality. Thus $Q$
must be constant.

Finally let $P$ be a polynomial of positive total degree $D$ in the
two variables $B,V_1$. Write its top homogeneous part in these two
variables as
\[
 \sum_{b+n=D}p_{b,n}(U,V_2,\ldots,V_m)B^bV_1^n,
 \qquad
 H_D(U,V_2,\ldots,V_m,S)
       =\sum_{b+n=D}p_{b,n}(U,V_2,\ldots,V_m)S^n.
\]
At the pole $X_{1,j}^{\rm pole}$ of Lemma~\ref{dcr:charts}, both
$\mathscr B_j$ and $\mathscr C_{1,j}$ have simple poles, related by
\eqref{dcr:beta-line}. The coefficient of the possible pole of order
$D$, measured in powers of $\mathscr B_j$, is
\[
 H_D\bigl(\eta_j(X_{1,j}^{\rm pole}),
       (\mathscr C_{l,j}(X_{1,j}^{\rm pole}))_{l\ge2},\sigma_j\bigr).
\]
Divide this expression by the lowest power of $\sigma_j$ occurring
in $H_D$. The resulting expression tends to a nonzero coefficient
$p_{b,n}(u_1^*,C_2,\ldots,C_m)$, by the independence proved in
Theorem~\ref{dcr:main}. Consequently the pole of order $D$ survives
for all large $j$. If instead $P$ is independent of both $B,V_1$,
Theorem~\ref{dcr:main} already shows that it can be integral only
when it is a polynomial in $\eta$.

The reverse inclusion holds because $\eta$ is an $E$-value. This
proves \eqref{dcr:one-beta-intersection}; the infinite-family assertion
follows because each rational expression involves finitely many scales.
\end{proof}

\begin{corollary}[A global derivative field with one harmonic coordinate]
\label{dcr:one-beta-global}
For every fixed positive algebraic $s$,
\[
 \Ebar\cap
 k\bigl(\beta(s),C(t):t>0\text{ algebraic}\bigr)=k.
\]
\end{corollary}
\begin{proof}
Adjoin $\Ein(a)$ for any positive algebraic $a$ and apply
Theorem~\ref{dcr:one-beta}. Joint independence makes $\Ein(a)$
transcendental over the displayed field: any algebraic relation would
involve only finitely many of its generators. Its polynomial ring
therefore intersects that field in $k$.
\end{proof}

\begin{remark}[The remaining dependent-scale question]
For example, the theorem gives
\[
 \Ebar\cap k\bigl(\Ein(1),\beta(2),C(2),C(4)\bigr)=k[\Ein(1)].
\]
It does not classify the field with two harmonic coordinates
$k(\beta(2),C(2),\beta(4),C(4))$. The unresolved part of the
dependent-scale question therefore concerns simultaneous harmonic
coordinates, rather than arbitrary derivative coordinates or one
harmonic coordinate accompanied by them.
\end{remark}

% Combined proof-bearing extension, 22 October 2026.
% Owned by resonant_closure.  Requires fre:main and dcr:one-beta.
% No authoritative manuscript or other fragment is modified by this file.

\section{A global derivative field with an exponential and one harmonic value}
\label{gde:section}

The exponential extension and the dependent derivative-scale theorem
can be combined.  The result gives one exact intersection for the
entire family of positive algebraic logarithmic-derivative arguments,
while retaining one harmonic Bessel coordinate.

\begin{theorem}[Exponential--Ein intersection over a derivative family]
\label{gde:main}
Let $a,s,t_1,\ldots,t_m$ be positive algebraic numbers, with the $t_i$
pairwise distinct.  Put
\[
 k=\overline{\mathbb Q},\qquad E=e^a,\qquad
 \eta=\Ein(a),\qquad \beta_s=\beta(s),\qquad C_i=C(t_i).
\]
Then
\begin{equation}\label{gde:intersection}
 \boxed{\quad
 \overline{\mathbf E}\cap
 k(E,\eta,\beta_s,C_1,\ldots,C_m)
       =k[E,E^{-1},\eta].
 \quad}
\end{equation}
No rational-independence hypothesis on the arguments is required.
\end{theorem}

\begin{proof}
Adding $C(s)$ enlarges the field, so it suffices to prove the result
after doing so and relabelling the distinct arguments such that
$t_1=s$.  Retain the notation and synchronized winding charts of
Lemma~\ref{dcr:charts}, with
\[
 p=i\pi,\qquad \tau=2\pi i,\qquad
 r_i=\frac{2t_i}{a},\qquad u_i^*=\eta-\frac{p}{r_i}.
\]
Write $\chi_j=\Phi^{[j]}(E)$, $\eta_j=\Phi^{[j]}(\eta)$,
$\mathscr B_j=\Phi^{[j]}(\beta_s)$, and
$\mathscr C_{i,j}=\Phi^{[j]}(C_i)$.

\paragraph{Independence and local limits.}
The joint theorem over the field generated by algebraic exponentials,
and the transcendence of $E$, give joint independence over $k$ of
\begin{equation}\label{gde:independence}
 E,\eta,\beta_s,C_1,\ldots,C_m.
\end{equation}
Let $F=k(E,C_1,\ldots,C_m)$.  All its displayed generators are real,
so $F$ is stable under conjugation.  Each $u_i^*$ is transcendental
over $F$: algebraicity of $u_i^*$ would imply algebraicity of its
conjugate, and then of their average $\eta$, contrary to
\eqref{gde:independence}.  In particular, for every subtuple of the
$C_l$, the numbers $E,u_i^*$ and that subtuple are jointly
independent over $k$.

The exponential image is always
\[
                    \chi_j(X)=e^{aX}.
\]
All the local charts below have $X\to1$ uniformly on their fixed
closed disks.  Consequently they all have $\chi_j\to E$ uniformly.
Together with the previously proved chart formulas, this gives:
\begin{enumerate}
\item In the chart $\mathscr B_j=v$, with $v$ bounded,
\[
 (\chi_j,\eta_j,\mathscr C_{1,j},\ldots,\mathscr C_{m,j})
       \longrightarrow(E,\eta,C_1,\ldots,C_m).
\]
\item In the chart $\eta_j=u$, on disks avoiding every $u_i^*$,
\[
 \chi_j\longrightarrow E,\qquad
 \mathscr C_{i,j}\longrightarrow C_i,\qquad
 \frac{\mathscr B_j}{j}\longrightarrow
       \frac{\tau r_1(u-\eta)}{p+r_1(u-\eta)}.
\]
If the disk also avoids $\eta$, the last limiting function is
bounded away from zero.
\item In the chart $\mathscr C_{i,j}=w$, on disks avoiding $C_i$,
\[
 \chi_j\longrightarrow E,\qquad \eta_j\longrightarrow u_i^*,
 \qquad \mathscr C_{l,j}\longrightarrow C_l\quad(l\ne i).
\]
For $i\ne1$ the additional limit is
\[
 \frac{\mathscr B_j}{j}\longrightarrow
              \frac{\tau(s/t_i)}{s/t_i-1}\ne0.
\]
For $i=1$ it is
\[
 \frac{\mathscr B_j}{j^2}\longrightarrow
              -4\pi^2sI_0(s)^2(w-C_1),
\]
which is bounded away from zero on the chosen disk.
\end{enumerate}
These are \eqref{dcr:beta-chart} and
\eqref{dcr:beta-growth-eta}--\eqref{dcr:beta-growth-C-own},
with the additional exponential limit.
All local inverse derivatives are nonzero for sufficiently large
$j$, as proved with those charts, so poles in the target parameters
are poles at finite values of $X$.

\paragraph{Reduction of a rational denominator.}
Take a reduced nonzero rational expression
\[
 \xi=\frac{P(E,\eta,\beta_s,C_1,\ldots,C_m)}
          {Q(E,\eta,\beta_s,C_1,\ldots,C_m)}
       \in\overline{\mathbf E},
 \qquad
 P,Q\in k[S,U,B,V_1,\ldots,V_m].
\]
All nonzero denominators at the actual tuple are nonzero by
\eqref{gde:independence}.  Continued phase embeddings therefore
act on the field in question.  Their images of $\xi$ cannot have
finite poles.

If $Q$ has positive $B$-degree, its resultant in $B$ with $P$ is
nonzero over $k(S,U,V_1,\ldots,V_m)$, because the two polynomials
are coprime.  Clearing denominators gives a nonzero polynomial in
the remaining variables.  Independence of $E,\eta,C_1,\ldots,C_m$
preserves this resultant and the leading coefficient after
specialization.  Thus the nonconstant polynomial
$Q(E,\eta,B,C_1,\ldots,C_m)$ has a root where the specialized
numerator is nonzero.  The finite-$\mathscr B_j$ chart and
Rouch\'e's theorem produce an actual denominator zero with
nonzero numerator for every sufficiently large $j$.
This is a forbidden pole.  Hence $Q$ is independent of $B$.

Choose an irreducible factor $q(S,U,V_1,\ldots,V_m)$ of $Q$.
Expand
\[
 P=\sum_b p_b(S,U,V_1,\ldots,V_m)B^b,
\]
and let $d$ be the largest index such that $q\nmid p_d$.
It exists by coprimality.  At every exact zero of the embedded
factor $q$, all terms with $b>d$ vanish exactly.

Suppose first that $\deg_U q>0$.  By Gauss's lemma, $q$ and $p_d$
are coprime over $k(S,V_1,\ldots,V_m)$ as polynomials in $U$.
Their resultant remains nonzero after the independent
specialization $(S,V_1,\ldots,V_m)=(E,C_1,\ldots,C_m)$.
The specialized $q(E,U,C_1,\ldots,C_m)$ is nonconstant.
Choose a root $u_0$; then
\[
 p_d(E,u_0,C_1,\ldots,C_m)\ne0.
\]
The root $u_0$ is algebraic over $F$, whereas $\eta$ and all
$u_i^*$ are transcendental over $F$.  Thus it differs from those
finitely many exceptional values.  Choose a small closed disk
around $u_0$ avoiding them, on which the displayed numerator
coefficient stays nonzero and whose boundary contains no root of
the limiting denominator.

On the Ein chart, uniform convergence and Rouch\'e's theorem
give exact zeros of the embedded $q$ in this disk.  At these zeros
the terms above degree $d$ vanish exactly, the degree-$d$
coefficient stays bounded away from zero, all lower coefficients
are bounded, and $|\mathscr B_j|\asymp j$.  The full numerator is
therefore nonzero for large $j$.  This produces a forbidden pole.
Hence no irreducible factor of $Q$ can have positive $U$-degree,
and $Q\in k[S,V_1,\ldots,V_m]$.

Suppose next that an irreducible factor $q$ of $Q$ has positive
degree in some $V_i$.  Put
\[
 F_i=k(E,C_l:l\ne i).
\]
After localizing in the other independent variables and specializing
them, $q$ becomes an irreducible nonconstant polynomial $q_i(V)$
over $F_i$, up to multiplication by a nonzero constant.
The coefficient $p_d$ becomes a polynomial $p_{d,i}(U,V)$ not
divisible by $q_i$.  Choose a root $\alpha$ of $q_i$.  Then
$\alpha\ne C_i$, since $C_i$ is transcendental over $F_i$.
Moreover
\begin{equation}\label{gde:root-numerator}
                         p_{d,i}(u_i^*,\alpha)\ne0.
\end{equation}
Indeed, $u_i^*$ is transcendental over the algebraic extension
$F_i(\alpha)$; vanishing in \eqref{gde:root-numerator} would
therefore imply the polynomial identity $p_{d,i}(U,\alpha)=0$.
The minimal polynomial $q_i$ would then divide every coefficient
in $U$, a contradiction.

Use the $\mathscr C_{i,j}$ chart in a small closed disk about
$\alpha$ avoiding $C_i$.  The denominator factor converges to
$q_i(w)$, and its selected numerator coefficient converges to
$p_{d,i}(u_i^*,w)$, which is nonzero near $\alpha$.
Rouch\'e's theorem gives exact denominator-factor zeros in this
disk.  At those zeros the higher $B$-coefficients vanish, and
$\mathscr B_j$ diverges at order $j$ if $i\ne1$, or at order $j^2$
if $i=1$, with a nonzero limiting coefficient.
The highest surviving coefficient therefore makes the numerator
nonzero.  This is another forbidden pole.  We conclude that
$Q\in k[S]$.

\paragraph{Only powers of the exponential may remain in the denominator.}
If $\alpha\in k^\times$ were a root of $Q$, coprimality would give
$P(\alpha,U,B,V_1,\ldots,V_m)\ne0$.
Choose a finite $X_0$ with $e^{aX_0}=\alpha$.
For every $l$,
\[
 \left|e^{2t_l(1-X_0)}\right|
       =\left(\frac{E}{|\alpha|}\right)^{2t_l/a}\ne1.
\]
Here $|\alpha|$ is algebraic and $E$ is transcendental.  Thus none
of the corresponding phase coordinates equals $1$.
The exact winding formulas at this fixed point give
\[
 \eta_j(X_0)=\tau j(1-E/\alpha)+O(1),\qquad
 \mathscr B_j(X_0)=\tau j+O(1),\qquad
 \mathscr C_{l,j}(X_0)\longrightarrow C_l.
\]
All Bessel images are finite at $X_0$ for large $j$.
Let $D$ be the total degree of
$P(\alpha,U,B,V_1,\ldots,V_m)$ in $(U,B)$, and let $H_D$ be its
nonzero top homogeneous part in those variables.
Independence of $E,C_1,\ldots,C_m$ implies
\[
 H_D(1-E/\alpha,1,C_1,\ldots,C_m)\ne0.
\]
Indeed $H_D(T,1,V_1,\ldots,V_m)$ is a nonzero polynomial.
The embedded numerator, divided by $(\tau j)^D$, tends to this
nonzero number.  Its denominator vanishes at $X_0$, and the
exponential image has derivative $a\alpha\ne0$ there.
This is an uncancelled finite pole.  Therefore $Q$ has no nonzero
root, and $Q$ is a constant times a power of $S$.
We have reduced $\xi$ to
\[
                k[E,E^{-1},\eta,\beta_s,C_1,\ldots,C_m].
\]

\paragraph{Eliminating polynomial dependence on the Bessel coordinates.}
Suppose this Laurent polynomial has positive total degree $D$
in the two coordinates $\beta_s,C_1$.  Its top homogeneous part
in formal variables $B,V_1$ has coefficients in
$k[S,S^{-1},U,V_2,\ldots,V_m]$.
At the common pole $X_{1,j}^{\rm pole}$ of the two images,
all other images are holomorphic and
\[
 \chi_j\longrightarrow E,\qquad
 \eta_j\longrightarrow u_1^*,\qquad
 \mathscr C_{l,j}\longrightarrow C_l\quad(l\ge2).
\]
The exact line identity is
\[
 \mathscr C_{1,j}
   =C_1+\sigma_j(\mathscr B_j-\beta_s),
 \qquad \sigma_j\ne0,\quad \sigma_j\longrightarrow0.
\]
In the coefficient of the possible order-$D$ pole, divide by the
lowest power of $\sigma_j$ occurring in the top homogeneous part.
The quotient tends to a nonzero Laurent polynomial evaluated at
$E,u_1^*,C_2,\ldots,C_m$.  It is nonzero by the independence
established above.  Thus an order-$D$ pole survives, a contradiction.
The Laurent polynomial must be independent of both $\beta_s,C_1$.

If it still depends on some $C_i$ with $i\ge2$, use the pole
$X_{i,j}^{\rm pole}$ from Lemma~\ref{dcr:charts}.
The leading coefficient in this coordinate tends to a nonzero
Laurent polynomial in $E,u_i^*$ and the remaining actual derivative
values; these are jointly independent.  The other coordinates
are holomorphic at the point, and the exponential image is
nonzero.  Hence the positive degree in $C_i$ again creates an
uncancelled pole.  This excludes every remaining derivative
coordinate.

It follows that $\xi\in k[E,E^{-1},\eta]$.  The reverse inclusion
holds because $e^a,e^{-a}$ and $\Ein(a)$ are $E$-values.
This proves \eqref{gde:intersection}.
\end{proof}

\begin{corollary}[A single exact ring over all positive derivative scales]
\label{gde:global}
For every fixed pair of positive algebraic numbers $a,s$,
\[
 \overline{\mathbf E}\cap
 k\bigl(e^a,\Ein(a),\beta(s),C(t):
                 t>0\text{ algebraic}\bigr)
       =k[e^a,e^{-a},\Ein(a)].
\]
Moreover
\[
 \overline{\mathbf E}\cap
 k\bigl(e^a,\beta(s),C(t):t>0\text{ algebraic}\bigr)
       =k[e^a,e^{-a}].
\]
\end{corollary}
\begin{proof}
Every rational expression involves only finitely many derivative
arguments, so the first identity follows from the theorem.
For the second, $\Ein(a)$ is transcendental over the indicated
field by the finite-subset form of joint independence.
Thus a polynomial in $\Ein(a)$ with coefficients in
$k[e^a,e^{-a}]$ can belong to that field only if it is constant
in $\Ein(a)$.
\end{proof}

\begin{remark}[The remaining harmonic-coordinate restriction]
The theorem allows all positive algebraic derivative scales
simultaneously and arbitrary rational expressions in the displayed
generators.  It retains only one harmonic coordinate $\beta(s)$.
It does not settle the field obtained by adjoining two different
harmonic coordinates such as $\beta(2)$ and $\beta(4)$.
The result is an internally audited consequence of the continued
phase embeddings and the established joint independence theorem;
external expert verification and historical priority remain separate.
\end{remark}

\subsection{The polynomial consequence at a single positive scale}

The full rational classification contains the earlier polynomial result.
In the notation of the matched-scale section, $\eta_a=\Ein(a)$,
$B_a=\beta(2a)$, and $C_a=C(2a)$.

\begin{corollary}[One-scale polynomial classification]
\label{mre:polynomial}\label{etaclass:main}
For every positive algebraic number $a$,
\begin{equation}\label{mre:polynomial-equation}
 \Ebar\cap\Qa[\eta_a,B_a,C_a]=\Qa[\eta_a].
\end{equation}
\end{corollary}
\begin{proof}
Apply Theorem~\ref{fri:main} with the Ein argument $a$ and Bessel
argument $2a$, and restrict its rational field to the displayed
polynomial ring. The reverse inclusion follows from $\eta_a\in\E$.
\end{proof}

\section{Integral classifications in the three proposed coordinates}

The general positive-scale classification specializes at $a=1$ to
$\eta=\Ein(1)$, $\beta=\gamma+\kappa$, and $C=I_1(2)/I_0(2)$.
It yields an exact statement in the original proposed coordinates.

\begin{corollary}[Exact classification in the proposed coordinates]
\label{etaclass:coordinates}
Let $\ell=\gamma-2\kappa$ and
\[
 A=\delta/e-\kappa=\eta-\beta,\qquad
 U=\ell+3\delta/e=3\eta-2\beta,\qquad
 V=\ell+3\kappa=\beta.
\]
Then
\[
 U=V+3A,\qquad
 \Ebar\cap\Qa(A,U,V,C)=\Qa[U+2V]
 =\Qa[A+V].
\]
Equivalently, a rational expression in $A,U,V,C$ is integral over the
E-value ring exactly when, after imposing $U-V-3A=0$, it equals
a polynomial in $U+2V=3\eta$. Every nonconstant rational function
in $A$ alone, $U$ alone, or $V$ alone is non-integral.
\end{corollary}
\begin{proof}
The invertible change $(\eta,\beta)\leftrightarrow(A,V)$ gives
$U=3A+V$ and $A+V=\eta$.  Apply
Theorem~\ref{fri:main} at $(a,t)=(1,2)$; multiplication of $\eta$
by the nonzero algebraic scalar $3$ does not change its polynomial ring.
\end{proof}


\begin{remark}
The coefficient ring in the theorem is $\Qa[\eta]$, not the
whole E-value ring: allowing arbitrary E-value coefficients
would make $I_0(2)\beta=Y(2)\in\E$ an immediate counterexample.
The rational-field classification therefore keeps its specified
algebraic coefficient field.
\end{remark}

\begin{theorem}[All nontrivial Laurent monomials are excluded]
\label{etaclass:laurent}
For integers $a,u,v$, not all zero,
\[
 A^aU^uV^v\notin\Ebar.
\]
In particular $A,U,V$ are multiplicatively independent, and every
nonzero positive or negative integer power of any one of them is
non-integral over the E-value ring.
\end{theorem}
\begin{proof}
Independence of $\eta,\beta$ identifies their polynomial ring with
$k[X,Y]$. Its three linear factors $X-Y,3X-2Y,Y$ are pairwise
nonassociate irreducibles. The valuations of
$(X-Y)^a(3X-2Y)^uY^v$ at these factors are exactly $a,u,v$.
Every nonzero polynomial in $X$ has valuation zero at all three
factors. Thus, unless all exponents vanish, the monomial does not
belong to $k[\eta]$. Theorem~\ref{fri:main} excludes its integrality.
The same observation excludes equality to $1$ and proves
multiplicative independence.
\end{proof}




\begin{corollary}[Rational-function exclusion for mixed linear forms]
\label{etaclass:rational}
Let $k_0$ be any subfield of $\Qa$, let $a,b\in k_0^\times$, and put
$Z=a\eta+b\beta$. Then
\[
 \Ebar\cap k_0(Z)=k_0.
\]
In particular,
\[
 \Ebar\cap\Qa(A)=\Qa,\qquad
 \Ebar\cap\Qa(U)=\Qa,\qquad
 \Ebar\cap\Qa(V,C)=\Qa.
\]
\end{corollary}
\begin{proof}
The first assertion follows from Theorem~\ref{mre:main} at the
single scale $1$, together with the fact that the algebraic
constants of the purely transcendental field $k_0(Z)$ are precisely
$k_0$.  Take $(a,b)=(1,-1)$ and $(3,-2)$ for $A$ and $U$.
The last equality, which closes the previously missing field for
$V=\beta(2)$, is Theorem~\ref{brp:main} at $t=2$.
\end{proof}

% Revision insert: exact units, radicals, and finite real-adjunction cost.
% Depends on ncr:main, fri:main, anchor:field, and depth:values.

\section{Units, radicals, and algebraic extensions of the exact intersections}
\label{unit:section}

The exact intersections determine more than membership in the integral
closure. They also determine units and pure radicals inside the specified
rational fields. Throughout this section $k=\overline{\mathbb Q}$,
$\E$ is the ring of $E$-values, and $\Ebar$ is its integral closure
inside $\mathbb C$. The notation $R^\times$ denotes the group of units
of a ring $R$.

\begin{proposition}[Units in the mixed rational field]\label{unit:mixed}
Use the hypotheses and notation of Corollary~\ref{ncr:main}, and set
\[
 A=k[E_1^{\pm1},\ldots,E_m^{\pm1},\eta_1,\ldots,\eta_m,
                                      Z_1^{\pm1},\ldots,Z_r^{\pm1}].
\]
Then
\begin{equation}\label{unit:mixed-equality}
 \E^\times\cap\mathcal F
 =\Ebar^\times\cap\mathcal F
 =A^\times
 =\left\{\alpha\exp\!\left(\sum_{i=1}^m n_i a_i+
                         \sum_{l=1}^r m_l b_l\right):
       \alpha\in k^\times,\ n_i,m_l\in\mathbb Z\right\}.
\end{equation}
For example, for $\mathcal F=k(e,\Ein(1),\beta(\sqrt2),C(\sqrt2))$,
both intersections of unit groups are precisely
$\{\alpha e^n:\alpha\in k^\times,\ n\in\mathbb Z\}$.
\end{proposition}
\begin{proof}
If $u\in\Ebar^\times\cap\mathcal F$, both $u$ and $u^{-1}$
belong to $\Ebar\cap\mathcal F=A$. Thus $u\in A^\times$.
The generators of $A$ are algebraically independent with the indicated
Laurent variables, so its units are exactly scalar Laurent monomials in
the $E_i,Z_l$. Each such monomial and its inverse are $E$-values.
This proves all the equalities. The example satisfies the required phase
independence because $1$ and $2\sqrt2$ are linearly independent over
$\mathbb Q$.
\end{proof}

Fischler--Rivoal conjecture that every unit of $\E$ has the form
$\alpha e^\beta$ with $\alpha\in k^\times$ and $\beta\in k$
\cite[Conjecture~1.2]{FR26}. Proposition~\ref{unit:mixed} verifies
the predicted description after intersection with the particular fields
specified there. It does not classify all units of $\E$.

\begin{proposition}[Pure radicals and monic polynomials]
\label{unit:radicals}
Let $F\subseteq\mathbb C$ be a field containing $k$, and suppose
that an exact intersection $\Ebar\cap F=A$ is known. For $y\in\mathbb C$
and $d\ge1$ with $f=y^d\in F$, one has
\begin{align}
 y\in\Ebar&\quad\Longleftrightarrow\quad f\in A,
                                      \label{unit:radical-integral}\\
 y\in\Ebar^\times&\quad\Longleftrightarrow\quad f\in A^\times
                    \qquad(f\ne0).\label{unit:radical-unit}
\end{align}
More generally, let $P(T)\in F[T]$ be monic of positive degree.
All its roots belong to $\Ebar$ if and only if all its coefficients
belong to $A$. All its roots belong to $\Ebar^\times$ if and only
if its coefficients belong to $A$ and its constant coefficient belongs
to $A^\times$.
\end{proposition}
\begin{proof}
If $y$ is integral over $\E$, so is $y^d$, giving the forward
implication in \eqref{unit:radical-integral}. Conversely, $f\in A$
is integral over $\E$, and $y$ satisfies the monic equation
$T^d-f=0$ over $\Ebar$. Transitivity of integrality makes $y$
integral over $\E$. Applying the same argument to $y^{-1}$ and
$f^{-1}$ proves \eqref{unit:radical-unit}.

If all roots of $P$ are integral, their elementary symmetric functions
are integral and lie in $F$, hence in $A$. Conversely, coefficients
in $A$ make every root integral over $\Ebar$, and then over $\E$.
If all roots are units, their product makes the constant coefficient a
unit of $A$. For the reverse implication, the reciprocal roots satisfy
the reversed polynomial divided by its constant coefficient; this is
again monic with coefficients in $A$. Thus the roots and their inverses
are integral, proving the unit assertion.
\end{proof}

For clarity, define the set
\[
 \operatorname{Rad}(F)=
 \{y\in\mathbb C:y^d\in F\text{ for some integer }d\ge1\}.
\]
It is not being called a field. For the mixed field $\mathcal F$
of Proposition~\ref{unit:mixed}, the radical-unit conclusion becomes
\begin{equation}\label{unit:rational-exponents}
 \E^\times\cap\operatorname{Rad}(\mathcal F)
 =\Ebar^\times\cap\operatorname{Rad}(\mathcal F)
 =\left\{\alpha e^\lambda:
       \alpha\in k^\times,\quad
       \lambda\in\operatorname{span}_{\mathbb Q}
                    \{a_1,\ldots,a_m,b_1,\ldots,b_r\}\right\}.
\end{equation}
Indeed, a root of a unit in \eqref{unit:mixed-equality} is an
algebraic multiple of the exponential with the exponent divided by its
root degree, and every such number is an $E$-unit. The converse follows
by clearing the rational denominators in the exponent.

\begin{corollary}[Units and radicals at resonant positive scales]
\label{unit:resonant}
For positive algebraic $a,t$, put
$F=k(\Ein(a),\beta(t),C(t))$. Then
\[
 \E^\times\cap F=\Ebar^\times\cap F=k^\times,
 \qquad
 \Ebar^\times\cap\operatorname{Rad}(F)=k^\times.
\]
If $y^d=f\in F$, then $y\in\Ebar$ if and only if
$f\in k[\Ein(a)]$. If instead an exact intersection
$\Ebar\cap F=k$ is known, then
\[
                  \Ebar\cap\operatorname{Rad}(F)=k.
\]
\end{corollary}
\begin{proof}
Theorem~\ref{fri:main} gives $A=k[\Ein(a)]$, whose units are
$k^\times$. Apply Proposition~\ref{unit:radicals}; a complex root
of a nonzero algebraic number is again in $k^\times$. In the last
assertion the same argument applies to $y^d\in k$, including zero.
\end{proof}

\begin{proposition}[Integral units which are not E-values]
\label{unit:quadratic-boundary}
Let $\eta=\Ein(1)$ and let $y_+,y_-$ be the two roots of
\[
                         T^2-\eta T+1=0.
\]
Then both roots are transcendental, both belong to $\Ebar^\times$,
and neither belongs to $\E$. Explicitly,
\[
                    y_\pm=\frac{\eta\pm\sqrt{\eta^2-4}}2.
\]
\end{proposition}
\begin{proof}
The monic equation proves integrality, and $y_+y_-=1$ proves that
both are units in $\Ebar$. If a root were algebraic, then
$\eta=y+y^{-1}$ would be algebraic, contrary to the joint theorem.

Suppose a root $y$ belonged to $\E$. The other root is $\eta-y$,
so $y^{-1}\in\E$ as well. Its image under the zero-winding phase
embedding is therefore a unit of
$\mathcal R=\mathbb C[k_{\rm add}]$. The units of this group
algebra are precisely $c e^{\rho X}$, with $c\in\mathbb C^\times$
and $\rho\in k$: the supports of a unit and its inverse generate
a finite-rank free subgroup, and units of its Laurent polynomial ring
are scalar monomials.

On the positive real branch the exponential-integral decomposition gives
\[
                  \Phi_0(\eta)=\gamma+\delta e^{-X}.
\]
Consequently $y+y^{-1}=\eta$ would imply
\[
                c e^{\rho X}+c^{-1}e^{-\rho X}
                      =\gamma+\delta e^{-X}.
\]
The two coefficients on the right are nonzero, since $\gamma,\delta>0$.
For $\rho=0$ the left side is constant; for $\rho\ne0$ its phase
support is $\{\rho,-\rho\}$, which cannot equal $\{0,-1\}$.
Linear independence of distinct exponential functions gives a
contradiction in either case.
\end{proof}

This example also specifies the limit of
\eqref{unit:rational-exponents}: its equality of the two unit groups
is a statement about pure radicals of the prescribed rational field.
It does not extend to an arbitrary algebraic extension of that field.

\begin{corollary}[Units with the matching exponential]
\label{unit:matching-exponential}
Let $a,t$ be arbitrary positive algebraic numbers and put
\[
 F=k(e^a,\Ein(a),\beta(t),C(t)),\qquad k=\overline{\mathbb Q}.
\]
Then
\begin{align*}
 \E^\times\cap F=\Ebar^\times\cap F
   &=\{\alpha e^{na}:\alpha\in k^\times,\ n\in\mathbb Z\},\\
 \E^\times\cap\operatorname{Rad}(F)
   =\Ebar^\times\cap\operatorname{Rad}(F)
   &=\{\alpha e^{qa}:\alpha\in k^\times,\ q\in\mathbb Q\}.
\end{align*}
These conclusions require no independence condition on $a$ and $t$.
In particular they apply to $a=1$, $t=2$.
\end{corollary}
\begin{proof}
Theorem~\ref{fre:main} gives
$\Ebar\cap F=A=k[e^a,e^{-a},\Ein(a)]$.
The joint independence theorem makes $A$ a Laurent polynomial ring
in $e^a$ and a polynomial variable $\Ein(a)$, so
$A^\times=k^\times e^{a\mathbb Z}$.
A unit and its inverse lie in $A$, which proves the first equality;
the displayed monomials and their inverses are $E$-values.
If $y\in\operatorname{Rad}(F)$ is an $\Ebar$-unit, choose $d\ge1$
with $y^d\in F$. Then $y^d=\alpha e^{na}$ by the first part, whence
$y=\alpha' e^{na/d}$ with $\alpha'\in k^\times$.
Conversely, clearing the denominator of $q$ shows that every
$\alpha e^{qa}$ belongs to $\operatorname{Rad}(F)$, and it and its
inverse are $E$-values.
\end{proof}


\part{Critical asymptotics}
\section{Critical factorial and Newton--Laguerre limits}\label{sec:limits}
\subsection{The factorial limit and its first correction}
For integers $n,m\ge1$, define the rational numbers
\begin{align}
f_k&=\frac{(m+k)!}{k!^2(n-k)!^2},&
h_k&=H_{m+k}-2H_k+2H_{n-k},\nonumber\\
r_{n,m}&=\frac{\sum_{k=0}^nf_kh_k}{\sum_{k=0}^nf_k}.\label{eq:r}
\end{align}
The leading critical limit appears in manuscript [019]\cite{Source019};
it was also recorded as an observation in the parent intersection manuscript.
The following proof supplies a stronger expansion with a uniform first correction.

\begin{theorem}\label{thm:limit}
Fix $0<a_-<a_+<\infty$. Uniformly for integers $n,m$ satisfying
$a_-\le a=m/n^2\le a_+$, put $s=1/a=n^2/m$. As $n\to\infty$,
\begin{equation}\label{eq:expansion}
r_{n,m}=\gamma+\frac{2K_0(2\sqrt s)}{I_0(2\sqrt s)}
       +\frac{s-1}{n I_0(2\sqrt s)^2}+O(n^{-2}).
\end{equation}
Consequently, if $m/n^2\to a_0>0$, then
\[
r_{n,m}\longrightarrow
\gamma+\frac{2K_0(2/\sqrt{a_0})}{I_0(2/\sqrt{a_0})}.
\]
On the exact square diagonal,
\begin{equation}\label{eq:square}
r_{n,n^2}=\gamma+2\kappa+O(n^{-2}).
\end{equation}
\end{theorem}
\begin{proof}
Write $j=n-k$ and normalize by $f_n$. The exact weights are
\begin{equation}\label{eq:weights}
w_{n,j}=\frac{f_{n-j}}{f_n}
 =\frac{s^j}{(j!)^2}\prod_{r=0}^{j-1}
       \frac{(1-r/n)^2}{1+s/n-sr/n^2},\qquad 0\le j\le n.
\end{equation}
In particular $0\le w_{n,j}\le s^j/(j!)^2$. Uniformly for
$j\le n^{1/4}$, expanding the logarithm of the product yields
\begin{align}
w_{n,j}&=\frac{s^j}{(j!)^2}
\left(1-\frac{\mathcal B_j}{n}
       +O\left(\frac{(j+1)^4}{n^2}\right)\right),
& \mathcal B_j&=j^2+(s-1)j.\label{eq:wexp}
\end{align}
Indeed the logarithmic first term is
$-[2\sum_{r<j}r+sj]/n=-\mathcal B_j/n$, and its remainder is
$O((j+1)^3/n^2)$; exponentiation gives \eqref{eq:wexp}.
The expansion $H_N=\log N+\gamma+1/(2N)+O(N^{-2})$ gives, on the same range,
\begin{equation}\label{eq:hexp}
h_{n-j}=\log a-\gamma+2H_j
       +\frac{s-1+2j}{n}
       +O\left(\frac{(j+1)^2}{n^2}\right).
\end{equation}
All implied constants may depend on $a_-,a_+$ but not on $m,n,j$.

We justify summation before computing coefficients. The weights are
dominated by $B^j/(j!)^2$ for a fixed $B>0$. Also
\[
h_{n-j}=A_n+\Delta_{n,j},\qquad
A_n=H_{m+n}-2H_n=O(1),\qquad
|\Delta_{n,j}|\le5H_j.
\]
For example, $H_N-H_{N-j}\le H_j$ for $N\ge j$, which bounds each
harmonic difference. The tails beyond $n^{1/4}$ of the dominating
factorial series, multiplied by any fixed polynomial in $j$ and by
$1+H_j$, are $O(n^{-M})$ for every fixed $M$. Thus
\eqref{eq:wexp}--\eqref{eq:hexp} may be summed with error $O(n^{-2})$.

Set $D(s)=I_0(2\sqrt s)$ and introduce the probability distribution
\[
\Pr(J=j)=\frac{s^j}{(j!)^2D(s)},\qquad j\ge0.
\]
Expanding the numerator and denominator of \eqref{eq:r} now gives
\begin{equation}\label{eq:covcoeff}
r_{n,m}=\log a-\gamma+2\Ex H_J
 +\frac{s-1+2\Ex J-2\Cov(H_J,B_J)}{n}+O(n^{-2}).
\end{equation}
The relevant moments are
\begin{equation}\label{eq:covs}
\Ex J^2=s,\qquad
\Cov(H_J,J^2)=\Ex J,\qquad
\Cov(H_J,J)=\frac{1-D(s)^{-2}}2.
\end{equation}
The first two follow from $j^2s^j/(j!)^2=s\,s^{j-1}/((j-1)!)^2$
and $H_{j+1}=H_j+1/(j+1)$. For the third, differentiating the
normalized series gives
\[
\Cov(H_J,J)=s\frac{d}{ds}\frac{Y(2\sqrt s)}{I_0(2\sqrt s)}.
\]
By \eqref{eq:beta} and \eqref{eq:wronskian},
\[
\frac{d}{dt}\left(\frac{Y(t)}{I_0(t)}\right)
=\frac1t-\frac{1}{tI_0(t)^2},
\]
which proves the third identity in \eqref{eq:covs}.
The coefficient in \eqref{eq:covcoeff} therefore simplifies to
\[
s-1+2\Ex J
-2\left(\Ex J+\frac{s-1}{2}(1-D^{-2})\right)
=\frac{s-1}{D^2}.
\]
Finally,
$\Ex H_J=Y(2\sqrt s)/I_0(2\sqrt s)$ and
$\log a=-\log s$ turn the constant term in \eqref{eq:covcoeff}
into the leading term of \eqref{eq:expansion}. For $m=n^2$,
$s=1$, so its entire $1/n$ correction vanishes.
\end{proof}

\begin{remark}
In the uniform expansion, $s=n^2/m$ is the current parameter. If only
$m/n^2\to a_0$ is known, substituting $1/a_0$ for $s$ while retaining
the $O(n^{-2})$ error is not justified without a rate of convergence.
Equation \eqref{eq:square} has no such parameter error.
\end{remark}

\subsection{Connecting two actual rational approximation limits}
Let
\[
L_+=\lim_{n\to\infty}r_{n,n^2}=\gamma+2\kappa.
\]
The parent manuscript's critical Van Assche--Wolfs limit is
$L_-=\gamma-2\kappa$ \cite[Section 3.5]{Intersection}. For precision, its
approximants are
\[
b_{n,p,k}=\frac{\binom nk^2\binom{n+k}{k}^{p}}{k!},\qquad
T_{n,p}=\frac{\sum_{k=0}^n b_{n,p,k}
 ((p+3)H_k-2H_{n-k}-pH_{n+k})}{\sum_{k=0}^n b_{n,p,k}},
\]
and the input identification is $\lim_{q\to\infty}T_{2^q,3q}=L_-$.
The arithmetic consequences of the two mirror values are
Corollary~\ref{cor:limits} and the preceding one-defect statement.

\subsection{A distinct critical Newton--Laguerre limit}
The following finite sums are the Newton--Laguerre approximants of \cite{Gamma}.
\begin{proposition}[Exact critical Newton--Laguerre limit]
\label{prim:limit}
Let
\[
 L_n(x)=\sum_{j=0}^n(-1)^j\binom nj\frac{x^j}{j!},
 \qquad R_N(x)=-\sum_{n=1}^N\frac{L_n(x)}n.
\]
For every $\lambda>0$,
\[
 R_N(\lambda/N)-\log(\lambda/N)\longrightarrow
 \mathcal L(\lambda):=
 \sum_{j\ge1}\frac{(-1)^{j-1}\lambda^j}{j(j!)^2}
       -\gamma-\log\lambda.
\]
In particular
\[
 \mathcal L:=\mathcal L(1)=B-\gamma
 =0.3066231395146690006884066448658841691\ldots.
\]
\end{proposition}
\begin{proof}
Since $\binom nk/n=\binom{n-1}{k-1}/k$, the hockey-stick identity gives
\[
 \sum_{n=k}^N\frac{\binom nk}{n}=\frac1k\binom Nk.
\]
Consequently
\[
 R_N(x)=-H_N+
  \sum_{k=1}^N\frac{(-1)^{k-1}\binom Nk x^k}{k\,k!}.
\]
At $x=\lambda/N$, the absolute value of each summand is at most
$\lambda^k/[k(k!)^2]$, a summable sequence, and it converges to the
corresponding signed term. Dominated convergence and
$H_N-\log N\to\gamma$ give the stated limit.
\end{proof}


\part{Second-kind Bessel values and density}
% Unified manuscript, part 05: affine circuits, density, and second-kind witnesses.
% Extracted with proofs from source 5; conditional density extension appended.
% Source originals are unchanged.

\section{Affine multiplicative circuits and second-kind Bessel values}
\label{kd:section}

All results of this section are deductions from
Theorem~\ref{cf:joint}. They are not assertions of historical priority.
Write $k=\overline{\mathbb Q}$ and
$K_{\exp}=k(e^\alpha:\alpha\in k)$.  For distinct positive algebraic
numbers $t_1,\ldots,t_m$, the part of Theorem~\ref{cf:joint} used here says that
\[
 I_j=I_0(t_j),\qquad C_j=I_1(t_j)/I_0(t_j),\qquad
 B_j=\beta(t_j)
 \quad(1\leq j\leq m)
\]
are jointly algebraically independent over $K_{\exp}$.  The connection
identity is
\begin{equation}\label{kd:connection}
 B_j=\gamma+\log(t_j/2)+\frac{K_0(t_j)}{I_j}.
\end{equation}
Only positive arguments and real logarithms occur in this section.

For a finite set $T=\{t_1,\ldots,t_m\}$, define its space of affine
multiplicative circuits by
\[
 W_T=\left\{q\in\mathbb Q^m:
       \sum_jq_j=0,\quad\sum_jq_j\log t_j=0\right\},
 \qquad d_T=m-\dim_{\mathbb Q}W_T.
\]
The logarithmic condition is equivalently $\prod_j t_j^{q_j}=1$.
If the logarithms of the $t_j$ span a rational vector space of dimension
$r$, then $d_T\leq r+1$; $d_T$ records the potentially sharper affine
rank.

\begin{theorem}[Affine rank bound]\label{kd:rank}
For every finite set $T$ as above,
\begin{equation}\label{kd:rank-bound}
 \operatorname{trdeg}_{K_{\exp}}
 K_{\exp}\bigl(K_0(t):t\in T\bigr)\geq |T|-d_T.
\end{equation}
More precisely, if $q^{(1)},\ldots,q^{(m-d_T)}$ form a rational basis of
$W_T$, then the numbers
\begin{equation}\label{kd:residuals}
 R_a=\sum_{j=1}^m q_j^{(a)}\frac{K_0(t_j)}{I_0(t_j)}
 \qquad(1\leq a\leq m-d_T)
\end{equation}
are algebraically independent over
$K_{\exp}(I_1,C_1,\ldots,I_m,C_m)$.
\end{theorem}

\begin{proof}
Put $F=K_{\exp}(I_1,C_1,\ldots,I_m,C_m)$.  Theorem~\ref{cf:joint} makes the
$B_j$ algebraically independent over $F$.  A matrix whose rows are the
$q^{(a)}$ has full row rank, so its linear combinations of the $B_j$
are algebraically independent over $F$: extend its rows to an invertible
matrix and change coordinates.  Both the common $\gamma$ term and the
logarithms in \eqref{kd:connection} cancel along each row.  Consequently
these linear combinations are precisely \eqref{kd:residuals}.
They belong to $F(K_0(t_1),\ldots,K_0(t_m))$, whence this field has
transcendence degree at least $m-d_T$ over $F$.  Extending the base field
cannot increase the transcendence degree of a fixed tuple, proving
\eqref{kd:rank-bound}.
\end{proof}

\begin{corollary}[Restrictions on algebraic values]\label{kd:exceptions}
If all $K_0(t_j)$ are algebraic, then $W_T=0$.  Thus the set of positive
algebraic arguments at which $K_0$ is algebraic contains no affine
multiplicative circuit.  Its intersection with a multiplicative group
whose real logarithms have finite rational dimension $r$ has at most
$r+1$ elements.
\end{corollary}

\begin{proof}
Algebraicity makes the left side of \eqref{kd:rank-bound} zero, forcing
$m=d_T$.  Any $r+2$ points in the indicated group would have
$d_T\leq r+1<m$.
\end{proof}

\begin{corollary}[Geometric progressions]\label{kd:geometric}
Let $a,r>0$ be algebraic, with $r\ne1$.  Among the values
\[
                  K_0(ar^n),\qquad n\in\mathbb Z,
\]
at most two are algebraic.  For any $m\geq2$ distinct integer exponents
$n_1,\ldots,n_m$,
\[
 \operatorname{trdeg}_{K_{\exp}}
 K_{\exp}\bigl(K_0(ar^{n_1}),\ldots,K_0(ar^{n_m})\bigr)
 \geq m-2.
\]
In particular, at least one of $K_0(2),K_0(4),K_0(8)$ is transcendental.
More strongly, one can delete at most two values from the entire
family indexed by $n\in\mathbb Z$ so that all the remaining values are
algebraically independent over $K_{\exp}$.
\end{corollary}

\begin{proof}
On the hyperplane $\sum q_j=0$, the logarithmic condition becomes
$(\log r)\sum n_jq_j=0$.  Since $\log r\ne0$ and the exponents are
distinct, these two rational equations are independent and $d_T=2$.
Theorem~\ref{kd:rank} proves the rank bound; three algebraic values
would contradict it.  For three consecutive exponents the circuit is
$(1,-2,1)$.  No invocation of Hermite--Lindemann is needed for this
deduction from Theorem~\ref{cf:joint}.

To obtain the last assertion, enumerate the integer exponents and
greedily retain a value precisely when it is transcendental over the
field of previously retained values.  At each finite stage with $m$
values the retained tuple is a transcendence basis of the generated
field.  At most two values have been rejected, because its rank is at
least $m-2$.  The total number of rejections can therefore never exceed
two.  The retained infinite family is algebraically independent.
\end{proof}

\begin{corollary}[A finite-rank multiplicative coset]
\label{kd:coset}
Let $G\subset k\cap\mathbb R_{>0}$ be a multiplicative group such that
$\dim_{\mathbb Q}\operatorname{span}_{\mathbb Q}\{\log g:g\in G\}=r<\infty$,
and let $a>0$ be algebraic.  From any family of distinct arguments in
$aG$ one can delete at most $r+1$ arguments so that the remaining
$K_0$ values are algebraically independent over $K_{\exp}$.
\end{corollary}

\begin{proof}
Every finite set $T\subset aG$ has $d_T\leq r+1$: on the hyperplane
$\sum q_j=0$, the common term $\log a$ cancels.  Thus its $K_0$ values
have rank at least $|T|-r-1$.  The family is countable since its
arguments are algebraic.  Apply the same greedy construction as in
Corollary~\ref{kd:geometric}; at most $r+1$ rejections are possible.
\end{proof}

\begin{theorem}[Density one and an independent subsequence]
\label{kd:density}
Let $\pi(N)$ be the number of primes at most the positive integer $N$.
Then
\begin{equation}\label{kd:integer-rank}
 \operatorname{trdeg}_{K_{\exp}}
 K_{\exp}\bigl(K_0(1),\ldots,K_0(N)\bigr)
 \geq N-\pi(N)-1.
\end{equation}
In particular,
\[
 \#\{1\leq n\leq N:K_0(n)\in k\}\leq\pi(N)+1,
\]
so $K_0(n)$ is transcendental for a set of positive integers of natural
density one.  There is, moreover, a single subset $S\subseteq\mathbb N$
such that
\begin{equation}\label{kd:single-subsequence}
 \#(S\cap[1,N])\geq N-\pi(N)-1\quad(N\geq1)
\end{equation}
and the entire family $\{K_0(n):n\in S\}$ is algebraically independent
over $K_{\exp}$.  Thus $S$ has natural density one.
\end{theorem}

\begin{proof}
The primes up to $N$ generate a multiplicative group containing
$1,\ldots,N$.  The augmented exponent vectors of $1$ and of these
primes are linearly independent, and all remaining exponent vectors
belong to their span.  Hence $d_{\{1,\ldots,N\}}=\pi(N)+1$, and
Theorem~\ref{kd:rank} proves \eqref{kd:integer-rank}.  Corollary
\ref{kd:exceptions}, applied to the algebraic values in this interval,
proves the counting bound.

For the stronger assertion, process $n=1,2,\ldots$ in order.  Include
$n$ in $S$ exactly when $K_0(n)$ is transcendental over the field
generated over $K_{\exp}$ by the previously included values.  Every
finite selected tuple is independent.  Every rejected value is
algebraic over the previously selected tuple, so at every stage $N$
the selected tuple is a transcendence basis for the field generated by
the first $N$ values.  Its size is the left side of
\eqref{kd:integer-rank}; this proves \eqref{kd:single-subsequence}.
Finally $\pi(N)=o(N)$.  This argument proves existence of $S$; it does
not provide an effective test for membership or identify an individual
transcendental value such as $K_0(2)$.
\end{proof}

\subsection{The value at two as an additional witness}

Put
\[
 I=I_0(2),\quad K=K_0(2),\quad
 \kappa=K/I,\quad C=I_1(2)/I,
 \qquad H_4=k(\gamma,\delta,\kappa,C).
\]
The relevant part of Theorem~\ref{cf:joint} says that
\begin{equation}\label{kd:five-input}
 e,\ I,\ \eta=\Ein(1)=\gamma+\delta/e,
 \ C,\ \beta(2)=\gamma+\kappa
\end{equation}
are algebraically independent over $k$.

\begin{proposition}[A second-kind value as a rank witness]
\label{kd:K-witness}
If $K_0(2)$ is algebraic over $H_4$, then
\[
                  \operatorname{trdeg}_k H_4=4.
\]
In particular, algebraicity of $K_0(2)$ over $k$ would close the original
four-constant independence problem.  Conversely, if the original four
constants are algebraically dependent, then $e$ and $K_0(2)$ are jointly
algebraically independent over $H_4$.  In that case $K_0(2)$ is
transcendental even over $H_4$.
\end{proposition}

\begin{proof}
All five independent numbers in \eqref{kd:five-input} belong to
$H_4(e,I)$.  If $K$ is algebraic over $H_4$, then so is $I=K/\kappa$,
since $\kappa>0$.  Therefore
\[
 5\leq\operatorname{trdeg}_k H_4(e,I)
   \leq\operatorname{trdeg}_k H_4+1\leq5,
\]
forcing $\operatorname{trdeg}_k H_4=4$.

More generally \eqref{kd:five-input} gives
$\operatorname{trdeg}_k H_4\geq3$.  If the original four are dependent,
this degree is exactly three.  The field $H_4(e,I)$ then has
transcendence degree exactly five over $k$, and hence degree two over
$H_4$.  Thus $e,I$ are independent over $H_4$.  Replacing $I$ by the
nonzero $H_4$-multiple $K=\kappa I$ proves the assertion.
\end{proof}

\begin{corollary}[Six-constant consequence]\label{kd:six-witness}
Using the complex-scale conclusion of Theorem~\ref{cf:joint}, put
\[
 \nu=-\frac{\pi}{2}\frac{Y_0(2)}{J_0(2)},\qquad
 C_J=\frac{J_1(2)}{J_0(2)}.
\]
If $K_0(2)$ is algebraic over $H_4$, then
$\gamma,\delta,\kappa,C,\nu,C_J$ are jointly algebraically independent
over $K_{\exp}$.
\end{corollary}

\begin{proof}
Write $H_6=k(\gamma,\delta,\kappa,C,\nu,C_J)$ and $J=J_0(2)$.
The seven numbers
\[
 \eta,I,C,\gamma+\kappa,J,C_J,\gamma+\nu
\]
are independent over $K_{\exp}$.  They belong to
$K_{\exp}H_6(I,J)$, in which $I$ is algebraic over $K_{\exp}H_6$
by the hypothesis.  Consequently
\[
 7\leq\operatorname{trdeg}_{K_{\exp}}K_{\exp}H_6(I,J)
 \leq\operatorname{trdeg}_{K_{\exp}}K_{\exp}H_6+1\leq7.
\]
The six generators of $H_6$ are therefore independent over $K_{\exp}$.
\end{proof}

\section{A smaller circuit defect if the original four are dependent}
\label{cden:section}

The affine circuit theorem has a sharper relative form in the
rank-three alternative. The improvement comes from putting the
common constant $\gamma$ in the coefficient field: circuit rows then
need to cancel only logarithms, not their coefficient sum.

Put
\[
 k=\overline{\mathbb Q},\qquad
 K_{\exp}=k(e^\alpha:\alpha\in k),\qquad
 H_4=k(\gamma,\delta,\kappa,C),\qquad F=K_{\exp}H_4,
\]
where $\kappa=K_0(2)/I_0(2)$ and $C=I_1(2)/I_0(2)$.
Throughout this section assume
\begin{equation}\label{cden:rank-three}
                 \operatorname{trdeg}_k H_4=3.
\end{equation}
Theorem~\ref{cf:joint} supplies the independence, over $K_{\exp}$, of
$\eta=\Ein(1)$ together with
$I_0(t),C(t),\beta(t)$ at any finite set of distinct positive
algebraic arguments. Recall that
\[
 \eta=\gamma+\delta/e,\qquad
 \beta(t)=\gamma+\log(t/2)+\frac{K_0(t)}{I_0(t)}.
\]
All logarithms in this section are real logarithms of positive numbers.

\begin{lemma}[The relative joint input in the deficient case]
\label{cden:relative-input}
Under \eqref{cden:rank-three},
$\operatorname{trdeg}_{K_{\exp}}F=3$. Moreover,
\[
 I_0(2),\qquad I_0(t),\ C(t),\ \beta(t)
 \quad(t\in k\cap\mathbb R_{>0},\ t\ne2)
\]
form an algebraically independent family over $F$, in the
finite-subset sense. The same assertion holds with $I_0(2)$ replaced
by $K_0(2)$.
\end{lemma}

\begin{proof}
The three independent values $\eta,\beta(2),C$ belong to $F$,
because $e\in K_{\exp}$. Thus
$\operatorname{trdeg}_{K_{\exp}}F\ge3$. Base extension cannot
increase the transcendence degree of the fixed tuple generating
$H_4$, so \eqref{cden:rank-three} gives equality.
For $m$ distinct additional scales $t_j\ne2$, the joint input gives
$4+3m$ independent values over $K_{\exp}$:
\[
 \eta,\beta(2),C,I_0(2),\qquad
 I_0(t_j),C(t_j),\beta(t_j)\quad(1\le j\le m).
\]
They lie in a field obtained from $F$, of relative transcendence
degree three over $K_{\exp}$, by adjoining the displayed $1+3m$
additional generators. The tower formula forces their relative
transcendence degree over $F$ to be exactly $1+3m$.
Finally $K_0(2)=\kappa I_0(2)$, with $0\ne\kappa\in F$, is an
invertible change of coordinate over $F$.
\end{proof}

\begin{theorem}[Logarithmic rank bound over the target field]
\label{cden:rank}
Let $T$ be a finite set of distinct positive algebraic numbers, and put
\[
 r_2(T)=\dim_{\mathbb Q}
       \operatorname{span}_{\mathbb Q}\{\log(t/2):t\in T\}.
\]
Under \eqref{cden:rank-three},
\begin{equation}\label{cden:bound}
 \operatorname{trdeg}_{F}
 F\bigl(K_0(t):t\in T\bigr)\ge |T|-r_2(T).
\end{equation}
\end{theorem}

\begin{proof}
Put $K_2=K_0(2)$, $F_0=F(K_2)$, and
$T'=T\setminus\{2\}$. Write $m=|T'|$ and enumerate its points as
$t_1,\ldots,t_m$. The logarithms for $T'$ have rank $r_2(T)$,
since the removed logarithm, if present, was zero.
By Lemma~\ref{cden:relative-input}, the $m$ values $\beta(t_j)$
are independent over
\[
 G=F_0\bigl(I_0(t_j),C(t_j):1\le j\le m\bigr).
\]
Choose a rational basis of the kernel of
\[
 \mathbb Q^m\longrightarrow\mathbb R,
 \qquad(q_j)\longmapsto\sum_jq_j\log(t_j/2).
\]
This basis has $m-r_2(T)$ rows. Its linear combinations of the
independent $\beta(t_j)$ remain independent over $G$. Subtracting
$\gamma\sum_jq_j\in F\subset G$ from each combination gives the
same number of independent residuals
\[
                 \sum_jq_j\frac{K_0(t_j)}{I_0(t_j)}.
\]
All lie in $G(K_0(t_1),\ldots,K_0(t_m))$. Hence
\[
 \operatorname{trdeg}_{F_0}
 F_0(K_0(t):t\in T')\ge m-r_2(T),
\]
using base-extension monotonicity as in the affine circuit argument.
If $2\notin T$, extending the base from $F$ to $F_0$ cannot
increase the rank of this tuple, and this proves
\eqref{cden:bound}. If $2\in T$, the element $K_2$ contributes one
independent dimension over $F$ by
Lemma~\ref{cden:relative-input}. The tower formula gives
$1+m-r_2(T)=|T|-r_2(T)$, as required.
\end{proof}

\begin{corollary}[A one-deletion geometric family]
\label{cden:geometric}
Let $r>0$ be algebraic with $r\ne1$. Under
\eqref{cden:rank-three}, one can delete at most one value from
\[
                   \{K_0(2r^n):n\in\mathbb Z\}
\]
so that the remaining entire family is algebraically independent
over $F=K_{\exp}H_4$. The retained family can be chosen to include
$K_0(2)$. At most one member of the displayed family is algebraic
over $F$.
More generally, if a multiplicative subgroup
$G_0\subset k\cap\mathbb R_{>0}$ has logarithmic rational rank $r$,
then any family of distinct arguments in $2G_0$ has a deletion budget
of at most $r$ over $F$.
\end{corollary}

\begin{proof}
For every finite set of exponents,
$\log((2r^n)/2)=n\log r$ spans a rational vector space of dimension
at most one. Theorem~\ref{cden:rank} therefore gives a rank deficit
of at most one for every finite subfamily. Enumerate $\mathbb Z$
starting with zero, retaining a value exactly when it is
transcendental over the previously retained field over $F$.
The first value $K_0(2)$ is retained by
Lemma~\ref{cden:relative-input}. At each finite stage, the retained
values form a transcendence basis of all values seen so far:
rejected values remain algebraic over the enlarging retained field.
Thus at most one rejection can ever occur. Every finite retained
tuple is independent, proving the infinite assertion. Two values
algebraic over $F$ would contradict the rank bound for that pair.
For $2G_0$, use $r_2(T)\le r$ and the identical greedy argument.
The families are countable because their arguments are algebraic.
\end{proof}

\begin{corollary}[A sharper density-one family containing the point two]
\label{cden:density}
Under \eqref{cden:rank-three}, for every integer $N\ge2$,
\begin{equation}\label{cden:integer-bound}
 \operatorname{trdeg}_F F(K_0(1),\ldots,K_0(N))
                       \ge N-\pi(N).
\end{equation}
There is a single set $S\subset\mathbb N$, with $2\in S$, such that
\[
 \#(S\cap[1,N])\ge N-\pi(N)\quad(N\ge2),
\]
and $\{K_0(n):n\in S\}$ is algebraically independent over
$F=K_{\exp}H_4$. In particular $S$ has natural density one.
At most $\pi(N)$ of the first $N$ values can be algebraic over $F$.
\end{corollary}

\begin{proof}
For $N\ge2$, the rational span of
$\log(n/2)$, $1\le n\le N$, is exactly the span of
$\log p$ for primes $p\le N$. One inclusion follows by factorization;
for the other, $\log(1/2)=-\log2$ is available and
$\log p=\log(p/2)+\log2$. Unique factorization makes the prime
logarithms rationally independent, so the rank is $\pi(N)$.
Theorem~\ref{cden:rank} proves \eqref{cden:integer-bound}.
Process the arguments in the order $2,1,3,4,\ldots$, using the same
greedy rule as above. The first value is retained, and for every
$N\ge2$ the first $N$ processed arguments form exactly $\{1,\ldots,N\}$.
Their retained count equals their transcendence degree, giving the
stated bound. Since $\pi(N)=o(N)$, the density is one.
For the final assertion, apply \eqref{cden:bound} to the subset of
arguments whose values are algebraic over $F$; its logarithmic rank
is at most $\pi(N)$.
\end{proof}

These are conditional statements throughout. They improve the
unconditional affine deletion budget $r+1$ to $r$ on the centered
coset $2G_0$, and the unconditional integer defect $\pi(N)+1$ to
$\pi(N)$, while enlarging the coefficient field to $K_{\exp}H_4$.
They do not establish \eqref{cden:rank-three} for the actual constants,
do not identify an unconditional individual transcendental
$K_0(n)$, and give no effective membership test for $S$.
Equivalently, a failure of any stated rank bound over $K_{\exp}H_4$
would force $\operatorname{trdeg}_kH_4=4$, using the already known
lower bound three.

\section{A common logarithmic defect for six families}
\label{jden:section}

The modified- and ordinary-Bessel blocks share the same logarithmic
coordinate. Combining them with the exponential integral and the critical
primitive limit therefore costs one common affine logarithmic rank,
not a separate rank for each family. The results below are consequences
of the joint-value theorems already proved; no assertion of historical
priority is intended.

For a positive algebraic number $\lambda$, write
\begin{align}
 h_\lambda&=E_1^{\rm tail}(\lambda)
       =\int_\lambda^\infty\frac{e^{-u}}u\,du,
 &g_\lambda&=\gamma+\log\lambda,\label{jden:tail}\\
 B(\lambda)&=\sum_{n\ge1}\frac{(-1)^{n-1}\lambda^n}{n(n!)^2},
 &\mathcal L(\lambda)&=B(\lambda)-g_\lambda.\label{jden:limit}
\end{align}
The tail $h_\lambda$ is not the scaled quantity
$e^\lambda E_1^{\rm tail}(\lambda)$. Put
\begin{equation}\label{jden:coordinates}
 Z_\lambda=
 \bigl(h_\lambda,K_0(2\lambda),K_1(2\lambda),
             \pi Y_0(2\lambda),\pi Y_1(2\lambda),
             \mathcal L(\lambda)\bigr).
\end{equation}
Here $Y_r$ is the classical Bessel function of the second kind. For a
finite set $T=\{\lambda_1,\ldots,\lambda_m\}$ of distinct positive
algebraic numbers, use the affine logarithmic rank
\begin{equation}\label{jden:affine}
 W_T=\left\{q\in\mathbb Q^m:
       \sum_jq_j=0,\quad\sum_jq_j\log\lambda_j=0\right\},
 \qquad d_T=m-\dim_{\mathbb Q}W_T.
\end{equation}
This agrees with the notation in Section~\ref{kd:section}. Multiplying
all arguments by one fixed positive algebraic number leaves $W_T$ and
$d_T$ unchanged.

\begin{theorem}[Six families with one affine defect]\label{jden:rank}
For every finite set $T$ as above,
\begin{equation}\label{jden:sixrank}
 \trdeg_{\Kexp}\Kexp(Z_{\lambda_j}:1\le j\le m)
 \ge 6m-d_T.
\end{equation}
In particular, each of the following two $4m$-tuples has transcendence
degree at least $4m-d_T$ over $\Kexp$:
\begin{align}
 &(h_{\lambda_j},K_0(2\lambda_j),\pi Y_0(2\lambda_j),
                          \mathcal L(\lambda_j))_{j=1}^m,
                          \label{jden:fourrank}\\
 &(K_0(2\lambda_j),K_1(2\lambda_j),
                  \pi Y_0(2\lambda_j),\pi Y_1(2\lambda_j))_{j=1}^m.
                          \label{jden:pairedrank}
\end{align}
The tuple $(K_0(2\lambda_j),K_1(2\lambda_j))_{j=1}^m$ has
transcendence degree at least $2m-d_T$, and the same bound holds for
the ordinary-Bessel pair with the factor $\pi$.
\end{theorem}

\begin{proof}
Set $t_j=2\lambda_j$, $I_{r,j}=I_r(t_j)$,
$J_{r,j}=J_r(t_j)$, $P_{r,j}=\pi Y_r(t_j)$, and
$g_j=g_{\lambda_j}$. Theorems~\ref{cf:joint} and
\ref{msprim:values} show that the following $8m$ coordinates are
jointly algebraically independent over $\Kexp$:
\begin{equation}\label{jden:eight}
 \bigl(\Ein(\lambda_j),I_{0,j},I_{1,j},\beta(t_j),
          J_{0,j},J_{1,j},\beta(it_j),B(\lambda_j)\bigr)_{j=1}^m.
\end{equation}
Indeed, the Bessel scales $t_j,it_j$ are distinct modulo sign, and the
primitive scales $2i\sqrt{\lambda_j}$ are distinct modulo sign.
A primitive scale is allowed to coincide with an old Bessel scale:
the primitive theorem gives independence over all the old coordinates.
Replacing the quotients $C$ by $I_1$ or $J_1$ is an invertible rational
change since $I_{0,j}$ and $J_{0,j}$ are nonzero.

The connection identities are
\begin{align}
 \Ein(\lambda_j)&=g_j+h_{\lambda_j},&
 B(\lambda_j)&=g_j+\mathcal L(\lambda_j),\label{jden:recovery1}\\
 \beta(t_j)&=g_j+K_0(t_j)/I_{0,j},&
 \beta(it_j)&=g_j-P_{0,j}/(2J_{0,j}).\label{jden:recovery2}
\end{align}
The Wronskians, with their signs fixed by
\cite[\S\S10.5, 10.28]{DLMF}, give
\begin{equation}\label{jden:wronskians}
 I_{0,j}K_1(t_j)+I_{1,j}K_0(t_j)=1/t_j,
 \qquad
 J_{1,j}P_{0,j}-J_{0,j}P_{1,j}=2/t_j.
\end{equation}
Since $K_0(t_j)>0$, adjoining $I_{0,j}$ suffices to recover $I_{1,j}$
from the modified-Bessel pair. If $P_{0,j}\ne0$, adjoining $J_{0,j}$
suffices to recover $J_{1,j}$. If $P_{0,j}=0$, then
$P_{1,j}\ne0$ and $J_{0,j}=-2/(t_jP_{1,j})$; in this case adjoin
$J_{1,j}$ instead. Thus exactly one freely adjoined first-kind
coordinate per Bessel block suffices, including the possible zero case.

Every vector in $W_T$ gives a rational linear relation among the $g_j$.
Consequently
$\trdeg_{\Kexp}\Kexp(g_1,\ldots,g_m)\le d_T$.
The field generated by \eqref{jden:eight} is contained in the field
generated by the $Z_{\lambda_j}$, the $g_j$, and at most $2m$
additional first-kind coordinates. Hence
\[
 8m\le
 \trdeg_{\Kexp}\Kexp(Z_{\lambda_j}:1\le j\le m)+d_T+2m,
\]
which proves \eqref{jden:sixrank}. Deleting $r$ entries from a tuple
can lower its transcendence degree by at most $r$. Deleting $2m$ or
$4m$ entries proves all the stated projections.
\end{proof}

\begin{corollary}[No exceptional positive algebraic point]
\label{jden:everypoint}
For every positive algebraic $\lambda$, the six entries of
$Z_\lambda$ have transcendence degree at least five over $\Kexp$,
and
\begin{equation}\label{jden:everyfour}
 \trdeg_{\Kexp}\Kexp
 \bigl(E_1^{\rm tail}(\lambda),K_0(2\lambda),
             \pi Y_0(2\lambda),\mathcal L(\lambda)\bigr)\ge3.
\end{equation}
If $g_\lambda$ is algebraic over $\Kexp$, all six entries are
algebraically independent over $\Kexp$; in particular the four
entries in \eqref{jden:everyfour} are independent.
\end{corollary}

\begin{proof}
A singleton has $d_T=1$. If $g_\lambda$ is algebraic over $\Kexp$,
the recovery proof costs no transcendence degree for that coordinate,
and yields the stronger lower bound six.
\end{proof}

These conclusions identify neither a prescribed independent triple nor
a prescribed independent five-tuple. In particular, they assert that
at most one of the four, or of the six, displayed entries is algebraic
over $\Kexp$. There is no exceptional $\lambda$ in this statement.
An exception count for algebraic $\gamma+\log\lambda$ is relevant to
independent tuples that include that logarithmic coordinate, but is
unnecessary after projecting to the four scalar outputs.

\begin{lemma}[Integer-order changes of coordinates]\label{jden:orders}
Fix a positive algebraic $t$. For every integer $r\ge0$, each of
$K_r(t)$ and $\pi Y_r(t)$ is a nonzero $k$-linear form in its
respective order-zero and order-one pair. Every adjacent pair of
orders is obtained from that basic pair by an invertible $k$-linear
transformation.
\end{lemma}

\begin{proof}
The standard recurrences \cite[\S\S10.6, 10.29]{DLMF} read
\[
 K_{r+2}(t)=K_r(t)+\frac{2(r+1)}tK_{r+1}(t),\qquad
 Y_{r+2}(t)=-Y_r(t)+\frac{2(r+1)}tY_{r+1}(t).
\]
Their transition matrices on adjacent pairs have determinants $-1$
and $1$, respectively. Starting with orders zero and one, every
finite product is therefore invertible over $k$. Each row is nonzero,
which proves the assertion about individual linear forms as well.
\end{proof}

\begin{theorem}[A simultaneous density-one family]\label{jden:density}
For every positive integer $N$,
\begin{equation}\label{jden:integer}
 \trdeg_{\Kexp}\Kexp(Z_{n/2}:1\le n\le N)
       \ge6N-\pi(N)-1.
\end{equation}
There exists a single set $S\subseteq\mathbb N$ such that
\begin{equation}\label{jden:commonS}
 \#(S\cap[1,N])\ge N-\pi(N)-1\qquad(N\ge1)
\end{equation}
and all six coordinates of $Z_{n/2}$, as $n$ ranges over $S$, are
jointly algebraically independent over $\Kexp$. In particular, for
every $n\in S$ all the numbers
\begin{equation}\label{jden:allorders}
 E_1^{\rm tail}(n/2),\quad\mathcal L(n/2),\quad
 K_r(n),\quad\pi Y_s(n)\qquad(r,s\ge0)
\end{equation}
are individually transcendental over $\Kexp$.
For each selected $n$, either basic Bessel pair may be replaced by
any adjacent pair of nonnegative integer orders, and the resulting
six-coordinate family is still jointly independent. These choices
may vary with $n$.
\end{theorem}

\begin{proof}
The affine logarithmic rank of $\{1,\ldots,N\}$ is $\pi(N)+1$,
as proved in Theorem~\ref{kd:density}. The common scaling by $1/2$
preserves that rank, so Theorem~\ref{jden:rank} proves
\eqref{jden:integer}.

Process the six coordinates of $Z_{1/2}$, then those of $Z_{2/2}$,
and so on, in any fixed order within each block. Retain a coordinate
if and only if it is transcendental over the field generated over
$\Kexp$ by previously retained coordinates. At stage $N$ the
retained coordinates form a transcendence basis of the field generated
by the first $N$ blocks. Thus at most $\pi(N)+1$ coordinates have
been rejected. Define $S$ to consist of the indices of blocks in
which no coordinate was rejected. Each omitted block accounts for
at least one rejected coordinate, proving \eqref{jden:commonS}.
Every finite subset of the retained coordinates is independent, so
the entire six-coordinate family on $S$ is independent.

Lemma~\ref{jden:orders} gives invertible changes for adjacent pairs.
It also expresses each higher-order value as a nonzero algebraic
linear form in an independent pair. Such a form is transcendental
over $\Kexp$, proving \eqref{jden:allorders}. Finally
$\pi(N)=o(N)$ implies that $S$ has natural density one.
\end{proof}

\begin{corollary}[Counting all exceptional orders at once]
\label{jden:badcount}
Let $\mathcal A_N$ be the set of integers $1\le n\le N$ for which
at least one of the following conditions holds:
\begin{enumerate}
\item $E_1^{\rm tail}(n/2)$ or $\mathcal L(n/2)$ is algebraic over $\Kexp$;
\item $K_r(n)$ is algebraic over $\Kexp$ for some integer $r\ge0$;
\item $\pi Y_s(n)$ is algebraic over $\Kexp$ for some integer $s\ge0$.
\end{enumerate}
Then $\#\mathcal A_N\le\pi(N)+1$.
\end{corollary}

\begin{proof}
Put $m=\#\mathcal A_N$ and
$T=\{n/2:n\in\mathcal A_N\}$. At every such point, the six entries
of $Z_{n/2}$ satisfy a nonzero polynomial relation over $\Kexp$.
For the higher-order cases this follows from
Lemma~\ref{jden:orders}: substitute its nonzero linear form into a
minimal polynomial over $\Kexp$. Hence the transcendence degree of
the union of these $m$ blocks is at most $5m$. On the other hand,
Theorem~\ref{jden:rank} gives a lower bound $6m-d_T$. Therefore
$m\le d_T\le\pi(N)+1$, as required. The case $m=0$ is immediate.
\end{proof}

\begin{corollary}[A common finite deletion on a geometric progression]
\label{jden:geometric}
Let $a,q>0$ be algebraic with $q\ne1$. From the entire family of
six-coordinate blocks $Z_{aq^n}$, $n\in\mathbb Z$, one can delete
at most two blocks so that every remaining coordinate is jointly
algebraically independent over $\Kexp$. At every remaining argument,
the corresponding tail, critical limit, and all nonnegative integer
orders of both second-kind Bessel families are individually
transcendental over $\Kexp$.
\end{corollary}

\begin{proof}
Every finite set of at least two distinct arguments in this progression
has affine rank two. Enumerate the integer exponents and apply the
coordinate-by-coordinate construction of Theorem~\ref{jden:density}.
Theorem~\ref{jden:rank} permits at most two rejected coordinates in
any finite stage, and hence at most two rejected blocks in the entire
family. Lemma~\ref{jden:orders} supplies the last assertion.
\end{proof}

The results subsume the separate paired $K$-value bound and the separate
ordinary-Bessel density bound while retaining only one common
logarithmic defect. The infinite-order conclusion is individual
transcendence, not independence of all orders at one point: the
recurrences impose relations. Even replacing a basic ordinary-Bessel
pair by two arbitrary nonadjacent orders need not be invertible. For
example, at $t=\sqrt8$ one has $Y_3(t)=-4Y_0(t)/t$.
The density-one set is existential; these arguments do not provide an
effective membership test or identify a particular integer at which
a prescribed second-kind value is transcendental.

\part{Further applications}
\section{A joint classification of two different rank defects}

The linear-rank classification in \cite{Source021} can now be combined with the Bessel independence statement. Let
\[
 \rho=\dim_k\operatorname{span}_k\{1,e,\delta,e\eta\},\qquad
 H=k(\gamma,\delta,\kappa,C).
\]
The linear rank $\rho$ and the transcendence degree of $H$ are different
invariants. Since $e,\eta$ are algebraically independent, $1,e,e\eta$ are
linearly independent, and therefore $\rho\in\{3,4\}$.

\begin{theorem}[Rank-three linear defect versus joint algebraic defect]
\label{transfer:rank3}
Assume $\rho=3$. Then precisely the following alternatives hold:
\begin{enumerate}
\item $\gamma$ is algebraic; the triple $\delta,\kappa,C$ is algebraically
independent, and $\operatorname{trdeg}_k H=3$.
\item $\delta$ is algebraic; the triple $\gamma,\kappa,C$ is algebraically
independent, and $\operatorname{trdeg}_k H=3$.
\item Both $\gamma$ and $\delta$ are transcendental; the four numbers
$\gamma,\delta,\kappa,C$ are algebraically independent, and
$\operatorname{trdeg}_k H=4$.
\end{enumerate}
Consequently, if $H$ has a nonzero polynomial relation and both $\gamma$
and $\delta$ are transcendental, then $\rho=4$.
\end{theorem}
\begin{proof}
The condition $\rho=3$ and the independence of $1,e,e\eta$ give a unique
triple $a,b,c\in k$ such that
\begin{equation}\label{transfer:rank-param}
 \gamma=a+b/e+c\eta,\qquad
 \delta=(1-c)e\eta-ae-b.
\end{equation}
Conjugation and uniqueness make $a,b,c$ real. Independence of $e,\eta$
shows that $\gamma\in k$ occurs exactly when $b=c=0$, and
$\delta\in k$ exactly when $a=0,c=1$. The two cases are disjoint.

Outside them, eliminate $\eta$ in \eqref{transfer:rank-param} to obtain
\[
 e\bigl((1-c)\gamma-a\bigr)=c\delta+b.
\]
The denominator $(1-c)\gamma-a$ is nonzero: otherwise either $c=1,a=0$,
or $\gamma=a/(1-c)$ is algebraic, both excluded. Thus
\[
 e=\frac{c\delta+b}{(1-c)\gamma-a},\qquad
 \eta=\gamma+\delta/e,
\]
and $k(\gamma,\delta)=k(e,\eta)$. Write $\beta=\beta(2)=\gamma+\kappa$.
Theorem~\ref{cf:joint} makes $e,\eta,\beta,C$ algebraically independent. Therefore in the
nonexceptional case
\[
 H=k(e,\eta,\beta,C),
\]
which proves independence of its four displayed original generators.

If $\gamma=a\in k$, the birational substitution
$\delta=e(\eta-a)$ shows that $e,\delta,\beta,C$ are independent; replacing
$\beta$ by $\kappa=\beta-a$ proves the first case. If $\delta=-b\in k$, the
substitution $\gamma=\eta+b/e$ shows that $e,\gamma,\beta,C$ are independent;
replacing $\beta$ by $\kappa=\beta-\gamma$ proves the second case. Finally, the
last assertion is the contrapositive of case three, since $\rho\ge3$.
\end{proof}

This excludes simultaneous \emph{generic} defects of the two ranks. It does
not establish $\rho=4$ for the actual constants, does not establish a
polynomial relation in $H$, and does not establish the individual
transcendence of $\gamma$ or $\delta$. A generic linear rank-three defect is
compatible with, and indeed forces, maximal joint algebraic independence.
\section{An exact coefficient-field obstruction for a moment integral}
\label{sec:hamiltonian-transfer}

The positive-parameter Hamiltonian moment model in \cite{Hamiltonian}
provides a different application of the four-constant rank bound.  For
$a=(a_1,\ldots,a_n)\in\mathbb R_{>0}^n$, put
\begin{equation}\label{ham:model}
 F_a(t)=\int_{[0,1]^n}\exp\!\left(t\sum_{j=1}^n a_jx_j\right)\,dx
       =\prod_{j=1}^n\frac{e^{a_jt}-1}{a_jt},\qquad F_a(0)=1.
\end{equation}
The singularities in this product at $t=0$ are removable.  The following
criterion concerns the field of coefficients of a homogeneous scalar
linear differential equation, rather than its order.

\begin{lemma}[The visible phases generate the lengths]\label{ham:span}
After combining equal subset sums, write
\[
 G_a(t):=\prod_{j=1}^n(e^{a_jt}-1)
       =\sum_{\lambda\in\Sigma_a}c_\lambda e^{\lambda t},
 \qquad c_\lambda\in\mathbb Z\setminus\{0\}.
\]
Then
\[
 \operatorname{span}_{\mathbb Z}\Sigma_a
 =\operatorname{span}_{\mathbb Z}\{a_1,\ldots,a_n\}.
\]
\end{lemma}
\begin{proof}
Each visible phase is a subset sum, giving one inclusion.  For the
reverse inclusion, induct on the distinct lengths $b$ in increasing
order.  If the coefficient at exponent $b$ survives, then $b\in\Sigma_a$.
Otherwise the singleton subsets with length $b$ contribute a nonzero
integer of one common sign.  Their cancellation therefore requires a
subset of at least two lengths whose sum is $b$.  Positivity makes every
length in that subset strictly smaller than $b$; the induction hypothesis
expresses each of them in the integer span of $\Sigma_a$, and hence does
the same for $b$.  This also establishes the initial case, since the
smallest length cannot be canceled in this way.
\end{proof}

\begin{theorem}[Exact coefficient-field criterion]\label{ham:field}
Let $K\subset\mathbb C$ be a field containing $\mathbb Q$.  The function
$F_a$ satisfies a nonzero homogeneous scalar linear differential equation
over $K(t)$ if and only if every $a_j$ is algebraic over $K$.
\end{theorem}
\begin{proof}
Multiplication by $t^n$ preserves the existence of such an equation:
the corresponding operators are conjugate by $t^n$.  Multiplication by a
nonzero constant also preserves a homogeneous equation.  Since
$G_a=(\prod_j a_j)t^nF_a$, it is therefore enough to treat $G_a$; this
argument does not require $\prod_j a_j\in K$.

Suppose that an equation exists.  Clearing its denominators gives a
nonzero operator
\[
 L=\sum_{j=0}^r p_j(t)\partial_t^j,\qquad p_j(t)\in K[t],
 \qquad LG_a=0.
\]
Distinct real exponentials are linearly independent over $\mathbb C[t]$.
Indeed, in a putative relation, divide by the exponential of largest
rate and let $t\to+\infty$: its polynomial coefficient tends to zero
and must vanish; induction removes the remaining terms.  Consequently
\[
 P(t,\lambda):=\sum_{j=0}^r p_j(t)\lambda^j=0
 \quad\text{identically in }t\quad(\lambda\in\Sigma_a).
\]
Write the nonzero polynomial $P(t,X)$ as $\sum_d t^dR_d(X)$.
Any nonzero $R_d\in K[X]$ vanishes at every $\lambda\in\Sigma_a$.
Thus every visible phase is algebraic over $K$; Lemma~\ref{ham:span}
shows that every length is algebraic over $K$ as well.

Conversely, if all lengths are algebraic over $K$, so are the finitely
many visible phases.  A nonzero polynomial $R\in K[X]$ vanishing on
$\Sigma_a$ exists, for example the product of their minimal polynomials.
Then $R(\partial_t)G_a=0$, and the nonzero operator
$t^{-n}R(\partial_t)t^n\in K(t)[\partial_t]$ annihilates $F_a$.
\end{proof}

\begin{corollary}[The four-constant coefficient cost]\label{ham:application}
Let
\[
 H=\Qa(\gamma,\delta,\kappa,C),\qquad
 \delta=\int_0^\infty\frac{e^{-u}}{1+u}\,du,\quad
 \kappa=\frac{K_0(2)}{I_0(2)},\quad C=\frac{I_1(2)}{I_0(2)},
\]
where $\gamma$ is Euler's constant, with the rank bound
$\trdeg_{\Qa}H\ge3$ established in Theorem~\ref{thm:many}.
For $a=(\gamma,\delta,\kappa,C)$, let $\mathcal K_a$ denote the fields
$\Qa\subset K\subset\mathbb C$ over which $F_a$ has a nonzero homogeneous
linear differential equation.  Then
\begin{equation}\label{ham:minimum}
 \min_{K\in\mathcal K_a}\trdeg_{\Qa}K
   =\trdeg_{\Qa}H\in\{3,4\}.
\end{equation}
In particular, $F_a$ has no such equation over $\Qa(e,\pi)(t)$.
\end{corollary}
\begin{proof}
The four lengths are positive.  For $K\in\mathcal K_a$,
Theorem~\ref{ham:field} says that the compositum $HK$ is algebraic over
$K$.  Therefore
\[
 \trdeg_{\Qa}H\le\trdeg_{\Qa}HK=\trdeg_{\Qa}K.
\]
Taking $K=H$ realizes equality.  The rank input and the four generators
give the two possibilities in \eqref{ham:minimum}.  Finally,
$\trdeg_{\Qa}\Qa(e,\pi)\le2$, which is too small; no assertion about
the algebraic independence of $e$ and $\pi$ is needed.
\end{proof}

The conclusion is a coefficient-field obstruction for an explicit
entire moment integral.  An equation over $H(t)$, hence over
$\mathbb C(t)$, still exists.  This application does not increase the
rank bound from three to four or settle the individual arithmetic
nature of $\gamma$ or $\delta$.

\section{Scope, unresolved questions, and source concordance}
\label{sec:scope}
\subsection{What the enlarged families do and do not decide}
The exact ring intersections and the simultaneous primitive theorem
close specific questions about rational expressions and relative
independence. They do not establish
\[
 \trdeg_k k(\gamma,\delta,\kappa,C)=4.
\]
Nor do they establish the irrationality of either $\gamma$ or $\delta$.
The phrase ``at most one dependence'' is literal: no nonzero relation
among the original four is constructed, and no proof that all such
relations are absent is given. A rank witness is a sufficient condition,
not evidence that the condition is actually satisfied.

An elementary model explains one logical limitation of dimension
counting. Start with algebraically independent formal coordinates
$E,\eta,I_j,C_j,\beta_j$, put $\gamma_0=1/2$,
$\delta_0=E(\eta-1/2)$ and $\mu_j=\beta_j-1/2$.
The indicated joint-coordinate independence survives, while
$k(\gamma_0,\delta_0,\mu_1,C_1)$ has transcendence degree three.
This is a model of the algebraic input alone, not a model of the analytic
definitions of the actual constants. Additional analytic or arithmetic
information can therefore matter, but coordinate changes alone cannot
settle the four-number target.

The following questions remain beyond the proofs in this manuscript:
\begin{enumerate}
\item the individual arithmetic nature of $\gamma$, $\delta$, and any
specified member of $K_0(2),K_0(4),K_0(8)$;
\item full algebraic independence of $\gamma,\delta,\kappa,C$;
\item full rational integral-intersection classifications with two or
more harmonic Bessel coordinates at dependent scales, for example
$k(\beta(2),C(2),\beta(4),C(4))$;
\item an effective criterion selecting the existential density-one
independent set of integer arguments.
\end{enumerate}
The old matched-phase question is resolved by Theorems~\ref{fri:main}
and~\ref{fre:main}:
\[
 \Ebar\cap k(\eta,\beta(2),C(2))=k[\eta],\qquad
 \Ebar\cap k(e,\eta,\beta(2),C(2))=k[e,e^{-1},\eta].
\]
The derivative-coordinate part of the dependent-scale question is also
resolved, even in the presence of one harmonic coordinate
(Theorem~\ref{gde:main}), together with the matching exponential.
These closures concern integrality over
the E-value ring; they do not remove the arithmetic questions about
$\gamma$ and $\delta$ listed above.
In particular, adjoining $e$ to a coefficient field that already contains
$\Kexp$ does not add an independent coordinate. Also
$2\operatorname{Chi}(1)=\operatorname{Ei}(1)-E_1^{\rm tail}(1)$,
so these three classical values must not be counted as three independent
coordinates.

\subsection{Attribution and the restricted common-logarithm obstruction}
Mahler's theorem \cite[Theorem 3 and p.~170]{Mahler} proves the
transcendence, for nonzero algebraic $\alpha$ and compatible branches, of
\[
 \frac{\pi Y_0(\alpha)}{2J_0(\alpha)}
          -\log(\alpha/2)-\gamma.
\]
This includes the classical one-point context of the combinations used
above. The present synthesis concerns joint families, integral
intersections, relative fields, and the additional primitive coordinates;
it does not reclassify a known one-point result as a new discovery.

The common-logarithm statement in manuscript [030] \cite{Source030}
concerns meromorphic germs in a specified field with the coordinate
$T=\gamma+\log z$. Its finite-monodromy conclusion is a theorem about
that prescribed germ field. It explicitly does not assert the analogous
numerical conclusion after specialization at $z=1$. Thus it is not a
proof that every possible E-value method has a rank-three ceiling.
The numerical identity $\gamma=\Ein(1)-\delta/e$ is already a finite
algebraic combination; what remains missing is the necessary arithmetic
information about its terms. No universal impossibility theorem is used
in this paper.

The approximation results also have a precise scope. A rational sequence
can converge to a rational, algebraic irrational, or transcendental
number. The uniform $O(n^{-2})$ error on the square diagonal supplies
neither sufficiently small integer linear forms nor the denominator
estimates needed to infer irrationality.

\subsection{Concordance of the five predecessor notes}
This is a mathematical synthesis of the five notes below, with one
notation system and internal proofs. It is not a concatenation of their
abstracts. Superseded proofs are replaced by stronger theorems; distinct
hypotheses and applications are retained.
\begin{longtable}{@{}>{\raggedright\arraybackslash}p{0.24\textwidth}>{\raggedright\arraybackslash}p{0.70\textwidth}@{}}
\toprule
Predecessor & Material retained or strengthened \\
\midrule
\endhead
Intersection extensions \cite{Extensions} &
Uniform finite and infinite relation ideals; affine-slope alternatives;
arbitrary-distinct-scale linear nonintegrality; uniform factorial
asymptotics and the first correction; mirror-limit consequences. \\
\addlinespace
Joint closures \cite{Joint} &
All-order primitive towers and solvable PV closure; conjugation and
Hankel reconstruction; relative coefficient fields; linear-rank
classification; exact moment-integral coefficient cost. The positive
joint theorem is replaced by its proved complex version, and the mixed
polynomial exclusion by the full rational theorem under the same scale
independence hypothesis. \\
\addlinespace
$A,U,V$ addendum \cite{AUV} &
Exact relation ideal and rational-shift classification; polynomial
integrality; all nontrivial Laurent monomials; mixed rational subfields.
The former open $k(V)$ intersection is closed by the stronger
$k(V,C)$ theorem. \\
\addlinespace
Primitive and multiscale note \cite{Primitive} &
The single primitive is strengthened to simultaneous primitives at every
scale and every depth. The distinct Newton--Laguerre limit, its independent combinations,
matched multiscale rational exclusions, general one-point polynomial
classification, and the two root witnesses are retained. \\
\addlinespace
Density and monodromy note \cite{Density} &
Affine circuits; finite-deletion and density-one independence;
second-kind Bessel witnesses; continued embeddings and one-point
rational closures; independent-block equivalences and relative six-tuples.
The weaker factorial limit proof is replaced by the uniform expansion. \\
\bottomrule
\end{longtable}
The additional proofs in this revision include the all-depth primitive
extension, complex multiscale rational intersections, the full
Ein--Bessel intersection with its matching exponential, arbitrary
derivative scales with one harmonic coordinate, exact unit and radical
criteria, the finite real-adjunction defect, and a simultaneous six-family
density theorem valid at every algebraic point before selecting a common
integer subsequence. The earlier multivariate exponential-integral
classification, reconstruction, and critical asymptotics are retained.
These
statements are research results of this manuscript and remain subject
to external mathematical review; numerical checks are not substitutes
for their displayed proofs.


\clearpage
\begin{thebibliography}{99}
\footnotesize\raggedright
\setlength{\itemsep}{3pt}
\setlength{\parskip}{0pt}
\bibitem{Intersection}
Keio, \emph{What partial results prove together:
intersection theorems from a corpus of 244 checkpoint papers},
supplied unpublished manuscript, October 2026, 25 pp.;
\texttt{intersection\_theorems.tex}. Theorem 3.4, Proposition 4.1,
Theorem 4.2, and Remark 3.7 are the specific source inputs and target.

\bibitem{DLMF}
NIST Digital Library of Mathematical Functions, Sections 6.2, 6.5--6.6, 6.12, 10.5--10.6, 10.8, 10.27--10.34, and 10.40. \url{https://dlmf.nist.gov/}.

\bibitem{FR26}
S. Fischler and T. Rivoal, \emph{The ring of values of $E$-functions
is not a field}, preprint, October 2026,
\href{https://arxiv.org/abs/2610.16603}{arXiv:2610.16603}.
The component-continuation property is in Section 2; the phase
embedding is Proposition 3.1.

\bibitem{Beukers}
F. Beukers,
\emph{A refined version of the Siegel--Shidlovskii theorem},
Ann. of Math. \textbf{163} (2006), 369--379.
\url{https://doi.org/10.4007/annals.2006.163.369}.

\bibitem{Kolchin}
E. R. Kolchin,
\emph{Algebraic groups and algebraic dependence},
Amer. J. Math. \textbf{90} (1968), 1151--1164.

\bibitem{Rivoal}
T. Rivoal,
\emph{On the arithmetic nature of the values of the gamma function,
Euler's constant, and Gompertz's constant},
Michigan Math. J. \textbf{61} (2012), 239--254.
\url{https://rivoal.perso.math.cnrs.fr/articles/gammater.pdf}.

\bibitem{Extensions}
Keio,
\emph{Extensions of a Joint Independence Theorem for Euler--Gompertz
and Bessel Values}, unpublished research note, 22 October 2026, 9 pp.

\bibitem{Periods}
Keio,
\emph{Exponential-Integral Periods across Arithmetic, Pad\'e, Hankel,
Reciprocal Geometry, and Frobenius Framing}, supplied manuscript [001].
The operative sections are ``Multipoint one-level polyexponentials of
arbitrary weight'' and ``Positive mixed Hankel determinants and finite
arithmetic reconstruction.''

\bibitem{Source021}
Keio,
\emph{Arithmetic Alternatives for Euler's Constant}, supplied manuscript
[021], section on the complete classification at linear rank three.

\bibitem{Hamiltonian}
Keio,
\emph{Arithmetic Rank and Hamiltonian Moment Integrals}, supplied
unpublished manuscript, Proposition on the visible support lattice.

\bibitem{Joint}
Keio, \emph{Explicit Independence and Polynomial Exclusion for Euler--Gompertz and Bessel Constants}, unpublished research note, October 2026, 15 pp.

\bibitem{AUV}
Keio, \emph{Polynomial Integrality and Rational Independence for Three Euler--Bessel Combinations}, unpublished research note, October 2026, 8 pp.

\bibitem{Vertex}
Keio, \emph{Exact Recurrences, Algebraic Independence, and Certified
Statistical Recovery}, supplied manuscript [036], 2026.

\bibitem{Gamma}
Keio, \emph{Gamma-Type Functions: Arithmetic and Differential
Independence, Probability Laws, and Interpolation}, supplied
manuscript [025], 2026.

\bibitem{Source030}
Keio, supplied manuscript [030], \emph{Arithmetic and Observation
Classifications, Recovery, and Transcendence Barriers} (2026),
corollary \texttt{p07:rgam:commonlog} and the following scope paragraph.

\bibitem{Mahler}
K. Mahler, Applications of a theorem by A. B. Shidlovski,
\emph{Proceedings of the Royal Society of London, Series A}
\textbf{305} (1968), 149--173.
\href{https://doi.org/10.1098/rspa.1968.0111}{doi:10.1098/rspa.1968.0111}.
See Theorem 3 and its explicit consequence on page 170.

\bibitem{Primitive}
Keio, \emph{Bessel Primitives, Critical Laguerre Limits, and Multiscale E-Value Exclusion}, unpublished research note, October 2026, 11 pp.

\bibitem{Source019}
Keio, supplied manuscript [019], \emph{Euler's Constant Complete
Integrated Research Manuscript} (2026), theorem
\texttt{ec:moving-theorem}.

\bibitem{Source033}
Keio, supplied manuscript [033], \emph{Euler's Constant Arithmetic,
Monodromy, and Gamma-Related Structures} (2026), section
\texttt{p5:pm:section} and theorem \texttt{p5:pm:embedding}.
Its continued phase construction is proved explicitly above.

\bibitem{Density}
Keio, \emph{Density-One Independence and Monodromy Closures for Bessel Values}, unpublished research note, October 2026, 15 pp.
\end{thebibliography}
\end{document}
